% !TEX root = mixed-sklyanin-cubics.tex
\documentclass[11pt,a4paper]{ltjsarticle}
\usepackage{amsmath,amssymb,amsthm,mathtools,bm,mathrsfs}
\usepackage{luatexja-fontspec}
\usepackage{enumitem,booktabs,longtable,array,needspace}
\usepackage{xcolor}
\usepackage{etoolbox,xr-hyper}
\usepackage[colorlinks=true,linkcolor=blue!45!black,citecolor=blue!45!black,urlcolor=blue!45!black]{hyperref}
\DeclareMathSizes{9}{9}{6}{5}
\DeclareMathSizes{9.85501}{9}{6}{5}
\IfFileExists{/usr/share/fonts/opentype/noto/NotoSerifCJK-Regular.ttc}{%
  \setmainjfont[Path=/usr/share/fonts/opentype/noto/]{NotoSerifCJK-Regular.ttc}%
}{\setmainjfont{HaranoAjiMincho-Regular}}
\setlength{\textwidth}{15.6cm}
\setlength{\oddsidemargin}{0.25cm}
\setlength{\evensidemargin}{0.25cm}
\setlength{\textheight}{22.6cm}
\setlength{\topmargin}{-0.2cm}
\setlength{\parindent}{1em}
\setlength{\parskip}{1pt}
\numberwithin{equation}{section}
\newtheorem{theorem}{定理}[section]
\newtheorem{maintheorem}{主定理}
\renewcommand{\themaintheorem}{\Alph{maintheorem}}
\newtheorem{lemma}[theorem]{補題}
\newtheorem{proposition}[theorem]{命題}
\newtheorem{corollary}[theorem]{系}
\newtheorem{remark}[theorem]{注意}
\newtheorem{definition}[theorem]{定義}
\newtheorem{example}[theorem]{例}
\newtheorem{conjecture}[theorem]{予想}
\renewcommand{\proofname}{証明}
% 三稿で共通に用いる記号。本文と証明はそれぞれ一つの稿に置く。
\providecommand{\ZZ}{\mathbb Z}
\providecommand{\cK}{\mathcal K}
\providecommand{\cP}{\mathcal P}
\providecommand{\cD}{\mathcal D}
\providecommand{\cC}{\mathcal C}
\providecommand{\cE}{\mathcal E}
\providecommand{\cF}{\mathcal F}
\providecommand{\cM}{\mathcal M}
\providecommand{\fm}{\mathfrak m}
\providecommand{\Perf}{\operatorname{Perf}}
\providecommand{\Ku}{\operatorname{Ku}}
\providecommand{\MF}{\operatorname{MF}}
\providecommand{\HH}{\operatorname{HH}}
\providecommand{\Dsg}{D_{\mathrm{sg}}}
\providecommand{\RHom}{\operatorname{RHom}}
\providecommand{\REnd}{\operatorname{REnd}}
\providecommand{\Ind}{\operatorname{Ind}}
\providecommand{\Def}{\operatorname{Def}}
\providecommand{\thick}{\operatorname{thick}}
\providecommand{\Spf}{\operatorname{Spf}}
\providecommand{\Id}{\operatorname{Id}}
\providecommand{\sym}{\operatorname{sym}}
\providecommand{\Cone}{\operatorname{Cone}}
\providecommand{\Pic}{\operatorname{Pic}}
\providecommand{\otL}{\mathbin{\otimes^{\mathbf L}}}
\newtheorem{problem}[theorem]{問題}

\newcommand{\C}{\mathbb C}
\newcommand{\CC}{\mathbb C}
\newcommand{\PP}{\mathbb P}
\newcommand{\OO}{\mathcal O}
\newcommand{\Sym}{\operatorname{Sym}}
\newcommand{\Spec}{\operatorname{Spec}}
\newcommand{\Proj}{\operatorname{Proj}}
\newcommand{\Hom}{\operatorname{Hom}}
\newcommand{\End}{\operatorname{End}}
\newcommand{\Ext}{\operatorname{Ext}}
\newcommand{\Aut}{\operatorname{Aut}}
\newcommand{\GL}{\operatorname{GL}}
\newcommand{\PGL}{\operatorname{PGL}}
\newcommand{\Gr}{\operatorname{Gr}}
\newcommand{\rank}{\operatorname{rank}}
\newcommand{\qgr}{\operatorname{qgr}}
\newcommand{\QGr}{\operatorname{QGr}}

% 他稿の番号には R (可換復元), M (モジュライ・降下), S (混合族) を添える。
% clean build では scripts/build_papers.py が三稿の aux を順に生成・収束させる。
\newcommand{\SetPaper}[1]{%
  \def\CurrentPaper{#1}%
  \ifstrequal{#1}{main}{}{\externaldocument[main-][nocite]{build/nc-cubic}[nc-cubic.pdf]}%
  \ifstrequal{#1}{moduli}{}{\externaldocument[moduli-][nocite]{build/cubic-moduli-descent}[cubic-moduli-descent.pdf]}%
  \ifstrequal{#1}{mixed}{}{\externaldocument[mixed-][nocite]{build/mixed-sklyanin-cubics}[mixed-sklyanin-cubics.pdf]}%
}
\makeatletter
\def\PaperTag@main{R}
\def\PaperTag@moduli{M}
\def\PaperTag@mixed{S}
\newcommand{\PaperRefText}[2]{\textnormal{[\csname PaperTag@#1\endcsname,~\ref{#1-#2}]}}
\newcommand{\PaperEqRefText}[2]{\textnormal{[\csname PaperTag@#1\endcsname,~\eqref{#1-#2}]}}
\DeclareRobustCommand{\paperref}[2]{\ifmmode\text{\PaperRefText{#1}{#2}}\else\PaperRefText{#1}{#2}\fi}
\DeclareRobustCommand{\papereqref}[2]{\ifmmode\text{\PaperEqRefText{#1}{#2}}\else\PaperEqRefText{#1}{#2}\fi}
\DeclareRobustCommand{\paperpageref}[2]{\textnormal{[\csname PaperTag@#1\endcsname,~p.~\pageref{#1-#2}]}}
\newcommand{\PaperReferenceGuide}{%
  本稿と関連する二稿の参照では、R を可換復元稿\cite{NCCubicRealization}、
  M をモジュライ・降下稿\cite{NCCubicModuliDescent}、
  S を混合族稿\cite{NCCubicMixedSklyanin}の略号とする。
  他稿の結果は、例えば定理 [R, 2.1] のように稿の略号とその稿内の番号を併記する。
  略号のない番号は本稿内の参照である。
}
\makeatother

\SetPaper{mixed}

\title{混合 Sklyanin 三次の滑らかさと残余圏\\
\large 有限中心・局所代数・可換比較の候補}
\author{}
\date{2026年11月22日・研究稿 v50}
\hypersetup{pdftitle={混合 Sklyanin 三次の滑らかさと残余圏},
pdfsubject={有限中心・局所代数・可換比較の候補}}

\begin{document}
\maketitle
\begin{abstract}
中心一次変数を共有する Sklyanin ブロックの三次中心商を調べる。
$S(1,-1,-1)$ の全定数型変形について普遍中心元を明示し、
有限中心性、PI 次数六、Azumaya locus と正則中心の一致を証明する。
指定した二ブロック五様体族では、各因子の十四節点の完備局所代数と
中央切断を決定し、非 smooth 係数を六十四点の集合として厳密に求める。
全 smooth locus 上の相対残余圏を構成し、Hochschild homology と
有理数値 Euler 形式を計算する。
また同じ判別集合を持つ可換三次族と、零係数における Eckardt 根スタック模型を与える。
この明示候補への Ku 同値は一般点でも未証明である。
一方、分離型比較と可換実現 locus の開性を用いれば、
任意の固定した定数型混合係数に対する余有限の可換 Ku 実現は得られる。
点スキームと Poisson 零スキームの一致を一般に期待できないことも、
被約な孤立定常点を持つ例として本文で説明する。
\end{abstract}

\setcounter{tocdepth}{1}
\tableofcontents

\section{序論}
\label{mix:sec:introduction}

\subsection{対象と比較問題}

非可換三次の Kuznetsov 成分を可換三次の成分として実現する問題では、
二次関係に中心変数が現れる混合表示が、楕円因子の分離テンソル比較の
次に調べるべき具体的な対象となる。
本稿は、そのうち一つの中心一次変数を共有する
$S(1,-1,-1)$ の定数型変形を扱う。
基礎体は $\C$ とし、$\alpha=(\alpha_1,\alpha_2,\alpha_3)\in\C^3$ に対して
\[
 \mathcal B_\alpha
 =\C\langle x,y,z,u\rangle/
 \bigl([y,z]-x^2+\alpha_1u^2,\,
       [z,x]-y^2+\alpha_2u^2,\,
       [x,y]-z^2+\alpha_3u^2,\,
       [u,x],[u,y],[u,z]\bigr)
\]
と置く。各生成元の内部次数は一である。
式\eqref{cmo:eq:mixed-center}の正則な三次中心元を $\Omega_\alpha$ とし、
$k\ge2$、$a_i\in\C^\times$ に対して
\[
 A_c=
 \left(\bigotimes_{\C[u],\,i=1}^k\mathcal B_{\alpha^{(i)}}\right)
 \Big/\left(\sum_i a_i\Omega_{\alpha^{(i)}}+cu^3\right)
 \qquad(c\in\C)
\]
を考える。滑らかな場合の次元は $n=3k-1$ である。
アファイン中心 chart の一因子を
$D_\alpha=\mathcal B_\alpha/(u-1)$ と書く。

$D^b(\operatorname{qgr}A_c)$ の全 coherent 残余圏と、
有限個の線束から生成する圏の残余部分は区別する。
前者を $\Ku_{\mathrm{coh}}(A_c)$ と書き、
$\langle\OO,\OO(1),\ldots,\OO(n-2)\rangle$ の右直交で定める。
後者の smoothness を証明した後に coherent 飽和を適用するのが、
本稿の可換比較と滑らかさの議論の順序である。
有限次元 chart 加群の射影次元、全 coherent 圏の proper 性、
線束生成圏の smoothness も同一の条件として扱わない。

\subsection{中心結果と成立範囲}

本稿の中心結果は次の三段階からなる。
完全な仮定と証明は各参照先に置く。

\begin{enumerate}
 \item \emph{全定数型変形の中心と、指定因子の局所構造。}
 定理\ref{cmo:thm:finite-center}は全 $\alpha\in\C^3$ に対して中心の生成元と
 有限中心性を与え、命題\ref{cmo:prop:center-hypersurface-degree}と
 定理\ref{cmo:thm:azumaya-regular}は PI 次数六と正則中心上の Azumaya 性を与える。
 さらに $\alpha^{(1)}=(1,4,9)$、$\alpha^{(2)}=(16,25,49)$ について、
 定理\ref{cmo:thm:fourteen-node-orders}は各因子の全十四節点で
 完備局所代数を $\operatorname{End}_R(R^{\oplus a}\oplus I^{\oplus b})$、
 $a+b=6$ と決定する。これにより単純加群と中央切断を調べる。

 \item \emph{指定五様体族の smooth locus と残余不変量。}
 上の二因子と重み $1,100$ を固定する。
 定理\ref{cmo:thm:exact-sixty-four}は
 $D^b(\operatorname{qgr}A_c)$ と $\Ku_{\mathrm{coh}}(A_c)$ が
 smooth proper である条件を、明示的な次数六十四多項式 $P_{64}$ による
 $P_{64}(c)\ne0$ と同値にする。
 この $P_{64}$ は、定常点の出現を測る有限 norm $N$ とは互いに素である。
 命題\ref{cmo:prop:full-relative-hh}は
 $U=\mathbb A^1\setminus V(P_{64})$ 全体で相対 smooth proper な模型を構成し、
 全 $c\in U$ に対して $HH_0=\C^2$、$HH_{\pm1}=\C^{21}$、
 他次数零を与える。系\ref{cmo:cor:rational-euler-comparison}の
 有理数値 Euler 形式も、後述の可換候補と一致する。

 \item \emph{可換実現の存在と、指定候補への比較の区別。}
 系\ref{cmo:cor:cofinite}は任意の固定した
 $k\ge2$、$\alpha^{(i)}\in\C^3$、$a_i\ne0$ に対し、
 有限個を除く $c$ で何らかの滑らかな可換三次への Ku 同値を与える。
 他方、定理\ref{cmo:thm:commutative-discriminant-candidate}は指定五様体族について
 同じ $P_{64}$ と局所通常二重点を持つ明示的な可換族 $X_c$ を構成するが、
 この特定の $X_c$ への同値は一般点でも証明していない。
 $c=0$ では定理\ref{cmo:thm:eckardt-root-residual}により、
 根スタック上の十六例外対象の右直交として
 $\Ku(X_0)^{\mu_2}$ を記述する。非可換側からこの圏への比較核は未構成である。
\end{enumerate}

点スキームは可換 Ku 実現に付随する零点不変量として無条件には使えない。
第\ref{mix:sec:poisson-limits}節では、
$N(c)=0$ の smooth proper な五様体について被約孤立定常点が現れ、
全ての滑らかな可換三次五様体の Poisson 零スキームと一致しないことを示す。
これは点スキームの一致という素朴な期待の限界を示す補助例であり、
可換 Ku 実現そのものを否定する例ではない。
全 Poisson 構造への排除は、未公表入力\cite{NCCubicA}の
全次元正規形に依存する。

\subsection{先行研究と本稿で加える具体化}

Sklyanin 代数の斉次中心の生成は Smith--Tate
\cite[Theorems~3.7, 4.8]{GRSklyaninST}、
potential の PBW 変形と三次中心元は Etingof--Ginzburg
\cite[Theorem~3.3.2]{NFMEG}に基づく。
Walton \cite[Definition~V.1, Remark~V.16]{CMOWaltonThesis}の
位数六の中心拡大に関する持上げの方針に対し、
本稿は斉次パラメータ $(1,-1,-1)$、線形補正が零、定数補正が任意という
部分族で六次中心元を明示し、その恒等式を検証する。
associated graded が PI であることだけから変形の有限中心性を推論しない。
Azumaya locus の議論には Brown--Yakimov
\cite[Corollary~5.2]{CMOBrownYakimov}と、
斉次 Sklyanin 代数についての Walton--Wang--Yakimov
\cite[Theorem~1.3(1), Lemma~5.4]{CMOWaltonWangYakimov}を用いる。

局所代数の記述には非可換クレパント解消の理論
\cite{CMOIyamaWemyssMM,CMOIyamaWemyssCA}を、
全 smooth locus 上の Hochschild 群には Mathew の相対退化定理
\cite[Theorem~1.3]{CMOMathewDegeneration}を使う。
Eckardt 模型の残余圏は Orlov の blow-up 公式
\cite[Theorem~4.3]{CMOOrlovBlowup}と根スタックの半直交分解
\cite[Theorem~1.6]{CMOIshiiUedaRoot}による。
これらの一般理論自体を本稿の成果とはしない。
本稿での内容は、指定族に対する中心・局所切断・対角・残余模型の具体的な組合せと、
その適用範囲の決定である。個々の具体化について後続文献まで含む
新規性の網羅的な確認はなお完了していない。

余有限可換実現には、可換復元を扱う主稿の分離型比較
（定理\paperref{main}{csc:thm:comparison}）、coherent 飽和
（定理\paperref{main}{crs:thm:coherent}）、可換実現 locus の開性
（定理\paperref{main}{grl:thm:open-locus}）を入力とする。
各々の具体的な仮定を本文の適用箇所で確認する。
有限中心元・有限行列・終結式などの計算証明書は本文に所在を記し、
計算の再現と大域的な圏同値の証明を区別する。

\subsection{構成}

第\ref{cmo:sec:points}節で点スキームを有限 norm から計算し、
第\ref{mix:sec:modules}節で小次元単純加群と中心方程式の高次作用を調べる。
第\ref{mix:sec:centers}--\ref{mix:sec:local-orders}節で
普遍中心、Azumaya 幾何、節点の局所代数を構成し、
第\ref{mix:sec:perfect-modules}節で有限台加群と残余 Ext を扱う。
第\ref{mix:sec:relative-smoothness}--\ref{mix:sec:smooth-locus}節が
相対 smoothness、余有限可換実現、正確な smooth locus と Hochschild 群を与える。
第\ref{mix:sec:commutative-candidate}--\ref{mix:sec:eckardt}節で
明示可換候補と Eckardt 模型を構成し、
第\ref{mix:sec:poisson-limits}節を点・零スキームの不一致の例に充てる。
最後に、得られた模型から進むための比較問題を述べる。

\PaperReferenceGuide

% !TEX root = ../../mixed-sklyanin-cubics.tex
\section{中心一次変数が一つの混合ブロックと有限 norm}
\label{cmo:sec:points}

分離表示の定理\paperref{main}{csc:thm:comparison}では、中心変数は各楕円因子の
二次関係に現れない。本稿では、この仮定を点関手の計算から取り除く。
結果は中心変数を一つ共有する表示についての定理であり、
一般の混合表示の点関手を計算した後、$S(1,-1,-1)$ の定数型変形については
有限中心性と収縮族の相対 smoothness を証明し、余有限可換 Ku 実現を得る。

\subsection{有限な定常環と point scheme}

基礎体を $\C$、$R=\C[u]$、$k\geq2$ とする。
各 $i$ について、次数一生成元 $x_i,y_i,z_i,u$ を持ち、
$u$ が中心である二次代数 $\mathcal B_i$ を取る。
$\mathcal B_i/(u)=S_i$ の双線型対応は、滑らかな平面三次 $E_i$ 上の
固定点なしの自己同型 $\sigma_i$ のグラフであると仮定する。
さらに可換化は
\begin{equation}\label{cmo:eq:abelian}
 (\mathcal B_i)_{\mathrm{ab}}
   =R[x_i,y_i,z_i]/(q_{i1},q_{i2},q_{i3}),\qquad
 q_{ia}\in R[x_i,y_i,z_i]_2
\end{equation}
であり、$q_{i1}|_{u=0},q_{i2}|_{u=0},q_{i3}|_{u=0}$ は
$\C[x_i,y_i,z_i]$ の正則列とする。
例えば、それらが滑らかな三次式の三つの偏微分ならばこの条件を満たす。
active 変数と $u$ の混合項は許す。

異なるブロックの元を可換とする $R$ 上の tensor 積
\[
 B=\mathcal B_1\otimes_R\cdots\otimes_R\mathcal B_k
\]
を考える。正則中心三次 $w\in Z(B)_3$ は、$u=0$ において
$\sum_i a_i\Omega_i$ に移ると仮定する。
ここで各 $\Omega_i$ は $S_i$ の全点加群を消す中心三次である。
$c\in\C$ に対し $A_c=B/(w+cu^3)$ と置く。
次の定理の点関手の主張自体には、$A_c$ の滑らかさを仮定しない。

\begin{theorem}[一中心変数での定常点の有限判定]
\label{cmo:thm:point-norm}
上の仮定の下で
\[
 T=R[x_1,y_1,z_1,\ldots,x_k,y_k,z_k]/(q_{ia})
\]
は $R$ 上階数 $L=8^k$ の有限自由加群である。
$Q=T/(u-1)$、$b=w_{\mathrm{ab}}|_{u=1}\in Q$ とすると、任意の試験環に対して
\begin{equation}\label{cmo:eq:point-scheme}
 \operatorname{Pt}(B)=\left(\coprod_i E_i\right)\amalg\operatorname{Spec}Q,
 \qquad
 \operatorname{Pt}(A_c)=\left(\coprod_i E_i\right)
                  \amalg\operatorname{Spec}(Q/(b+c)).
\end{equation}
shift は $E_i$ 上で $\sigma_i$、定常成分上で恒等である。

乗法 $m_b:Q\to Q$ の norm 多項式
\begin{equation}\label{cmo:eq:norm}
 N(c)=\det(c\operatorname{id}_Q+m_b)\in\C[c]
\end{equation}
は次数 $8^k$ の monic 多項式であり、
\[
 \operatorname{Pt}(A_c)=\coprod_iE_i
 \quad\Longleftrightarrow\quad N(c)\ne0.
\]
従って定常点が現れるパラメータは有限個で、
その零次 Fitting イデアルは $(N(c))$ である。
これは閉点集合だけでなく、$c$ に関する族の有限表示としての主張である。
\end{theorem}

\begin{proof}
$u,q_{11},\ldots,q_{k3}$ の同時零点は affine 空間の原点だけである。
これらは $3k+1$ 変数多項式環の斉次 system of parameters なので正則列である。
従って $T$ は Hilbert 級数
\[
 H_T(t)=\frac{(1+t)^{3k}}{1-t}
\]
を持つ。$T/uT$ の斉次基底を $T$ へ持ち上げると、
次数についての帰納法によりそれらは $R$-加群として $T$ を生成する。
基底の次数を保つ自由 $R$-加群からの全射の両辺の Hilbert 級数は同じなので、
その核は零である。これで階数 $8^k$ の自由性を得る。
また $C=\operatorname{Proj}T$ は $u=0$ と交わらず、
$u\ne0$ の chart で $\operatorname{Spec}Q$ と同一視される。

双線型関係を満たす二つの点を $p,p'$ と書く。
$u(p)\ne0$ なら、中心交換関係
$u(p)x'_i=x_i u(p')$ などと $p'\ne0$ から $u(p')\ne0$ であり、
両方を一に規格化すると全座標が等しい。
残る関係は式\eqref{cmo:eq:abelian}の $3k$ 本の二次式 $q_{ia}=0$ である。
よってこの chart の対応は $\Delta_C$ である。

$u(p)=0$ の幾何点では $u(p')=0$ となる。
二つ以上のブロックが非零ならブロック間交換関係が各ブロックの二点を比例させ、
$\sigma_i$ の固定点を生じるので不可能である。
従ってその幾何点は純ブロックの $\Gamma_{\sigma_i}$ にある。
そこでは固定点なしという仮定により、同一ブロックの二点の
二行二列小行列式の少なくとも一つが単元となる。
中心交換関係の対応する二式は $u,u'$ を環の中で消し、
他ブロックとの交換関係も同じ単元小行列式によって全ての他座標を消す。
残る対応は $\Gamma_{\sigma_i}$ そのものである。
従って横断方向の冪零構造もなく、
\[
 \Gamma_B=\left(\coprod_i\Gamma_{\sigma_i}\right)\amalg\Delta_C
\]
をスキームとして得る。この局所環上の消去は任意の試験環でも成立する。

対応が自己同型のグラフなので、無限点列は初期点から一意に定まる。
純ブロックでは $u=0$、かつ $\Omega_i$ は普遍点加群を消す。
定常列では三次関係の三重線型化はその可換化の評価であり、
$u=1$ で $b+c$ となる。これで式\eqref{cmo:eq:point-scheme}を得る。

普遍パラメータ $c$ を使うと、$Q[c]$ 上の乗法 $m_b+c$ は単射である。
実際、非零の多項式ベクトルの最高次係数は $c\operatorname{id}$ により残る。
従って
\[
 0\longrightarrow Q[c]\xrightarrow{m_b+c}Q[c]
   \longrightarrow Q[c]/(b+c)\longrightarrow0
\]
は階数 $L$ の正方表示を与え、その零次 Fitting イデアルは行列式である。
閉点 $c$ で商が零となることは $b+c$ が $Q$ の単元であること、
従って $N(c)\ne0$ と同値である。
\end{proof}

\subsection{滑らかな Poisson 標的での定常零点の消滅}

同じ変数分解の下で
\[
 H=\sum_{i=1}^kF_i(x_i,y_i,z_i,u)+G(u),\qquad
 X=V(H)\subset\mathbb P^{3k}
\]
とする。各 $F_i$ は斉次三次であり、$g_i=F_i|_{u=0}$ は滑らかな平面三次を定め、
$X$ も滑らかとする。
$P_i=\partial_{x_i}\wedge\partial_{y_i}\wedge\partial_{z_i}$、
$d_i\in\C^\times$ を取り、
\[
 \beta=\sum_i d_i\iota_{dF_i}P_i
\]
の射影化を $\pi$ と書く。共通の $u$ は各項の Casimir であり、
異なるブロックの Schouten 括弧も零なので $\beta$ は Poisson で、$H$ を保つ。

\begin{proposition}[一中心変数の Poisson 零スキーム]
\label{cmo:prop:poisson-zero}
上の仮定の下で、非被約構造を含めて
\begin{equation}\label{cmo:eq:poisson-zero}
 Z_X(\pi)=\coprod_iV(g_i)
\end{equation}
である。右辺はそれぞれの純ブロック平面に埋め込む。
\end{proposition}
\begin{proof}
Euler vector field を $\xi$ とする。
命題\paperref{main}{pps:prop:zero-scheme}と同じく零イデアルは $\xi\wedge\beta$ の係数である。
$u\ne0$ では $u\,d_i(F_i)_{x_i}$ などの係数から、
全ての active 偏微分が零となる。
$H=0$ と Euler 恒等式から $uH_u=0$ も従うので、
もしこの chart に零点があれば $X$ の特異点となる。
従って滑らかさの仮定によりこの chart の零スキームは空である。

$u=0$ の幾何点で $s$ 個のブロックが非零なら、$g_i$ の滑らかさから
$\beta$ の階数は $2s$ である。
射影化が零なら $\beta\in\xi\wedge V$ なので階数は高々二であり、$s=1$ である。
純ブロック $i$ では残る係数は $3d_i g_i$ である。
さらにその点の近傍では $(F_i)_{x_i},(F_i)_{y_i},(F_i)_{z_i}$ の一つが単元なので、
混合外積の係数は $u$ と全ての他ブロック座標を環の中で消す。
残る式は $g_i=0$ となり、式\eqref{cmo:eq:poisson-zero}を得る。
\end{proof}

\subsection{六十四個の相異なる例外値を持つ表示}

次の例では、定常成分が現れる全パラメータを有理数で記述できる。
これらの係数での coherent smoothness は
系\ref{cmo:cor:fixed-zero-norm-smooth}で確定する。
さらに系\ref{ppz:cor:reduced-point-poisson-counterexample}は、
そこで生じる孤立定常点から被約 point/Poisson 対応も排除する。

\begin{example}[混合二ブロックの完全な norm 分解]
\label{cmo:ex:sixty-four}
$\alpha=(\alpha_1,\alpha_2,\alpha_3)$ に対し、$u$を中心とする代数
\[
 \mathcal B_\alpha=\C\langle x,y,z,u\rangle/
 \left([y,z]-x^2+\alpha_1u^2,
       [z,x]-y^2+\alpha_2u^2,
       [x,y]-z^2+\alpha_3u^2,
       [u,x],[u,y],[u,z]\right)
\]
を考える。次の二組を用いる。
\[
 \alpha=(1,4,9),\qquad \alpha'=(16,25,49).
\]
各代数で、$\sym(xyz)$を六語の平均として
\begin{equation}\label{cmo:eq:mixed-center}
 \Omega_\alpha
 =\frac{x^3+y^3+z^3}{3}-\frac23\sym(xyz)
       -u^2(\alpha_1x+\alpha_2y+\alpha_3z)
\end{equation}
は中心元である。
$B=\mathcal B_\alpha\otimes_{\C[u]}\mathcal B_{\alpha'}$、
$w_c=\Omega_\alpha+100\Omega_{\alpha'}+cu^3$と置く。
定理\ref{cmo:thm:point-norm}の$Q$は$\C^{64}$と同型であり、
そのnormは
\begin{equation}\label{cmo:eq:explicit-norm}
 \begin{gathered}
 N(c)=\prod_{r\in\mathcal R}\prod_{s\in\mathcal S}(c+r+100s),\\
 \mathcal R=\{\pm28,\ \pm16,\ \pm56/3,\ \pm28/3\},\qquad
 \mathcal S=\{\pm448,\ \pm196,\ \pm176,\ \pm284/3\}
 \end{gathered}
\end{equation}
となる。六十四個の根は全て相異なる。
各根では点スキームにちょうど一つの被約な定常点が加わり、
根以外では二つの楕円曲線だけとなる。特に$c=0$では定常成分は空である。
\end{example}

\begin{proof}
$u=0$の代数はSklyanin parameters $(1:-1:-1)$を持ち、
点曲線は$E=V(x^3+y^3+z^3+xyz)$である。
$1+(1/3)^3\ne0$より滑らかで、
対角関係$x^2=y^2=z^2=0$は射影共通零点を持たない。
命題\paperref{main}{esi:prop:factor}と同じ行列の核の議論で、
双線型対応は固定点なしの楕円グラフである。
このSklyanin代数は\cite[Theorem~1.1, Corollary~1.3]{GRSklyaninTV}によりKoszulである。

$u=1$の関係は定数項だけを加えたPBW変形である。
実際、$r_x=yz-zy-x^2$とその巡回置換を用いると、
唯一の三次Koszul重なりは
\[
 x\otimes r_x+y\otimes r_y+z\otimes r_z
   =r_x\otimes x+r_y\otimes y+r_z\otimes z
\]
で生成される。
Braverman--Gaitsgoryの判定
\cite[Main Theorem~0.5, Lemma~3.3]{CMOBravermanGaitsgory}における
一次補正は零、定数補正は$r_a\mapsto-\alpha_a$である。
残る条件はこの重なり上の
$\beta\otimes1-1\otimes\beta=0$であり、定数の可換性から成立する。
従って表示した斉次代数はこのPBW変形のRees代数で、$\C[u]$上自由である。

式\eqref{cmo:eq:mixed-center}の中心性は、指定した二組の係数について
有理数上の有限な線型代数で直接検証する。
次数$d$の全語を基底に、二次関係を長さ$d-2$の全左右語で挟んだ行列を作る。
$d=3,4$の商次元は$20,35$であり、$\Omega_\alpha$の次数三剰余は非零、
四生成元との全交換子の次数四剰余は零となる。
再現コードは\texttt{verify\_point\_poisson\_obstruction.py}の
\texttt{verify\_mixed\_central\_norm}にあり、全行列を関係式から再生成して厳密消去する。
中心性は生成元とのこの有限個の交換子で尽くされ、
全次数PBW性は前段の判定で別に示している。
$u=0$では、$S=x^3+y^3+z^3$、
$C_+=xyz+yzx+zxy$、$C_-=yxz+xzy+zyx$と書くと
\[
 C_+-C_--S=0,\qquad 3c_3=2C_--2S,\qquad
 \Omega_0=\frac{3S-C_+-C_-}{9}=-\frac13c_3.
\]
ここで標準中心元$c_3$の式は
\cite[Section~2, p.~3]{GRSklyaninDLL}に$(a,b,c)=(1,-1,-1)$を代入したものである。
従って$\Omega_0$は全点加群を消す。
二因子の中心三次和は$B/(u)$上正則であり、
$u$の正則性と次数についての帰納法から、$w_c$も全ての$c$で正則である。

可換化して$u=1$とすると$x^2=\alpha_1,y^2=\alpha_2,z^2=\alpha_3$である。
各因子は八つの相異なる有理点に分解し、
その上の中心元の値は
\[
 b_\alpha=-\frac23(\alpha_1x+\alpha_2y+\alpha_3z+xyz)
\]
となる。$(x,y,z)=(\pm1,\pm2,\pm3)$での全評価値が$\mathcal R$、
$(\pm4,\pm5,\pm7)$での全評価値が$\mathcal S$である。
従って$Q=\C^{64}$上の乗法は対角となり、式\eqref{cmo:eq:explicit-norm}を得る。
$\mathcal S$の異なる二元の差の絶対値は少なくとも$20$、
$\mathcal R$の直径は$56$なので、$100\mathcal S$を中心とする八つの組は交わらない。
各組の八値も相異なるため、全ての根は単純である。
また$100\min_{s\in\mathcal S}|s|>\max_{r\in\mathcal R}|r|$より$N(0)\ne0$である。
点関手の結論は定理\ref{cmo:thm:point-norm}から従う。
\end{proof}

\begin{remark}[可換 Ku 実現と点スキーム比較との関係]
\label{cmo:rem:scope}
$k=2,m=1$ は $N=7,n=5$ に当たり、定理\ref{cmo:thm:point-norm}の
有限定常環の長さは $64$ である。
一般の中心係数では点スキームは二つの楕円曲線だけとなり、
命題\ref{cmo:prop:poisson-zero}の零スキームと同じ形を持つ。
しかし、一般の混合表示について、その楕円曲線を active 制限に持つ
\emph{同じ}可換 Ku 標的を構成したことにはならない。
分離表示では定理\paperref{main}{csc:thm:comparison}がこの不足を埋める。

定常点の出現だけから反例を結論するのではなく、
source の smoothness と全可換標的の排除を別に要する。
指定した六十四係数では前者を系\ref{cmo:cor:fixed-zero-norm-smooth}、
後者を孤立 Poisson 零点の次元制約
（定理\ref{ppz:thm:isolated-zero-dimension}）で満たす。
有限 norm による明示的な例外集合は、これらの次の検証対象を与える。
\end{remark}

% !TEX root = ../../mixed-sklyanin-cubics.tex
\section{小次元単純加群と中心方程式の高次作用}
\label{mix:sec:modules}

\subsection{例外的な定常点の自己 Ext と中心方程式の高次作用}
\label{cmo:subsec:character-ext}

定常点の出現と商の非滑らかさは同じ条件ではない。
本節の明示族では、例外値の定常 character の自己 Ext を全次数で計算できる。
結論は自己 Ext の次数付きベクトル空間についてであり、
その Yoneda 積、character の射影次元、商の smoothness は区別する。

高次 homotopy から超曲面商上の自由分解を作る一般構成は
Shamash--Eisenbud 理論に属する
\cite[Theorem~7.2]{NFMEisenbudCI}。
この可換環上の構成と、本稿の正則中心元に対する非可換版を区別し、
後者の完全性は補題\paperref{main}{mbnp:lem:shamash}の濾過による証明を用いる。
また非斉次三次 potential の Calabi--Yau 性、PBW 性と中心元は
Etingof--Ginzburg の既知の理論
\cite[Theorem~3.3.2, Sections~3.4,~8.3]{NFMEG}に位置付けられる。
本節では指定中心元の普遍な多項式証明書を作り、有限消去から
式\eqref{cmo:eq:kappa}を導いて六十四個の自己 Ext を全次数で決定する。
一般の分解構成や周期性自体の新規性は主張せず、
この具体式と自己 Ext 計算の包括的な新規性確認も未完了である。

$p=(a,b,c)\in(\C^\times)^3$ とし、
$\alpha=(a^2,b^2,c^2)$、
$D_p=\mathcal B_\alpha/(u-1)$ と置く。
$M_p=\C$ は $(x,y,z)$ が $(a,b,c)$ として作用する左 $D_p$-加群である。
式\eqref{cmo:eq:mixed-center}の $u=1$ への特殊化を $\Omega_p$ と書き、
\begin{equation}\label{cmo:eq:character-value}
 \lambda_p=-\frac23(a^3+b^3+c^3+abc),\qquad
 f_p=\Omega_p-\lambda_p,\qquad R_p=D_p/(f_p)
\end{equation}
とする。$f_p$ は $M_p$ を消す正則中心元である。
ここではこの中心元の正規化を固定する。

\begin{proposition}[character における中心三次の二次作用]
\label{cmo:prop:character-ext}
次の有理式を置く。
\begin{equation}\label{cmo:eq:kappa}
 \kappa(a,b,c)=
 \frac{2\bigl((ab)^3+(bc)^3+(ca)^3+9a^2b^2c^2\bigr)}{9abc}.
\end{equation}
自己 Ext を計算する cochain 複体として
\begin{equation}\label{cmo:eq:character-complex}
 \operatorname{RHom}_{R_p}(M_p,M_p)
  \simeq
  \bigl(\C[t]\oplus\C[t]\eta,\quad
        |t|=2,\ |\eta|=3,\ d(t)=0,\ d(\eta)=\kappa(p)t^2\bigr)
\end{equation}
がある。この表示の右辺の乗法を Yoneda 積と同定する主張は含めない。
従って $\kappa(p)\ne0$ なら
\[
 \operatorname{Ext}_{R_p}^j(M_p,M_p)=
 \begin{cases}\C,&j=0,2,\\0,&\text{その他},\end{cases}
\]
である。$\kappa(p)=0$ なら、自己 Ext は次数 $2q$ と $2q+3$
（$q\geq0$）に各一次元あり、上に有界でない。
中心方程式を $s f_p$、$s\in\C^\times$ に変更したとき、
同じ Shamash 変数の規約では式\eqref{cmo:eq:character-complex}の係数は
$s^2\kappa(p)$ となる。
\end{proposition}

\begin{proof}
$X=x-a,Y=y-b,Z=z-c$ と置く。関係式は
\[
 [Y,Z]=X^2+2aX,\qquad [Z,X]=Y^2+2bY,\qquad
 [X,Y]=Z^2+2cZ
\]
となる。左自由加群の係数を行として書き、写像の行列を右から掛ける規約で、
\begin{equation}\label{cmo:eq:character-resolution}
 0\longrightarrow D_p\xrightarrow{d_3}D_p^3
   \xrightarrow{d_2}D_p^3\xrightarrow{d_1}D_p
   \longrightarrow M_p\longrightarrow0,
 \quad
 \begin{gathered}
 d_1=(X,Y,Z)^{\mathsf t},\qquad d_3=(X,Y,Z),\\
 d_2=\begin{pmatrix}
  -X-2a&-Z&Y\\ Z&-Y-2b&-X\\ -Y&X&-Z-2c
 \end{pmatrix}
 \end{gathered}
\end{equation}
を得る。$d_2d_1=0$ と $d_3d_2=0$ は表示関係そのものである。
第 $i$ 項を $i$ だけずらした通常の次数濾過を取ると、associated graded は
$S(1,-1,-1)$ の augmentation の Koszul 分解である。
例\ref{cmo:ex:sixty-four}の PBW 論証と Koszul 性により、
この有界下濾過複体も完全である。
$\Omega_p$ の最高次項は正則中心元 $\Omega_0$ なので、
同じ濾過から $f_p$ の正則性を得る。
中心性は、以下の普遍な有限証明書でも同時に確認する。

$P_\bullet$ を式\eqref{cmo:eq:character-resolution}の自由分解とする。
$f_p$ の乗法に対する higher homotopies を
\[
 h_{1,i}:P_i\longrightarrow P_{i+1}\quad(0\leq i\leq2),
 \qquad h_2:P_0\longrightarrow P_3
\]
と書く。$\mathbb Q[a,b,c]$ 上で、係数がそれぞれ全次数二、三の行列を取り、
\begin{equation}\label{cmo:eq:character-homotopy}
 dh_1+h_1d=f_p\operatorname{id},\qquad
 dh_2+h_2d+h_1^2=0
\end{equation}
を満たすようにできる。長さ三の分解なので、これより高い恒等式は
写像先の次数の理由で零になる。
ここで使う具体的な多項式行列の全係数と、その有限証明書の検算を
\texttt{computations/verify\_mixed\_character\_ext.py} に収めた。
三つの中心パラメータも次数一とし、全次数 $2,3,4$ の全語に対し、
各関係の全左右語倍と全パラメータ倍を列挙する。
その有理係数関係行列で、$d^2=0$、$[f_p,X]=[f_p,Y]=[f_p,Z]=0$、
式\eqref{cmo:eq:character-homotopy}を全て零へ還元する。
この検算は有限個のパラメータの代入や補間ではなく、
$\mathbb Q[a,b,c]$ 上の多項式恒等式である。
全次数での分解の完全性は前段の PBW/Koszul 論証で別に保証した。

この証明書の character での値だけを記すと、
\begin{equation}\label{cmo:eq:character-homotopy-values}
 \begin{gathered}
 d_2(p)=\operatorname{diag}(-2a,-2b,-2c),\qquad h_{1,1}(p)=0,\\
 v=h_{1,0}(p)=-\frac23(bc,ac,ab+c^2),\\
 z=h_{1,2}(p)=\frac23(a^2-bc,-ac-b^2,-ab)^{\mathsf t},\qquad
 h_2(p)=\frac{16}{9}abc.
 \end{gathered}
\end{equation}
補題\paperref{main}{mbnp:lem:shamash}の分解構成は、内部次数を忘れた正則中心元と
任意の消される加群にも同じ証明で適用できる。
すなわち $\bigoplus_{q\geq0}(R_p\otimes_{D_p}P_{n-2q})$ に
$d,h_1,h_2$ を配置する。添字 $q$ による濾過の第一頁は
二項中心分解から $M_p$ を二つ持ち、$dh_1+h_1d=f_p$ による連結射が
隣接する項を恒等で相殺する。各全次数の濾過は有限なので、
これは $M_p$ の自由 $R_p$-分解となる。

この分解へ $\operatorname{Hom}_{R_p}(-,M_p)$ を施すと、
\[
 \left(\operatorname{Hom}_{D_p}(P_\bullet,M_p)[t],
       d^*+t h_1^*+t^2h_2^*\right)
\]
を得る。式\eqref{cmo:eq:character-homotopy-values}の可逆な $d_2(p)$ を
$\C[t]$ 上で消去する。$h_{1,1}(p)=0$ なのでこの消去は有限の行列操作であり、
形式冪級数への完備化を必要としない。
残る次数零と三の二つの基底間の係数は
\[
 \left(h_2(p)-v\,d_2(p)^{-1}z\right)t^2
   =\kappa(p)t^2
\]
である。最後の等式は式\eqref{cmo:eq:character-homotopy-values}の代入による
有理式恒等式であり、同じコードで記号的にも確認する。
これが式\eqref{cmo:eq:character-complex}を与える。
$\kappa\ne0$ なら残る cohomology は $\C[t]/(t^2)$、
$\kappa=0$ なら微分は零である。
また中心元を $s$ 倍すると式\eqref{cmo:eq:character-homotopy}では
$h_1$ を $s$ 倍、$h_2$ を $s^2$ 倍するので、係数は $s^2\kappa$ となる。
\end{proof}

\begin{corollary}[六十四例外点の自己 Ext は有界である]
\label{cmo:cor:sixty-four-ext}
例\ref{cmo:ex:sixty-four}の二ブロックについて、
$p=(\pm1,\pm2,\pm3)$、$p'=(\pm4,\pm5,\pm7)$ とし、
$c_0=-\lambda_p-100\lambda_{p'}$ と置く。
中心 chart の代数
\[
 D_{c_0}=(A_{c_0}[u^{-1}])_0
    \simeq (D_p\otimes D_{p'})/(f_p+100f_{p'})
\]
の character $M=M_p\otimes M_{p'}$ に対して、
\begin{equation}\label{cmo:eq:sixty-four-ext}
 \operatorname{Ext}_{D_{c_0}}^j(M,M)=
 \begin{cases}\C,&j=0,2,3,5,\\0,&\text{その他}\end{cases}
\end{equation}
である。従って全六十四例外値について、この定常 character の自己 Ext は
商の非滑らかさを示す「無限次数の自己 Ext」という障害を与えない。
\end{corollary}

\begin{proof}
二因子の式\eqref{cmo:eq:character-resolution}を tensor し、
homological Koszul 符号付きに higher homotopies を足す。
異なる因子の奇写像は反交換するため、
中心 $f_p+100f_{p'}$ の higher homotopies は、第二因子の
$h_1$ を $100$ 倍、$h_2$ を $10000$ 倍して得られる。
先の $\C[t]$ 上の有限消去を各因子に施すと、
自己 Ext を計算する複体は、次数付きベクトル空間として
\begin{equation}\label{cmo:eq:two-character-complex}
 \left(\C[t]\otimes\bigwedge(\eta_1,\eta_2),\quad
 |t|=2,\ |\eta_i|=3,\quad
 d(\eta_1)=\kappa(p)t^2,\quad
 d(\eta_2)=10000\kappa(p')t^2\right)
\end{equation}
となる。外積と Leibniz 則はこの計算複体の記法であり、
自己 Ext の Yoneda 積の同定は必要ない。

式\eqref{cmo:eq:kappa}の分子で、全ての符号の選び方に対して
\[
 9a^2b^2c^2=324>8+216+27=251
\]
が成立する。第二因子でも
\[
 9a^2b^2c^2=176400>8000+42875+21952=72827
\]
なので、両方の $\kappa$ は零でない。
$\eta=\eta_1/\kappa(p)$ と
$\zeta=\eta_2-10000\kappa(p')\eta_1/\kappa(p)$ を取ると
$d\eta=t^2,d\zeta=0$ である。従って全次数の cohomology は
$\C[t]/(t^2)\otimes\bigwedge(\zeta)$ となり、
式\eqref{cmo:eq:sixty-four-ext}を得る。
全十六因子 character と六十四組の係数・有限行列の検算も
上記再現コードに含めた。全次数の消滅は、この基底変換による証明であり、
有限次数までの計算からの外挿ではない。
\end{proof}

\begin{remark}[自己 Ext 判定の射程]
\label{cmo:rem:self-ext-scope}
可換化の norm $N(c)$ は定常 character の存在を検出する。
有限自己 Ext から有限射影次元を結論する一般論は使えないが、
以下では有限 ambient 分解と二周期性を追加して射影次元まで決定する。
その結論も、全加群の有限大域次元や全 coherent 圏の smoothness の判定とは異なる。
また計算複体に書いた外積を Yoneda 積と同一視する主張は含めない。
\end{remark}

\subsection{二次元単純加群の分類と中心固有値}
\label{cmo:subsec:two-dimensional}

定常 character が検出しない有限次元単純加群を、最小の次元で調べる。
本節では $D_\alpha=\mathcal B_\alpha/(u-1)$ とし、
$\Omega_\alpha$ は式\eqref{cmo:eq:mixed-center}の $u=1$ 特殊化を表す。
$\alpha=0$ における二次元既約表現の存在は
\cite[Proposition~4.1]{CMOWaltonRepresentations}で既知である。
以下は定数項を含む本表示について、全ての二次元単純加群と
指定中心元の値を直接計算するものであり、分類公式の新規性は未確定である。

\begin{proposition}[半跡による全二次元単純加群の分類]
\label{cmo:prop:two-dimensional-classification}
次の二次写像と三次式を用いる。
\[
 q(a,b,c)=(a^2-bc,b^2-ca,c^2-ab),\qquad
 F(a,b,c)=a^3+b^3+c^3-3abc.
\]
二次元単純 $D_\alpha$ 加群の同型類は
\begin{equation}\label{cmo:eq:two-dimensional-fibre}
 \{(a,b,c)\in(\C^\times)^3:q(a,b,c)=\alpha\}
\end{equation}
と一対一に対応する。
対応する表現の $x,y,z$ の半跡はそれぞれ $a,b,c$ であり、
$r^2=-bc$ を選ぶと代表は
\begin{equation}\label{cmo:eq:two-dimensional-matrices}
 \begin{split}
 X&=aI+\begin{pmatrix}r&0\\0&-r\end{pmatrix},\\
 Y&=bI+\begin{pmatrix}0&1\\-ca&0\end{pmatrix},\qquad
 Z=cI+\begin{pmatrix}0&r/c\\ra&0\end{pmatrix}
 \end{split}
\end{equation}
で与えられる。二つの $r$ の選択は同型な表現を与える。

特に $D=F(\alpha)\ne0$、$Q=q(\alpha)$ とすると、
二次元単純加群が存在するための必要十分条件は
$Q_1Q_2Q_3\ne0$ である。この場合、ちょうど二つの同型類があり、
それらの半跡は
\begin{equation}\label{cmo:eq:two-dimensional-inversion}
 (a,b,c)=\frac{Q}{\delta},\qquad \delta^2=D
\end{equation}
である。
\end{proposition}

\begin{proof}
$X=aI+U,Y=bI+V,Z=cI+W$ と分解し、
$U,V,W\in\mathfrak{sl}_2(\C)$ とする。
対称双線型形式を $\langle U,V\rangle=\frac12\operatorname{tr}(UV)$ と置く。
二次 Cayley--Hamilton 恒等式から
$U^2=\langle U,U\rangle I$ であり、関係式のスカラー部分と跡零部分は
\begin{equation}\label{cmo:eq:two-dimensional-traceless}
 \begin{gathered}
 \langle U,U\rangle=\alpha_1-a^2,\quad
 \langle V,V\rangle=\alpha_2-b^2,\quad
 \langle W,W\rangle=\alpha_3-c^2,\\
 [V,W]=2aU,\qquad [W,U]=2bV,\qquad[U,V]=2cW
 \end{gathered}
\end{equation}
となる。
$\operatorname{span}(U,V,W)$ は $\mathfrak{sl}_2$ の Lie 部分代数である。
真部分代数は可解であるから、Lie の定理により同時三角化される。
従って既約表現では $U,V,W$ は $\mathfrak{sl}_2$ の基底をなす。
非零な跡零二次行列の $\mathfrak{sl}_2$ 内の中心化部分代数は一次元なので、
式\eqref{cmo:eq:two-dimensional-traceless}から $abc\ne0$ が従う。
跡形式の不変性から $U,V,W$ は互いに直交する。

跡零二次行列の恒等式
\[
 [[U,V],V]
 =4\bigl(\langle V,V\rangle U-\langle U,V\rangle V\bigr)
\]
と式\eqref{cmo:eq:two-dimensional-traceless}を合わせると
\[
 U^2=-bcI,\qquad V^2=-caI,\qquad W^2=-abI
\]
を得る。従って $q(a,b,c)=\alpha$ である。

逆に式\eqref{cmo:eq:two-dimensional-fibre}を満たす点に対し、
式\eqref{cmo:eq:two-dimensional-matrices}の行列は
式\eqref{cmo:eq:two-dimensional-traceless}を直接満たす。
$r\ne0$ であり、$U$ の二つの固有直線は $V$ の不変直線ではないから既約である。
また任意の既約表現で $U$ を対角化すると、直交性により $V$ は非対角成分だけを持つ。
$V^2=-caI$ より両非対角成分が非零で、対角行列による共役で上右成分を $1$ にできる。
このとき下左成分は $-ca$、$W=[U,V]/(2c)$ となり、表示された形を得る。
$r$ の符号の変更は $V$ による共役で実現される。
半跡は共役不変なので、異なる $(a,b,c)$ は非同型である。

多項式恒等式
\begin{equation}\label{cmo:eq:quadratic-involution}
 q(q(a,b,c))=F(a,b,c)(a,b,c),\qquad
 F(q(a,b,c))=F(a,b,c)^2
\end{equation}
から、$q(a,b,c)=\alpha$、$F(\alpha)\ne0$ なら
$F(a,b,c)=\delta$、$\delta^2=D$、$(a,b,c)=Q/\delta$ となる。
逆にこの式により $q(Q/\delta)=q(Q)/D=\alpha$ である。
非零座標の条件は $Q_1Q_2Q_3\ne0$ と同値である。
\end{proof}

\begin{corollary}[退化係数を含む存在条件]
\label{cmo:cor:two-dimensional-boundary}
$F(\alpha)=0$ とする。
$q(\alpha)\ne0$ なら二次元単純 $D_\alpha$ 加群は存在しない。
$q(\alpha)=0$ なら存在し、次のように分類される。
\begin{enumerate}
\item $\alpha=0$ のとき、同型類の空間は
$(a,b,c)=(s,\zeta s,\zeta^2s)$、
$s\in\C^\times$、$\zeta^3=1$ による三つの $\C^\times$ の非交和である。
\item $\alpha\ne0$ のとき、$\alpha=h(1,\zeta,\zeta^2)$
（$h\ne0$, $\zeta^3=1$）であり、同型類の空間は
一つの $\C^\times$ から六つの相異なる点を除いた曲線である。
\end{enumerate}
\end{corollary}

\begin{proof}
式\eqref{cmo:eq:quadratic-involution}より
$F(\alpha)=0$ を満たす解では $F(a,b,c)=0$、従って $q(\alpha)=0$ である。
$q(\alpha)=0$ の非零解は表示された三本の直線上の点である。
原始三乗根 $\omega$ を固定し、
\[
 R=a+b+c,\quad S=a+\omega b+\omega^2c,\quad
 T=a+\omega^2b+\omega c
\]
と置く。$q(a,b,c)=\alpha$ は
\[
 ST=\alpha_1+\alpha_2+\alpha_3,\quad
 RS=\alpha_1+\omega\alpha_2+\omega^2\alpha_3,\quad
 RT=\alpha_1+\omega^2\alpha_2+\omega\alpha_3
\]
に等しい。右辺が全て零なら、$R,S,T$ の座標軸の非零部分を得る。
右辺の一つだけが非零なら、二変数の積を非零定数に固定し、残る変数を零とする双曲線を得る。
例えば $\alpha=(h,h,h)$ では $R=0$、$ST=3h$ であり、
$abc=(S^3+T^3)/27$ だから、除かれる点は $S^6=-27h^3$ の六解である。
他の二直線も Fourier 座標の巡回置換で同じ結論となる。
\end{proof}

\begin{proposition}[中心固有値と既約二次元加群の剛性]
\label{cmo:prop:two-dimensional-center-rigidity}
命題\ref{cmo:prop:two-dimensional-classification}の加群 $M_{a,b,c}$ において、
$\Omega_\alpha$ は次のスカラーで作用する。
\begin{equation}\label{cmo:eq:two-dimensional-lambda}
 \lambda_2(a,b,c)=-\frac23(a^3+b^3+c^3+abc).
\end{equation}
$D=F(\alpha)\ne0$、$Q=q(\alpha)$、$Q_1Q_2Q_3\ne0$ のとき、二つの値は
\begin{equation}\label{cmo:eq:two-dimensional-lambda-polynomial}
 \lambda_2=-\frac{2(D^2+4Q_1Q_2Q_3)}{3D\delta},\quad \delta^2=D,
 \qquad
 \lambda_2^2=\frac{4(D^2+4Q_1Q_2Q_3)^2}{9D^3}
\end{equation}
を満たす。分子が零なら、二つの非同型単純加群の中心値はともに零である。
また
\begin{equation}\label{cmo:eq:two-dimensional-self-ext}
 \operatorname{Ext}_{D_\alpha}^i(M_{a,b,c},M_{a,b,c})
 \cong
 \begin{cases}
 \C&i=0,3,\\
 \C^{e}&i=1,2,\\
 0&i>3,
 \end{cases}
 \qquad
 e=\begin{cases}0&F(\alpha)\ne0,\\1&F(\alpha)=0\end{cases}
\end{equation}
がベクトル空間として成立する。
\end{proposition}

\begin{proof}
中心性と Schur の補題により、中心値は半跡で計算できる。
前命題で得た積の公式から
\[
 \tfrac12\operatorname{tr}(X^3)=a^3-3abc,\qquad
 \tfrac12\operatorname{tr}\bigl(\sym(XYZ)\bigr)=abc,
 \qquad
 \alpha_1a+\alpha_2b+\alpha_3c=a^3+b^3+c^3-3abc
\]
となる。これを中心元へ代入すると式\eqref{cmo:eq:two-dimensional-lambda}を得る。
また $F(Q)=D^2$ より
$Q_1^3+Q_2^3+Q_3^3+Q_1Q_2Q_3=D^2+4Q_1Q_2Q_3$ であり、
式\eqref{cmo:eq:two-dimensional-lambda-polynomial}が従う。

一次変形を双対数環 $A=\C[\varepsilon]/(\varepsilon^2)$ 上で考える。
$abc$ と $r$ は単元なので、前命題の正規形の構成はこの環上でも成立する。
実際、$U$ の二つの固有値は差が単元であり、その固有基底を
特殊化時の基底から持ち上げられる。
直交性により $V$ は非対角形となり、その上右成分を単元の対角共役で $1$ にできる。
従って恒等写像へ特殊化する共役で割った一次変形は、
$q(a+\varepsilon\dot a,b+\varepsilon\dot b,c+\varepsilon\dot c)=\alpha$ と一対一に対応する。
これは
\[
 \operatorname{Ext}^1_{D_\alpha}(M_{a,b,c},M_{a,b,c})
 \cong\ker
 \begin{pmatrix}
 2a&-c&-b\\-c&2b&-a\\-b&-a&2c
 \end{pmatrix}
\]
を与える。この行列式は $-2F(a,b,c)$ である。
$abc\ne0$ では階数は少なくとも二である。
実際、階数一以下なら三つの主小行列から
$4ab=c^2,4bc=a^2,4ca=b^2$ となり、積を取ると $64=1$ を得る。
$F(\alpha)=F(a,b,c)^2$ により、表示した $e$ を得る。

本代数の三次 potential による長さ三の bimodule 分解は、
定数項の PBW 変形でも自己双対であり、$D_\alpha$ は三次元 Calabi--Yau である
\cite[Theorem~3.3.2]{NFMEG}。
その有限次元加群間の双対性
$\operatorname{Ext}^{3-i}(M,M)\cong\operatorname{Ext}^{i}(M,M)^*$、
Schur の補題、および同分解から、残る次数の主張を得る。
\end{proof}

\begin{example}[六十四 character の表示における二次元加群]
\label{cmo:ex:two-dimensional-sixty-four}
例\ref{cmo:ex:sixty-four}の各因子は、ちょうど二つずつの二次元単純加群を持つ。
中心固有値の集合を $\mathcal T,\mathcal T'$ とすると
\[
 \mathcal T=\left\{\pm\frac{192\sqrt{14}}{49}\right\},\qquad
 \mathcal T'=\left\{\pm\sqrt{\frac{37623589339922432}{748505944125}}\right\}.
\]
実際、$(D,Q)=(686,(-35,7,77))$、
$(78570,(-969,-159,2001))$ である。
これらの加群はいずれも周囲代数 $D_\alpha$ または $D_{\alpha'}$ 上で剛的である。

二つの因子の一次元または二次元単純加群のテンソル積は、
$D_\alpha\otimes D_{\alpha'}$ の単純加群である。
その次元は $1,2,4$ であり、中心商
$(D_\alpha\otimes D_{\alpha'})/(\Omega_\alpha+100\Omega_{\alpha'}+c)$
へ降りる条件は
\[
 c\in-(\mathcal R\cup\mathcal T)-100(\mathcal S\cup\mathcal T')
\]
である。ここで $\mathcal R,\mathcal S$ は式\eqref{cmo:eq:explicit-norm}の集合である。
これは単純加群の出現値の計算であり、この集合を非滑らかな中心値の集合と同一視しない。
二次元因子の剛性だけから中心商上の射影次元を結論することはできない。
命題\ref{cmo:prop:one-two-dimensional-pd}では、指定中心方程式の高次作用を
別に計算し、ここに現れる全百個の加群の射影次元を決定する。
\end{example}

\subsection{周期性を使った character の射影次元判定}

有限自己 Ext だけから有限射影次元を結論することは一般にはできない。
ここでは有限 ambient 分解から得られる二周期性と、中心商の Rees 同型を
追加してこの推論を正当化する。以下の一般論自体の新規性は主張しない。
Shamash 型分解は補題\paperref{main}{mbnp:lem:shamash}の証明を用いる。
非可換な正則 normal 商の周期分解については Cassidy--Conner--Kirkman--Moore
\cite[Theorem~4.2]{CMOPeriodicMF} も関連する。
ただし同論文の主定理は連結次数付きの設定なので、以下の affine 版は直接証明する。

\begin{lemma}[二周期性による有限射影次元の検出]
\label{cmo:lem:periodic-pd}
$R$ を環、$M$ を有限生成左 $R$-加群とする。
$M$ が有限生成射影加群による分解を持ち、その十分高い syzygy が
その分解の中で実際に二周期であると仮定する。
さらに $\operatorname{Ext}_R^i(M,R)=0$ が十分大きい $i$ で成立するとする。
このとき
\[
 \operatorname{pd}_R M<\infty
 \quad\Longleftrightarrow\quad
 \operatorname{Ext}_R^i(M,M)=0\quad(i\gg0).
\]
また $d\ge1$ とし、$\operatorname{Ext}_R^i(M,R)=0$ が $i\ge d$ で成立するなら、
有限の場合には $\operatorname{pd}_R M\le d-1$ である。
\end{lemma}
\begin{proof}
順方向は定義から従う。逆方向を示す。
$K=\Omega_R^n M$ を二周期部分に入った syzygy とする。
$M$ の分解の第 $i$ 項を $P_i$ と書く。
各 $P_i$ は有限生成射影なので、$\operatorname{Ext}_R^{\gg0}(M,P_i)=0$ である。
十分大きい $r$ に対して、第一変数、次いで第二変数の dimension shifting から
\[
 \operatorname{Ext}_R^{2r}(K,K)
 \simeq\operatorname{Ext}_R^{2r+n}(M,K)
 \simeq\operatorname{Ext}_R^{2r}(M,M)=0
\]
を得る。第二同型は有限個の完全列
$0\to\Omega_R^{i+1}M\to P_i\to\Omega_R^iM\to0$
に $\operatorname{Hom}_R(M,-)$ を適用して得られる。

$K$ の二周期分解における同型
$\theta:\Omega_R^{2r}K\xrightarrow{\sim}K$ を取る。
標準の表示
\[
 \operatorname{Ext}_R^{2r}(K,K)
 =\operatorname{coker}\!\left(
  \operatorname{Hom}_R(P_{n+2r-1},K)
  \longrightarrow\operatorname{Hom}_R(\Omega_R^{2r}K,K)\right)
\]
において $[\theta]=0$ なので、$\theta$ は $P_{n+2r-1}\to K$ へ延長する。
この延長と $\theta^{-1}$ の合成は
$\Omega_R^{2r}K\hookrightarrow P_{n+2r-1}$ の retraction となる。
従って $K\simeq\Omega_R^{2r}K$ は射影で、$M$ の射影次元は有限である。

最後に $m=\operatorname{pd}_R M<\infty$ とする。
もし $m>0$ かつ $\operatorname{Ext}_R^m(M,R)=0$ なら、長さ $m$ の
有限生成射影分解を双対化した最後の射
$P_{m-1}^{\vee}\to P_m^{\vee}$ は全射である。
$P_m^{\vee}$ は有限生成右射影加群だからこの全射は分裂する。
有限生成射影加群の bidual 同型を使うと
$P_m\to P_{m-1}$ が分裂単射となり、分解を短縮できて矛盾する。
従って $m>0$ なら $\operatorname{Ext}_R^m(M,R)\ne0$。
これより $m\le d-1$ を得る。
\end{proof}

\begin{corollary}[有限 ambient 分解を持つ中心超曲面加群]
\label{cmo:cor:hypersurface-pd}
$Q$ を環、$f\in Z(Q)$ を左右正則元、$R=Q/(f)$ とする。
$M\ne0$ を左 $R$-加群とし、$Q$ 上で長さ $d\ge1$ の
有限生成射影分解を持つと仮定する。
このとき前補題の同値と有限の場合の上限 $\operatorname{pd}_R M\le d-1$ が成立する。
左右を交換した主張も成立する。
\end{corollary}
\begin{proof}
中心元の二項分解から
\[
 \operatorname{RHom}_Q(R,Q)\simeq R[-1],\qquad
 \operatorname{RHom}_Q(M,Q)\simeq
 \operatorname{RHom}_R(M,R)[-1]
\]
を得る。すなわち Rees 同型
\[
 \operatorname{Ext}_R^i(M,R)\simeq
 \operatorname{Ext}_Q^{i+1}(M,Q)=0\quad(i\ge d)
\]
がある。ここで正則元の中央性により normalizing twist は現れない。

$Q$ 上の有限分解 $P_\bullet$ に $f$ の higher homotopies $s_a$ を選ぶ。
$s_a:P_j\to P_{j+2a-1}$ なので非零な $s_a$ は有限個である。
補題\paperref{main}{mbnp:lem:shamash}と同じ濾過証明により
\[
 C_n=\bigoplus_{q\ge0}(R\otimes_QP_{n-2q})
\]
は $M$ の有限生成射影分解になる。
十分大きい偶数 $n$ では $C_n=\bigoplus_{j\text{ even}}R\otimes_QP_j$、
十分大きい奇数 $n$ では奇数 $j$ の直和となり、微分も parity のみに依存する。
従ってこれは実際の二周期分解である。
ここでは内部次数を忘れた affine 加群を扱うので、内部次数 shift を
同型と取り違えているのではない。前補題が適用できる。
\end{proof}

\begin{proposition}[混合 character の射影次元の完全判定]
\label{cmo:prop:character-pd}
$p_i=(a_i,b_i,c_i)\in(\C^\times)^3$、$\beta_i\in\C^\times$ とし、
命題\ref{cmo:prop:character-ext}の記法で
\[
 Q=\bigotimes_{i=1}^kD_{p_i},\quad
 R=Q\Big/\Big(\sum_{i=1}^k\beta_if_{p_i}\Big),\quad
 M=\bigotimes_{i=1}^kM_{p_i}
\]
と置く。中心和は正則とする。この仮定は本節の混合族では PBW 濾過から成立する。
このとき
\[
 \operatorname{pd}_R M=
 \begin{cases}
  3k-1,&\text{少なくとも一つの }\kappa(p_i)\ne0,\\
  \infty,&\text{全ての }\kappa(p_i)=0.
 \end{cases}
\]
特に一因子の $\kappa\ne0$ では射影次元は $2$、
例\ref{cmo:ex:sixty-four}の六十四個の affine character では全て $5$ である。
\end{proposition}
\begin{proof}
各因子の三段自由分解を tensor すれば、$M$ は $Q$ 上で長さ $3k$ の
有限自由分解を持つ。二因子の証明と同じ有限消去により、自己 Ext の計算複体は
\[
 \C[t]\otimes\bigwedge(\eta_1,\ldots,\eta_k),\quad
 |t|=2,\quad|\eta_i|=3,\quad d\eta_i=\beta_i^2\kappa(p_i)t^2
\]
となる。全ての $\kappa(p_i)$ が零なら微分は零で、自己 Ext は無限次数にある。
従って射影次元は無限である。
少なくとも一つが非零なら、線型基底変換で
$d\eta=t^2,d\zeta_1=\cdots=d\zeta_{k-1}=0$ となる。
cohomology は次数付きベクトル空間として
\[
 \C[t]/(t^2)\otimes\bigwedge(\zeta_1,\ldots,\zeta_{k-1})
\]
であり、上に有界で最高非零次数は $3k-1$ である。
前系より射影次元は有限で $3k-1$ 以下、最高自己 Ext の非消滅より以上でもある。
\end{proof}

\subsection{二次元加群に対する中心方程式の高次作用}

二次元加群の中心商上の射影次元を判定するため、
character の証明書を形式的に代入する代わりに、その加群の自由分解を用いる。
以下では自由左加群の射を行行列で記述する。
表現行列を $X,Y,Z$ とし、
\[
 A=xI-X^{\mathsf t},\qquad B=yI-Y^{\mathsf t},\qquad C=zI-Z^{\mathsf t}
\]
と置く。$M$ には基底を選び、その双対基底で自由加群の行成分を添字付ける。

\begin{lemma}[行列係数の長さ三分解と高次作用の抽出]
\label{cmo:lem:two-dimensional-shamash}
任意の二次元 $D_\alpha$ 加群 $M$ に対し、
\begin{equation}\label{cmo:eq:two-dimensional-resolution}
 0\longrightarrow D_\alpha^2\xrightarrow{d_3}D_\alpha^6
 \xrightarrow{d_2}D_\alpha^6\xrightarrow{d_1}D_\alpha^2
 \longrightarrow M\longrightarrow0,
\end{equation}
\[
 d_1=\begin{pmatrix}A\\B\\C\end{pmatrix},\quad
 d_2=\begin{pmatrix}
 -xI-X^{\mathsf t}&-C&B\\
 C&-yI-Y^{\mathsf t}&-A\\
 -B&A&-zI-Z^{\mathsf t}
 \end{pmatrix},\qquad
 d_3=(A\ B\ C)
\]
は自由分解である。

$M=M_{a,b,c}$ が命題\ref{cmo:prop:two-dimensional-classification}の
$F(\alpha)\ne0$ を満たす単純加群で、
$f=\Omega_\alpha-\lambda_2(a,b,c)$ とする。
PBW 性と斉次 Sklyanin 代数の整域性より $f$ は正則中心元である。
次の恒等式を満たす行列 $h_0,h_1,h_2,g$ を取る。
\begin{equation}\label{cmo:eq:two-dimensional-homotopies}
 \begin{gathered}
 h_0d_1=fI_2,\qquad d_1h_0+h_1d_2=fI_6,\\
 d_2h_1+h_2d_3=fI_6,\qquad d_3h_2=fI_2,\\
 gd_3+h_0h_1=0,\qquad d_1g+h_1h_2=0.
 \end{gathered}
\end{equation}
行列の大きさは順に $2\times6,6\times6,6\times2,2\times2$ である。
各成分を $M$ に作用させた複素行列を
$D_i,H_j,G$ と書き、$v=(1,0,0,1)^{\mathsf t}$ とする。
$D_2b=H_2v$ を満たす $b\in\C^{12}$ に対し
\begin{equation}\label{cmo:eq:two-dimensional-kappa-extraction}
 Gv-H_0b=\kappa_2v
\end{equation}
で $\kappa_2\in\C$ を定める。このスカラーは $b$ の選択に依存しない。
また固定した $f$ と両端の規格化の下で、高次 homotopy と収縮の選択にも依存しない。
この規格化では、$D_\alpha/(f)$ 上の自己 Ext を計算する小複体は
\begin{equation}\label{cmo:eq:two-dimensional-effective}
 \C[T]\oplus\C[T]\eta,\qquad |T|=2,\quad|\eta|=3,\quad
 d\eta=\kappa_2T^2
\end{equation}
である。
半跡と表現を同時に $s\ne0$ 倍し、係数を $s^2\alpha$ とすると、
この規格化の係数は $s^3\kappa_2$ になる。
\end{lemma}

\begin{proof}
表示された行列について、$d_2d_1=d_3d_2=0$ は
$D_\alpha$ と $X,Y,Z$ の関係式から従う。
例えば $d_2d_1$ の最初の $2\times2$ ブロックは
\[
 -x^2I+(X^{\mathsf t})^2+[y,z]I+[Y^{\mathsf t},Z^{\mathsf t}]
 =-\alpha_1I+(X^2-[Y,Z])^{\mathsf t}=0
\]
である。各自由加群をその homological degree で濾過シフトすると、
随伴次数付き複体は斉次 Sklyanin 代数の Koszul 分解の二直和になる。
PBW 性と濾過の有界下性により式\eqref{cmo:eq:two-dimensional-resolution}は完全である。

$\operatorname{Hom}_{D_\alpha}(-,M)$ を適用した複体の大きさは
$4,12,12,4$ であり、その微分は $D_1,D_2,D_3$ である。
命題\ref{cmo:prop:two-dimensional-center-rigidity}から階数は $3,9,3$ となる。
最初と最後の四次元空間は $\operatorname{End}_{\C}(M)$ と同一視する。
$v$ は恒等行列を表し、$\ker D_1=\C v$、
$\operatorname{im}D_3$ は跡零行列の空間である。
従って $v$ は $H^0,H^3$ の生成元を指定する。
特に $H^3$ の規格化は商への射影 $U\mapsto\operatorname{tr}(U)/2$ である。

式\eqref{cmo:eq:two-dimensional-homotopies}は、長さ三の分解上の全高次 homotopy 恒等式である。
補題\paperref{main}{mbnp:lem:shamash}から中心商上の自由分解を得る。
その Hom 複体の微分は、周囲微分に $TH_j$ と $T^2G$ を加えたものである。
まず $H_2v$ は $\ker D_3=\operatorname{im}D_2$ に属し、$b$ が存在する。
また $D_1(Gv-H_0b)=0$ なので、式\eqref{cmo:eq:two-dimensional-kappa-extraction}が成立する。
$b$ を変えたときの差は $\ker D_2=\operatorname{im}D_1$ に属し、
$H_0D_1=0$ なので係数は変わらない。

有界な周囲複体を $H^0\oplus H^3$ に収縮する。
$H_j$ は元の cochain degree を一つ、$G$ は三つ下げるため、
収縮後の非零な微分は $H^3\to T^2H^0$ だけであり、
二段の消去式はちょうど $Gv-H_0b$ である。
収縮の補正級数は元の cochain degree を下げるので有限であり、
これは $\C[T]$ 上の homotopy 同値を与える。
従って式\eqref{cmo:eq:two-dimensional-effective}は全次数の計算である。
小複体の外積表示を Yoneda 積の主張とはしない。
この係数は中心商の change-of-rings spectral sequence において
$\operatorname{Ext}^3_{D_\alpha}(M,M)$ から
$T^2\operatorname{Ext}^0_{D_\alpha}(M,M)$ への最初に許される微分の係数である。
従って $f$ による周期変数 $T$ と両端の生成元を固定すれば、
高次 homotopy や収縮を変えても同じ係数になる。
実際、二つの高次 homotopy 系の比較射は $T$ による濾過を保ち、
$E_1=\operatorname{Ext}_{D_\alpha}^{\bullet}(M,M)[T]$ に恒等を誘導するので、
最初の非零微分を保存する。

最後に、座標変換 $x\mapsto sx,y\mapsto sy,z\mapsto sz$ では
$d_i$ は一次、$f$ は三次である。
式\eqref{cmo:eq:two-dimensional-homotopies}の $h_j$ は二次、$g$ は三次に変換できる。
式\eqref{cmo:eq:two-dimensional-kappa-extraction}における二項はともに三次に変換されるから、
$\kappa_2\mapsto s^3\kappa_2$ を得る。
表現の同時共役は両端の恒等行列を固定するので、この規格化に影響しない。
\end{proof}

\begin{proposition}[二次元加群の普遍な高次作用]
\label{cmo:prop:two-dimensional-universal-kappa}
$abcF(a,b,c)\ne0$、$\alpha=q(a,b,c)$ とし、
\begin{equation}\label{cmo:eq:two-dimensional-sextic}
 \begin{split}
 N_6(a,b,c)={}&a^6+b^6+c^6
   +66\bigl((ab)^3+(bc)^3+(ca)^3\bigr)\\
 &+26abc(a^3+b^3+c^3)+153a^2b^2c^2
 \end{split}
\end{equation}
と置く。
補題\ref{cmo:lem:two-dimensional-shamash}で規格化した高次作用は
\begin{equation}\label{cmo:eq:two-dimensional-universal-kappa}
 \boxed{\displaystyle
 \kappa_2(a,b,c)=-\frac{2N_6(a,b,c)}{9F(a,b,c)}}
\end{equation}
である。従って $M_{a,b,c}$ の
$D_\alpha/(\Omega_\alpha-\lambda_2)$ 上の射影次元は、
$N_6(a,b,c)\ne0$ なら二、$N_6(a,b,c)=0$ なら無限である。
\end{proposition}

\begin{proof}
まず $r=1,c=-1/b$ とした普遍表現を
$\mathbb Q(a,b)$ 上で考える。
再現コード \path{verify_mixed_two_dimensional_ext.py} の
\texttt{--universal} オプションは
式\eqref{cmo:eq:two-dimensional-resolution}と
式\eqref{cmo:eq:two-dimensional-homotopies}をこの有理関数体上で厳密に計算し、
全行列恒等式を関係式の全左右倍による剰余で検査する。
さらに微分の階数 $3,9,3$ と
式\eqref{cmo:eq:two-dimensional-kappa-extraction}を同じ体上で計算する。
具体的には
\[
 \begin{split}
 P(a,b)={}&a^6b^6-26a^4b^6+66a^3b^9-66a^3b^3+153a^2b^6\\
 &-26ab^9+26ab^3+b^{12}-66b^6+1
 \end{split}
\]
と置くと、得られる係数は
\[
 -\frac{2P(a,b)}
 {9b^3(ab+b^2-1)(a^2b^2-ab^3+ab+b^4+b^2+1)}
\]
である。これは式\eqref{cmo:eq:two-dimensional-universal-kappa}への
$c=-1/b$ の代入と、有理関数として恒等的に等しい。
この計算は有限個のパラメータでの補間ではない。

計算中に現れる補助的な分母の零点を含め、全ての
$abcF\ne0$ へ結論を延ばす理由を確認する。
$r^2=-bc$、$abcF$ を可逆にした係数環を $R$ とする。
式\eqref{cmo:eq:two-dimensional-matrices}は $R$ 上自由な階数二加群を与え、
式\eqref{cmo:eq:two-dimensional-resolution}も $R$ 上の相対 PBW 分解として完全である。
相対 Hom 複体の $H^0,H^3$ は恒等行列と半跡により自明な直線束であり、
全幾何点で $H^1=H^2=0$、微分の階数は $3,9,3$ で一定である。
従って係数環上で局所的に contraction を選べる。
中心方程式の null homotopy とその高次持上げも、この自由分解から局所的に取れる。
前補題の選択独立性により、これらから得た $\kappa_2$ は $\operatorname{Spec}R$ 上の正則関数として貼り合う。
一方、右辺 $-2N_6/(9F)$ も同じ基礎上で正則である。

$r\ne0$ なので全ての点で $s=1/r$ の同次変換が可能であり、
変換後は $r=1,c=-1/b$ となる。
前補題の三次変換則により、計算した有理関数の恒等式は $R$ の稠密開上で成立する。
$R$ は整域なので、二つの正則関数の一致は $\operatorname{Spec}R$ 全体へ延びる。
従って計算時の補助分母に追加の非消滅条件を課す必要はない。

$\kappa_2\ne0$ のとき、小複体の cohomology は $\mathbb C[T]/(T^2)$ である。
$\kappa_2=0$ のときは $\mathbb C[T]\oplus\mathbb C[T]\eta$ で非有界となる。
長さ三の周囲分解に対する中心商の周期性と Rees の議論
（命題\ref{cmo:prop:character-pd}の証明）から、表示した射影次元を得る。
\end{proof}

\begin{example}[character が検出しない二次元単純加群の無限射影次元]
\label{cmo:ex:two-dimensional-zero-action}
$a$ を $p(T)=T^3-2T^2-11T+8$ の任意の複素根とし、
\[
 \alpha=(a^2+1,a+1,1-a),\qquad
 \lambda=-\frac23(a^3-a)
\]
と置く。
半跡 $(a,1,-1)$ の二次元単純 $D_\alpha$ 加群 $M$ は剛的であるが、
$D_\alpha/(\Omega_\alpha-\lambda)$ 上では射影次元が無限である。
この中心値には $D_\alpha$ の character は一つも存在しない。
また $D_\alpha$ の八個の character は全て自身の中心商上で射影次元二を持つ。
従って character の計算だけでは、中心商における単純加群の有限射影次元を判定できない。
\end{example}

\begin{proof}
$p$ は有理数体上既約であり、
$\gcd(p(T),T(T^2+3)(1-T^4))=1$ である。
従って $a\ne0$、$F(a,1,-1)=a(a^2+3)\ne0$、
$\alpha_1\alpha_2\alpha_3=1-a^4\ne0$ となる。
命題\ref{cmo:prop:two-dimensional-classification}と
命題\ref{cmo:prop:two-dimensional-center-rigidity}から $M$ の存在と剛性を得る。
さらに
\[
 N_6(T,1,-1)
 =(T^3-2T^2-11T+8)(T^3+2T^2-11T-8)
\]
なので、命題\ref{cmo:prop:two-dimensional-universal-kappa}から
$\kappa_2(M)=0$ と無限射影次元を得る。

character に関する二つの主張は、有限エタール代数
\[
 E_a=\mathbb Q(a)[u,v,w]/
 (u^2-\alpha_1,v^2-\alpha_2,w^2-\alpha_3)
\]
における次の二元の可逆性である。
\[
 \begin{split}
 K={}&\alpha_1\alpha_2uv+\alpha_2\alpha_3vw
        +\alpha_3\alpha_1wu+9\alpha_1\alpha_2\alpha_3,\\
 L={}&\alpha_1u+\alpha_2v+\alpha_3w+uvw-(a^3-a).
 \end{split}
\]
ここで $K$ は character の $\kappa_1$ の非零定数倍の分子であり、
$L$ は $-3/2$ 倍した character の中心値と $\lambda$ の差である。
基底 $u^iv^jw^k$（$i,j,k\in\{0,1\}$）で乗法行列を作り、
その行列式を $\mathbb Q[T]$ 上で計算すると、いずれも $p(T)$ と互いに素になる。
従って任意の根 $a$ に対して $K,L$ は可逆である。
再現コード \texttt{verify\_mixed\_two\_dimensional\_ext.py --zero-example} は
これら二つの行列式と互除法を厳密に検査する。
同じコードは $\mathbb Q(a)$ 上で $M$ の全高次 homotopy 恒等式と
$\kappa_2(M)=0$ を、普遍式への代入から独立に再検査する。
\end{proof}

\begin{proposition}[指定二係数の二次元高次作用と百個の単純加群]
\label{cmo:prop:one-two-dimensional-pd}
例\ref{cmo:ex:sixty-four}の二係数に対し、
半跡の一方の符号を次のように選ぶと、高次作用は
\begin{equation}\label{cmo:eq:two-dimensional-exact-kappa}
 \begin{array}{c|c|c}
 \alpha&(a,b,c)&\kappa_2\\\hline
 (1,4,9)&(-5,1,11)/\sqrt{14}&
 \displaystyle\frac{1297144\sqrt{14}}{151263}\\[3pt]
 (16,25,49)&(-323,-53,667)/(3\sqrt{970})&
 \displaystyle\frac{2200813914116392\sqrt{970}}{121008460966875}
 \end{array}
\end{equation}
であり、いずれも非零である。
他方の半跡の符号では $\kappa_2$ の符号も反転する。

各因子の八個の character と二個の二次元単純加群から得られる
百個のテンソル積単純加群は、互いに相異なる百個の中心値
\[
 c\in-(\mathcal R\cup\mathcal T)-100(\mathcal S\cup\mathcal T')
\]
で対応する中心商へ降りる。
これらの各加群は当該中心商上で射影次元五を持ち、
自己 Ext の次元は $0,2,3,5$ 次で一、それ以外で零である。
本命題単独は次元 $1,2,4$ の指定単純加群を扱う。
各因子の全単純加群は定理\ref{cmo:thm:fourteen-node-orders}、
全係数での smoothness は定理\ref{cmo:thm:exact-sixty-four}で決定する。
character の Kuznetsov 射影後の Ext は命題\ref{cmo:prop:projected-point-ext}で扱うが、
二次元因子を含む残りの射影対象の Ext は未決定である。
\end{proposition}

\begin{proof}
式\eqref{cmo:eq:two-dimensional-matrices}を使い、まず半跡を整数にした二つの表現を計算する。
\[
 \begin{array}{c|c|c|c}
 (a,b,c)&r&\text{係数}&\kappa_2\\\hline
 (-5,1,11)&\sqrt{-11}&14(1,4,9)&5188576/3087\\
 (-323,-53,667)&\sqrt{35351}&8730(16,25,49)&
 8803255656465568/19053225
 \end{array}
\]
これらの値は普遍式\eqref{cmo:eq:two-dimensional-universal-kappa}への代入で得られ、
$\mathbb Q(\sqrt{-11})$、$\mathbb Q(\sqrt{35351})$ 上の厳密な有限線型代数でも独立に検査する。
再現コード \texttt{verify\_mixed\_two\_dimensional\_ext.py} は、
次数二以下の $h_j$ と次数三以下の $g$ を式\eqref{cmo:eq:two-dimensional-homotopies}から順に解き、
全恒等式を独立に剰余計算で検査する。
剰余計算は必要な次数以下の全語と全左右関係倍を列挙して行い、
剰余空間の次元も PBW の値と照合する。
同じコードで微分の階数 $3,9,3$ と式\eqref{cmo:eq:two-dimensional-kappa-extraction}を検査する。
これは補間や浮動小数点近似を用いない再生成可能な計算証明書である。
$s=1/\sqrt{14}$ と $s=1/\sqrt{8730}$ の斉次変換から
式\eqref{cmo:eq:two-dimensional-exact-kappa}を得る。
$s=-1$ により反対の符号の類も得られる。

二因子の単純加群を $M_1,M_2$ とする。
その像はそれぞれ全行列環なので、$M_1\otimes M_2$ は単純である。
八個の character に対する $\kappa_1\ne0$ は既に検査されており、
二次元加群については上で $\kappa_2\ne0$ を得た。
各因子の小複体を $\C[T]$ 上でテンソルすると
\[
 \C[T]\otimes\bigwedge(\eta_1,\eta_2),\qquad
 d\eta_1=\kappa(M_1)T^2,\qquad
 d\eta_2=100^2\kappa(M_2)T^2
\]
となる。
片方の奇基底を正規化し、他方からその倍を引けば、微分は
$d\eta_1=T^2,d\eta_2=0$ となる。
従って表示した自己 Ext を全次数で得る。
命題\ref{cmo:prop:character-pd}の証明で用いた中心商の周期性と Rees の議論は、
character であることではなく、長さ六の有限な周囲射影分解を使う。
式\eqref{cmo:eq:two-dimensional-resolution}のテンソル積がこの分解を与えるので、
同じ議論により射影次元は五である。

中心値の数は $8\cdot8+2\cdot8+8\cdot2+2\cdot2=100$ である。
最初の六十四値が相異なることは既証である。
残る値は、それぞれ $\mathbb Q+\mathbb Q^\times\sqrt{14}$、
$\mathbb Q+\mathbb Q^\times\sqrt{970}$、
$\mathbb Q^\times\sqrt{14}+\mathbb Q^\times\sqrt{970}$ に属する。
$1,\sqrt{14},\sqrt{970},\sqrt{14\cdot970}$ の一次独立性と、
各有理部分の相異性により、表示した百個の値は全て相異なる。
\end{proof}

\subsection{中心 chart の character 自己 Ext と qgr の比較}

以下では右加群の規約を使う。先の左 character の計算を移すときは、
有限次元 $M$ に対する自然なベクトル空間同型
\[
 \operatorname{Ext}_{R\text{-left}}^i(M,M)
 \simeq \operatorname{Ext}_{R\text{-right}}^i(M^*,M^*)
\]
を使える。実際、左射影分解 $P_\bullet\to M$ の体上の双対は
右単射分解 $M^*\to P_\bullet^*$ になる。
$P_i$ が平坦であることと tensor--Hom 随伴から $P_i^*$ の単射性が従い、
$\operatorname{Hom}_{R\text{-right}}(M^*,P_i^*)
\simeq\operatorname{Hom}_{R\text{-left}}(P_i,M^{**})$
で比較できる。一文字 character の場合は $M^*$ の座標も同じである。

\begin{lemma}[有限次元 affine 加群の coherent な延長]
\label{cmo:lem:chart-ext}
$A$ を連結局所有限な次数付き右 noetherian 代数とし、
$u\in Z(A)_1$ を正則元とする。
$L=A[u^{-1}]$、$D=L_0$ と置く。
$D$ 上の有限次元右加群 $M$ に対し、
\[
 P(M)=\bigoplus_{n\ge0}M u^n,
 \qquad (m u^n)a_d=(m(a_du^{-d}))u^{n+d}
\]
を次数付き $A$-加群とし、$\mathcal P(M)=\pi P(M)$ と置く。
これは coherent であり、有限次元 $D$-加群 $M,N$ に対し全ての $i\ge0$ で
\[
 \operatorname{Ext}^i_{\operatorname{qgr}A}(\mathcal P(M),\mathcal P(N))
 \simeq\operatorname{Ext}^i_D(M,N)
\]
が成立する。
\end{lemma}
\begin{proof}
$u$ は中心で次数一の可逆元なので $L\simeq D[u,u^{-1}]$ であり、
次数零部分を取る関手は $\operatorname{Gr}L\simeq\operatorname{Mod}D$ を与える。
平坦な中心局所化 $\operatorname{Gr}A\to\operatorname{Gr}L$ は exact であり、
torsion 加群を消すので exact な関手
\[
 j^*:\operatorname{QGr}A\longrightarrow\operatorname{Mod}D
\]
を誘導する。

その右随伴は
$j_*N=\pi\operatorname{Res}(N\otimes_\C\C[u,u^{-1}])$ である。
ここで adjunction が Serre 商に降りることを確認する。
$V\in\operatorname{Gr}L$ の $L$ 上の単射分解を restriction すると、
局所化が exact であるため $\operatorname{Gr}A$ 上の単射分解になる。
従って任意の torsion $T$ に対し
\[
 \operatorname{Ext}^q_{\operatorname{Gr}A}(T,\operatorname{Res}V)
 \simeq\operatorname{Ext}^q_{\operatorname{Gr}L}(T[u^{-1}],V)=0
 \quad(q\ge0).
\]
ゆえに $\operatorname{Res}V$ は torsion-local であり、
商写像の section と同一視できる。これで上の adjunction を得る。
$j_*$ は $\pi$, restriction, $N\mapsto N[u,u^{-1}]$ の合成なので exact であり、
$j^*j_*\simeq\operatorname{id}$ である。
さらに exact な左随伴 $j^*$ を持つので $j_*$ は単射加群を単射対象へ移す。
$N$ の単射分解を $j_*$ で移して計算すれば
\[
 \operatorname{Ext}^i_{\operatorname{QGr}A}(j_*M,j_*N)
 \simeq\operatorname{Ext}^i_D(M,N)
\]
を得る。

$M$ が有限次元なら $P(M)$ は次数零の有限個の元で生成される。
$P(M)\subset M[u,u^{-1}]$ の商は負次数だけにあり、
各斉次元は十分高い正次数部分により消されるので torsion である。
従って $j_*M=\pi P(M)=\mathcal P(M)$ は coherent である。

最後に $\operatorname{QGr}A$ は局所 noetherian Grothendieck 圏で、
その noetherian 対象の部分圏が $\operatorname{qgr}A$ である。
この包含は全ての Ext を保つ。
これは Positselski--Schn\"urer \cite[Proposition~2.1]{CMOSmallBig} の適用である。
同命題の仮定は、noetherian 対象への任意の全射について、その始域の
有限生成部分対象を取ってもなお全射となることから従う。
従って $D^-(\operatorname{qgr}A)\to D^-(\operatorname{QGr}A)$ は完全忠実であり、
上の比較を $\operatorname{qgr}$ に制限できる。
\end{proof}

\subsection{有理係数の非滑らかな混合中心商}

\begin{proposition}[無限自己 Ext を持つ有理係数例]
\label{cmo:prop:rational-singular}
$\alpha=(9,9,1)$ とし、独立な二組の変数を持つ
$\mathcal B_\alpha$ を中心変数 $u$ で tensor する。
\[
 B=\mathcal B_\alpha\otimes_{\C[u]}\mathcal B_\alpha',\qquad
 A=B\big/(\Omega_\alpha+100\Omega_\alpha'-\tfrac{1616}{3}u^3)
\]
と置く。$p=p'=(3,-3,1)$ の定常 character に対応する
$\mathcal P\in\operatorname{qgr}A$ に対し
\[
 \sum_{j\ge0}\dim_\C\operatorname{Ext}^j_{\operatorname{qgr}A}
       (\mathcal P,\mathcal P)\,s^j
 =\frac{(1+s^3)^2}{1-s^2}.
\]
従って $\operatorname{qgr}A$ の homological dimension は無限で、
$D^b(\operatorname{qgr}A)$ は proper でない。
線束生成部分 $\mathcal C_A=\operatorname{thick}\{\mathcal O(j):j\in\mathbb Z\}$ は
proper であるが smooth でなく、線束生成残余部分 $\mathcal K_A^{\rm lin}$ も smooth でない。
また $\mathcal P\notin\mathcal C_A$ なので
$D^b(\operatorname{qgr}A)/\mathcal C_A\ne0$ である。
\end{proposition}
\begin{proof}
本節の普遍 PBW 表示と中心元の証明により、$B$ は左右 noetherian な
Koszul AS-regular 代数で、global dimension と parameter は $7$ である。
中心三次元は正則で、$u$ も $A$ 上正則である。
後者二点は $u=0$ での $\Omega_0+100\Omega_0'$ の正則性からも直接従う。
実際、各 $S=S(1,-1,-1)$ は正則中心元 $\Omega_0$ 上自由であり、
$S\otimes S$ は $\C[\Omega_0,\Omega_0']$ 上自由なので、この中心線型和は正則である。
$B$ の $u$-進 lowest term で中心元の正則性を持ち上げ、
$u x=w y$ を $u$ で還元して $y=u z$ とすれば、$u$ の商上の正則性も得る。

指定 character では
\[
 (ab)^3+(bc)^3+(ca)^3=-729,\qquad
 9a^2b^2c^2=729,\qquad
 \kappa(p)=0,\quad\lambda_p=\frac{16}{3}.
\]
従って二因子の中心和はこの character を消し、自己 Ext の計算複体は
$\C[t]\otimes\bigwedge(\eta_1,\eta_2)$、微分零となる。
前補題によって affine 自己 Ext は全次数で coherent 自己 Ext と一致する。
これが表示した Hilbert 級数と最初の二つの結論を与える。

補題\paperref{main}{crs:lem:sections}と\paperref{main}{crs:lem:finite-lines}の仮定は満たされ、
$\mathcal C_A$ は $\bigoplus_{j=0}^6\mathcal O(j)$ で生成され proper である。
もし smooth なら定理\paperref{main}{crs:thm:coherent}により
$\mathcal C_A=D^b(\operatorname{qgr}A)$ となり、右辺も proper になるが、
$\mathcal P$ の無限自己 Ext と矛盾する。
さらに命題\paperref{main}{crs:prop:residual-smooth}の対偶から残余線束部分も smooth でない。
$\mathcal P$ 自身が proper な $\mathcal C_A$ に入ることも同じ理由で不可能である。
\end{proof}

\paragraph{結論の範囲.}
この例は滑らかな非可換三次多様体の仮定を満たさないため、
滑らかな対象についての可換 Ku 実現予想への反例ではない。
また $D^b(\operatorname{qgr}A)$ の非proper性と、同 dg 圏の smoothness 単独は区別する。
従って本命題は、全 coherent 圏そのものの dg smoothness 単独を否定する主張へ
言い換えない。

% !TEX root = ../../mixed-sklyanin-cubics.tex
\section{普遍中心と Azumaya 幾何}
\label{mix:sec:centers}

\subsection{定数型 Sklyanin 変形の普遍中心と有限中心性}
\label{cmo:subsec:pi-center}

本節では斉次 Sklyanin 入力を $S=S(1,-1,-1)$ に固定する。
その定数項を任意に変えた PBW 変形について、中心の生成元を明示する。
以下で使う $a,b,c$ は一ブロックの三つの変形係数であり、
中心和の切断係数とは区別する。

\begin{lemma}[普遍六次中心元]\label{cmo:lem:universal-sixth-center}
中心変数 $a,b,c$ を用いて
\[
 \mathcal D=\C[a,b,c]\langle x,y,z\rangle/
 \bigl([y,z]-x^2+a,\ [z,x]-y^2+b,\ [x,y]-z^2+c\bigr)
\]
と置く。$a,b,c$ に重み二、$x,y,z$ に重み一を与える。
次の元は $\mathcal D$ の中心に属する：
\begin{align}
 C_0={}&z(z^2-c)(xzy+ax+by)+abz^2,
 \label{cmo:eq:center-six-zero}\\
 C_+={}&yz^4y+zyz^3y+z^2y^2z^2+2z^2yz^2y-z^4y^2+z^5x
 \notag\\
 &+a(zyz^2+2z^2yz+2z^3y)
   +b(-yz^2x-2zy^3+4z^2yx+4z^4)
 \notag\\
 &+c(-y^2z^2-3yz^2y-z^2y^2-2z^3x)
 \notag\\
 &+a^2z^2+b^2(-yz+3zy)+c^2(2y^2+zx)
 \notag\\
 &+2abzx+ac(-yz-2zy)+bc(-yx-5z^2).
 \label{cmo:eq:center-six-positive}
\end{align}
さらに、積の順序を逆転し、$x,a$ を固定し、
$y,z$ および $b,c$ をそれぞれ交換する反自己同型を $\iota$ と置き、
$C_-=\iota(C_+)$ とする。これも中心元である。

$a=b=c=0$ での像を $C_{0,0},C_{+,0},C_{-,0}\in S_6$ と書く。
既述の三次中心元
\[
 \Omega_0=\frac{x^3+y^3+z^3}{3}-\frac23\operatorname{sym}(xyz)
\]
について、四元
\begin{equation}\label{cmo:eq:sixth-center-basis}
 \Omega_0^2,\quad C_{0,0},\quad C_{+,0},\quad C_{-,0}
\end{equation}
は $S_6$ 内で一次独立である。
\end{lemma}

\begin{proof}
三つの関係を $r_x,r_y,r_z$ と書く。
$\iota(r_x)=r_x$、$\iota(r_y)=r_z$、$\iota(r_z)=r_y$ であるから、
$\iota$ は実際に $\mathcal D$ の反自己同型である。
従って $C_-$ の中心性は $C_+$ の中心性からも従う。

中心性と一次独立性には、次の有限有理数計算を用いる。
重み $d$ における自由代数の基底を
\[
 \bigl\{a^{e_1}b^{e_2}c^{e_3}w\ \bigm|\
       2(e_1+e_2+e_3)+|w|=d\bigr\}
\]
とし、全ての二側関係倍
$a^{e_1}b^{e_2}c^{e_3}v r_j w$ を行とする。
行簡約は $\mathbb Q$ 上で行うので、零への簡約は二側関係イデアルへの
厳密な所属証明である。重み六と七での行列の大きさ、rank はそれぞれ
\[
\begin{array}{c|r|r|r|r}
 d&\text{行数}&\text{列数}&\text{rank}&\text{商の次元}\\ \hline
 6&1476&1036&917&119\\
 7&5454&3108&2919&189
\end{array}
\]
となる。$[x,C_\nu],[y,C_\nu],[z,C_\nu]$ は
$\nu\in\{0,+,-\}$ の全てについて重み七の関係行列で零へ簡約される。
同じ計算は
$\Omega=\Omega_0-ax-by-cz$ の中心性も重み四で確認する。

active word の長さを降順、同じ長さでは $x<y<z$ の辞書式順序とし、
同じ word に付くパラメータ指数を辞書式に並べる。
重み六の簡約後に $a=b=c=0$ とすると、$S_6$ の二十八個の自由列を得る。
式\eqref{cmo:eq:sixth-center-basis}を行に取り、そのうち
\[
 y^2z^3y,\quad yz^2y^3,\quad yz^4y,\quad z^3y^3
\]
の四列を選んだ係数行列は
\begin{equation}\label{cmo:eq:sixth-center-minor}
 \begin{pmatrix}
  4/9&0&0&4/3\\
  0&0&0&-1\\
  0&0&1&0\\
  0&-1&0&0
 \end{pmatrix},\qquad \det=-\frac49\ne0
\end{equation}
である。よって一次独立性が従う。
この計算を再生成する有限証明書が
\path{computations/verify_mixed_pi_center.py} である。
係数の代入・補間・浮動小数点 rank による検査ではなく、
三つのパラメータを独立変数に保つ普遍恒等式を確認している。
\end{proof}

\begin{theorem}[全定数変形での有限中心性]\label{cmo:thm:finite-center}
任意の $\alpha=(\alpha_1,\alpha_2,\alpha_3)\in\C^3$ に対して
\[
 D_\alpha=\C\langle x,y,z\rangle/
 \bigl([y,z]-x^2+\alpha_1,\ [z,x]-y^2+\alpha_2,
       \ [x,y]-z^2+\alpha_3\bigr)
\]
と置く。$C_{\nu,\alpha}$ を上の $C_\nu$ の係数特殊化とすると
\begin{equation}\label{cmo:eq:affine-center-generators}
 Z(D_\alpha)=\C[\Omega_\alpha,C_{0,\alpha},C_{+,\alpha},C_{-,\alpha}].
\end{equation}
ここで affine 表示では
$\Omega_\alpha=\Omega_0-\alpha_1x-\alpha_2y-\alpha_3z$ である。
$D_\alpha$ はこの可換有限型中心環上有限であり、従って PI である。

さらに、$u$ を次数一の中心変数とする斉次化 $\mathcal B_\alpha$ で
\[
 \widehat C_{\nu,\alpha}
     =C_\nu(x,y,z;\alpha_1u^2,\alpha_2u^2,\alpha_3u^2)
\]
と定めると
\begin{equation}\label{cmo:eq:homogeneous-center-generators}
 Z(\mathcal B_\alpha)=\C[u,\Omega_\alpha,
             \widehat C_{0,\alpha},\widehat C_{+,\alpha},\widehat C_{-,\alpha}].
\end{equation}
右辺の $\Omega_\alpha$ は $\mathcal B_\alpha$ の斉次三次中心元を表す。
$\mathcal B_\alpha$ もこの中心上有限である。
\end{theorem}

\begin{proof}
まず、全次数に必要な既知の中心生成定理を特定する。
$S=S(1,-1,-1)$ の点曲線は
$E=V(x^3+y^3+z^3+xyz)$ であり、その点 shift は位数六である。
後者は Walton \cite[Proposition~4.1]{CMOWaltonRepresentations} にも
同じパラメータで記載されている。
Smith--Tate \cite[Theorem~3.7, $r=1$]{GRSklyaninST} によれば、
位数 $n$ の三次元 Sklyanin 代数の中心は
次数三の元 $g$ と次数 $n$ の三元 $z_1,z_2,z_3$ で生成される。
この定理は $3\mid n$ の場合も含む。
本件の $n=6$ に適用すると
\[
 Z(S)=\C[g,z_1,z_2,z_3],\qquad
 \deg g=3,\quad\deg z_i=6.
\]
従って $\dim_\C Z(S)_3=1$ かつ $\dim_\C Z(S)_6\le4$ である。
$\Omega_0$ は非零である：例えば可換化
$S_{\mathrm{ab}}=\C[x,y,z]/(x^2,y^2,z^2)$ への像は
$-2xyz/3\ne0$ である。よって $\Omega_0$ は $g$ の非零 scalar 倍である。
補題\ref{cmo:lem:universal-sixth-center}は $Z(S)_6$ の四つの独立元を与えるので、
\emph{$\dim_\C Z(S)_6=4$ であり、式\eqref{cmo:eq:sixth-center-basis}はその全体の基底}
となる。特に各 $z_i$ はこの四元の線形結合であり、
\begin{equation}\label{cmo:eq:graded-center-generators}
 Z(S)=\C[\Omega_0,C_{0,0},C_{+,0},C_{-,0}].
\end{equation}
これは低次数の中心計算だけから全次数の生成性を推測した結論ではない。
全次数の生成上界が Smith--Tate の定理であり、有限計算はその全生成元の
持上げを確認している。

既述の PBW 判定により $D_\alpha$ の active 濾過の随伴代数は $S$ である。
任意の中心元の最高部分は $Z(S)$ に属する。
式\eqref{cmo:eq:graded-center-generators}の全生成元は、補題で与えた
$D_\alpha$ の中心元へ持ち上がっている。
従って任意の $h\in Z(D_\alpha)$ の最高部分と一致する中心生成元の多項式を
$h$ から引き、active 次数について帰納すれば
式\eqref{cmo:eq:affine-center-generators}を得る。

同じ Smith--Tate の定理により $S$ は $Z(S)$ 上有限である。
有限個の斉次加群生成元を取り、これらを $D_\alpha$ に持ち上げる。
任意の元の最高部分をそれらで表し、次数を下げる帰納法を繰り返すと、
同じ有限個の持上げが $D_\alpha$ を $Z(D_\alpha)$ 上生成する。

斉次化については、$C_\nu$ の普遍恒等式に
$(a,b,c)=(\alpha_1u^2,\alpha_2u^2,\alpha_3u^2)$ を代入すれば
$\widehat C_{\nu,\alpha}$ の中心性を得る。
$u$ は $\mathcal B_\alpha$ 上正則であり、$\mathcal B_\alpha/(u)=S$ である。
これらの中心元の像が $Z(S)$ を全て生成するため、
Smith--Tate \cite[Lemma~3.6]{GRSklyaninST} が適用でき、
式\eqref{cmo:eq:homogeneous-center-generators}と加群有限性を得る。
実際、この段階の帰納法は、中心元 $h$ の像を生成元の多項式で消して
$h-P=u h_1$ と書き、$u$ の正則性から $h_1$ の中心性を得るものである。
\end{proof}

\begin{corollary}[係数族に一様な有限中心模型]
\label{cmo:cor:relative-finite-center}
上の中心元は三つの定数係数に関して多項式であり、係数の分母は生じない。
したがって、係数環 $R$ 上で $\alpha_i\in R$ を指定し、既述の分裂した
PBW 模型を底変更して得る $\mathcal B_{\alpha,R}$ は、
\[
 R[u,\Omega_\alpha,
       \widehat C_{0,\alpha},\widehat C_{+,\alpha},\widehat C_{-,\alpha}]
\]
という可換有限型中心部分環上有限である。
特に $R=\C[t]$、$\alpha\mapsto t^2\alpha$ の収縮族に適用できる。
複数因子が同じ $u$ を共有する tensor 積も、各因子のこの有限中心部分環の
$R[u]$ 上の tensor 積の像に対して有限である。
\end{corollary}

\begin{proof}
$S$ の $Z(S)$ 上の有限斉次加群生成元を一度選び、その同じ語を全係数族で使う。
PBW によって $\mathcal B_{\alpha,R}/(u)=S\otimes_\C R$ であり、
式\eqref{cmo:eq:graded-center-generators}の係数拡大は、表示した中心部分環の像に
含まれる。任意の斉次元を $u=0$ でこの有限個の語の結合として表し、差を $u$ で割る。
内部次数が毎回一ずつ下がるため、有限回で停止する。
この証明は $R$ の点ごとの有限生成性からの推論ではなく、一組の有限生成語を
全係数族で共通に用いる。
因子ごとの有限生成系の tensor 積を取れば最後の主張が従う。
\end{proof}

\begin{remark}[先行研究との関係と新規性の範囲]
\label{cmo:rem:pi-center-literature}
Walton の学位論文 \cite[Definition~V.1, Remark~V.16]{CMOWaltonThesis} は、
三次元 Sklyanin 代数の一般中心拡大を扱い、位数六の場合について
三次・六次中心元を持ち上げる方法を提案している。
本節はそのうち、斉次パラメータが $(1,-1,-1)$、線形補正 $d_i x_i u$ が零で、
定数補正 $e_i u^2$ が任意の部分族について、その持上げを明示し検証したものである。
位数六の全 Sklyanin パラメータや、非零の線形補正を含む一般拡大を扱う主張ではない。

Etingof--Ginzburg \cite[Theorem~3.3.2 and Section~9]{NFMEG} の
三次中心元と PBW 変形の理論、および Smith--Tate の斉次中心生成定理は既知の入力である。
本節の具体的六次持上げとこの定数型部分族への適用について、
後続文献を含む初出の網羅的確認はまだ完了していない。
有限中心性は明示した中心元から証明しており、associated graded の PI 性だけから
PBW 変形の PI 性を結論していない。
\end{remark}

\subsection{中心の有限表示と二つの係数における臨界値}
\label{cmo:subsec:effective-center-geometry}

定理~\ref{cmo:thm:finite-center} の中心生成元を用い、
\[
 W=\Omega_\alpha^2,\qquad U=C_0,\qquad V=C_+,\qquad T=C_-
\]
とおく。これらのフィルター次数はすべて $6$ である。
以下の三次多項式を固定する：
\begin{equation}
\begin{split}
 P_3(W,U,V,T)={}&81W^2U+468WU^2+784U^3\\
 &-36WVT-176UVT+16V^3+16T^3.
\end{split}
\label{cmo:eq:center-cubic-top}
\end{equation}

\begin{proposition}[中心の一関係表示と PI 次数]
\label{cmo:prop:center-hypersurface-degree}
任意の $\alpha\in\mathbb C^3$ に対し、最高次部分が $P_3$ である
全次数 $3$ の多項式 $P_\alpha(W,U,V,T)$ が一意に存在し、
\begin{equation}
 Z(D_\alpha)\simeq
 \mathbb C[s,U,V,T]/\bigl(P_\alpha(s^2,U,V,T)\bigr),
 \qquad s\longmapsto\Omega_\alpha
 \label{cmo:eq:center-hypersurface}
\end{equation}
となる。さらに $D_\alpha$ の中心上の分数体階数は $36$、
PI 次数は $6$ である。
\end{proposition}

\begin{proof}
定理~\ref{cmo:thm:finite-center} の証明は
$\operatorname{gr}Z(D_\alpha)=Z(S)$ も与える。
Smith--Tate \cite[Theorem~4.8]{GRSklyaninST} によれば、
$S=S(1,-1,-1)$ の中心は次数 $3,6,6,6$ の生成元を持ち、
関係イデアルは次数 $18$ の一元で生成される。
補題~\ref{cmo:lem:universal-sixth-center} の四元の最高部は
この生成系の可逆な線型変更を与える。
後述する厳密計算から $P_3(\Omega_0^2,C_{0,0},C_{+,0},C_{-,0})=0$
が従うので、この非零多項式が最高次数での関係を生成する。

中心への全射はフィルター付きで狭義であるから、
この次数 $18$ の関係を低次数で補正して核の元に持ち上げられる。
$x,y,z\mapsto-x,-y,-z$ は $D_\alpha$ の自己同型であり、
$\Omega_\alpha\mapsto-\Omega_\alpha$、$C_\nu\mapsto C_\nu$ と作用する。
補正の偶数部分を取れば、関係は $s$ の偶数冪だけからなる。
従って $W=s^2$ を用いて全次数 $3$ 以下の $P_\alpha$ と書ける。
これより低いフィルター次数には関係がないため、最高部分を
$P_3$ と規格化した持上げは一意である。
任意の核の元の最高部分は $P_3$ で割り切れるので、
$P_\alpha$ の倍を引いてフィルター次数を下げる帰納法により、
核全体が $P_\alpha$ で生成される。

中心の Hilbert 級数は
\[
 H_{\operatorname{gr}Z(D_\alpha)}(t)=
 \frac{1-t^{18}}{(1-t^3)(1-t^6)^3},
 \qquad H_{\operatorname{gr}D_\alpha}(t)=\frac1{(1-t)^3}.
\]
有限中心性により $D_\alpha$ は中心上の有限フィルター加群であり、
両者の多重度の比は分数体階数に等しい。
従ってその階数は
\[
 \lim_{t\to1}
 \frac{(1-t^3)(1-t^6)^3}{(1-t)^3(1-t^{18})}
 =\frac{3\cdot6^3}{18}=36.
\]
$D_\alpha$ は PBW 性と $\operatorname{gr}D_\alpha=S$ の整域性から
整域である。中心の分数体に移した中心単純環の次元は PI 次数の
二乗なので、PI 次数は $6$ となる。
\end{proof}

以下では
\[
 \alpha^{(1)}=(1,4,9),\qquad
 \alpha^{(2)}=(16,25,49)
\]
とする。$P_i=P_{\alpha^{(i)}}$ は $P_3$ に次表の低次項を加えた
多項式である。表にない単項式の係数は零である。
\begin{center}
\small
\begin{tabular}{c|r|r}
単項式 & $P_1$ の係数 & $P_2$ の係数\\ \hline
$WU$ & $-59760$ & $-11301840$\\
$WV$ & $144$ & $230400$\\
$WT$ & $324$ & $451584$\\
$U^2$ & $-193984$ & $-39551680$\\
$UV$ & $-13952$ & $-3701824$\\
$UT$ & $-20752$ & $-3076672$\\
$V^2$ & $5152$ & $718992$\\
$VT$ & $13856$ & $2825120$\\
$T^2$ & $3472$ & $899856$\\
$W$ & $-47952$ & $-16719897600$\\
$U$ & $7099264$ & $180425753600$\\
$V$ & $-152048$ & $-43671732224$\\
$T$ & $-153248$ & $-48922542080$\\
$1$ & $997056$ & $448348972646400$
\end{tabular}
\end{center}

\begin{lemma}[有限計算による関係の証明]
\label{cmo:lem:center-relation-certificate}
上の二つの表は、それぞれ
$P_i(\Omega_{\alpha^{(i)}}^2,C_0,C_+,C_-)=0$ を与える。
全次数 $3$ 以下の四変数単項式 $35$ 個の間の関係空間は一次元である。
\end{lemma}

\begin{proof}
\path{computations/verify_mixed_center_geometry.py} は
既存の \path{verify_mixed_pi_center.py} の中心元を直接読み込み、
すべての演算を $\mathbb Q$ 上で行う。再現の線型代数を説明する。
フィルター次数 $n-1$ 以下の語代表基底と右乗法表が既に構成されたとする。
次数 $n$ の仮の最高部分には、次数 $n-1$ の基底語の右に
$x,y,z$ を付加した $3\binom{n+1}{2}$ 列だけを用いる。
低次数部分には既知の基底を用いる。
次数 $n-2$ の基底語 $w$ に対し
\[
 w(yz-zy-x^2+\alpha_1),\quad
 w(zx-xz-y^2+\alpha_2),\quad
 w(xy-yx-z^2+\alpha_3)
\]
を既知の乗法表で表し、その行空間で割る。
これらは実際の関係なので、得られた空間から $F_nD_\alpha$ への
全射がある。計算は各 $2\le n\le18$ で新しい最高部分の次元が
$\binom{n+2}{2}$ であり、低次数に新しい関係がないことを検査する。
既知の PBW 次元と一致するため、この全射は同型である。
したがって得られた乗法表は次数 $18$ までの正確な商計算を与える。

$n=18$ では最高部分の仮の列数は $513$、実際の最高部分の次元は
$190$、低次数を含む語代表は $1330$ 個である。
四中心元の全次数 $3$ 以下の単項式をこの基底に展開すると
$35$ 列の階数は $34$ となり、核は表示した $21$ 項の関係で張られる。
これは自由語 $3^{18}$ 個の列挙、浮動小数点の階数判定、係数の標本補間の
いずれにも依存しない有限の理想所属計算である。
全次数における中心表示には、命題の証明で Smith--Tate の
一関係定理を別に用いている。
\end{proof}

臨界値を表示するため、次の多項式を置く：
\begin{align*}
 H_1(W)={}&(W-784)(W-256)(9W-3136)(9W-784)(343W-73728),\\
 H_2(W)={}&(W-200704)(W-38416)(W-30976)(9W-80656)\\
 &\hspace{18mm}\cdot(748505944125W-37623589339922432),
\end{align*}
\begin{align*}
 Q_1(W)={}&129140163W^4-188108856360W^3+92316693932208W^2\\
 &-16917496764590592W+875501309973348352,\\
 Q_2(W)={}&129140163W^4-34079674491720W^3
       +2934863745734897712W^2\\
 &-83555308663913970224640W
       +165916680996850221150748672.
\end{align*}

\begin{proposition}[中心の特異点と八つの新しい Morse 臨界値]
\label{cmo:prop:center-critical-polynomials}
$X_i=\operatorname{Spec}Z(D_{\alpha^{(i)}})$ とし、
$s=\Omega_{\alpha^{(i)}}:X_i\to\mathbb A^1$ とする。
\begin{enumerate}
 \item $s^{-1}(\lambda)$ が特異であるための必要十分条件は
 \[
   \lambda H_i(\lambda^2)Q_i(\lambda^2)=0
 \]
 である。この条件は相異なる $19$ 個の値を与える。
 \item $X_i$ 自身の特異点は $14$ 個で、すべて通常二重点である。
 そのうち $10$ 個は $H_i(s^2)=0$ 上にあり、
 残りの $4$ 個は $s=0$ 上にある。
 \item $Q_i(s^2)=0$ に対応する $8$ 個の臨界点は $X_i$ の滑らかな部分にあり、
 それぞれ $s$ の Morse 臨界点である。
\end{enumerate}
\end{proposition}

\begin{proof}
$R_i=\mathbb Q[W,U,V,T]$ 内のイデアル
\[
 I_i=(P_i,\partial_UP_i,\partial_VP_i,\partial_TP_i)
\]
を用いる。同じ再現コードの有理数体上の Gr\"obner 基底計算は
\begin{equation}
 \dim_{\mathbb Q}R_i/I_i=18,\qquad
 I_i\cap\mathbb Q[W]=(W H_i(W)^2Q_i(W))
 \label{cmo:eq:center-critical-elimination}
\end{equation}
を、非零有理数倍を除いて検証する。
$Q_i$ は $\mathbb Q$ 上既約であり、$W,H_i,Q_i$ は互いに素、
$H_i$ は重根を持たない。
臨界スキームは零次元なので射影の像も閉であり、
消去多項式の根はちょうど特異なファイバーを与える。
これに $W=s^2$ を代入すると (1) を得る。

$P_{i,W}=\partial_WP_i$ と略記する。
$I_i+(P_{i,W})$ の商は長さ $5$ で、
$W$ の消去多項式は $H_i$、しかも四変数 Hessian の行列式は
この商で単元である。従って非零の $s$ 上の中心特異点は
$H_i(s^2)=0$ で与えられる $10$ 個の通常二重点である。
$W=0$ での $I_i$ の商は長さ $4$ の被約環であり、
$P_{i,W}$ と三変数 $U,V,T$ に関する Hessian の行列式はともに単元である。
$s=0$ における $P_i(s^2,U,V,T)$ の Hessian は、
$s$ 方向の成分 $2P_{i,W}$ とこの三変数 Hessian の直和なので可逆である。
これで (2) が従う。

最後に $I_i+(Q_i)$ の商は長さ $4$ である。
この商でも $P_{i,W}$ と三変数 Hessian の行列式は単元である。
$W\ne0$ なので $\partial_sP_i(s^2,U,V,T)=2sP_{i,W}\ne0$ となり、
対応する中心の点は滑らかである。
陰関数定理で $s$ を $U,V,T$ の関数とみれば、その Hessian は
非零スカラー倍の三変数 Hessian となる。従って八つの点は Morse 臨界点である。
各単元判定は、対応する Gr\"obner 基底に当該行列式または微分を加えた
イデアルが単位イデアルになることを有理数体上で検証している。
\end{proof}


\subsection{全定数変形の Azumaya locus}

中心の臨界点から非可換中心商の特異性を結論するには、
中心上の代数の局所構造を確認しなければならない。
本表示では、その入力を全定数係数について得られる。

\begin{theorem}[正則中心上の Azumaya 性]
\label{cmo:thm:azumaya-regular}
任意の $\alpha\in\C^3$ に対して、$D_\alpha$ は PI 次数六の
homologically homogeneous 代数である。
その中心 $Z_\alpha=Z(D_\alpha)$ は正規で、$D_\alpha$ は
$Z_\alpha$ 上 Cohen--Macaulay である。
さらに
\[
 \operatorname{Az}(D_\alpha)=
       (\operatorname{Spec} Z_\alpha)_{\mathrm{reg}}.
\]
特に中心の任意の正則閉点で、完備化した代数は
$\operatorname{Mat}_6(\widehat{Z_{\alpha,\mathfrak m}})$ と同型である。
\end{theorem}
\begin{proof}
三次 potential の長さ三の自己双対な双加群解消により
$D_\alpha$ は三次元 Calabi--Yau で大域次元は三以下である
\cite[Theorem~3.3.2]{NFMEG}。
随伴代数 $S=S(1,-1,-1)$ は domain なので、$D_\alpha$ も domain である。
定理\ref{cmo:thm:finite-center}より中心上有限であり、
全ての単純加群 $V$ は有限次元である。
Calabi--Yau 双対から
\[
 \operatorname{Ext}_{D_\alpha}^3(V,V)
       \simeq\operatorname{End}_{D_\alpha}(V)^*\ne0
\]
なので、全ての単純加群の射影次元は正確に三である。
従って $D_\alpha$ は homologically homogeneous である。
Brown--Hajarnavis による中心の正規性と中心上の Cohen--Macaulay 性を適用する。
ここではこの既知の結果の形を、
Brown--Yakimov \cite[\S3.8(b)(i), Proposition~3.10 and its proof]{CMOBrownYakimov}
に記載されたものに従う。
PI 次数六は、命題\ref{cmo:prop:center-hypersurface-degree}の Hilbert 級数の比から従う。

余次元一での Azumaya 性を、斉次化から示す。
$B=\mathcal B_\alpha$、$T=Z(B)$ とすると、全中心生成元の持上げと
$u$ の正則性により
\[
 B/(u)=S,\qquad T/(u)=Z(S),\qquad
 B[u^{-1}]=D_\alpha[u,u^{-1}],\qquad
 T[u^{-1}]=Z_\alpha[u,u^{-1}].
\]
$T/(u)$ への中心の全射性は定理\ref{cmo:thm:finite-center}で証明済みであり、
その核が $uT$ であることは $ub\in Z(B)\Rightarrow b\in Z(B)$ による。
$T$ は次元四の連結正次数付き可換有限型 domain で、$B$ はその上有限である。

$u=0$ 上で $S$ が Azumaya である閉点 $z$ を取る。
局所環 $T_z$ 上で $B_z$ と $T_z$ に $u$ は正則に作用し、
$B_z/uB_z$ は $T_z/(u)$ 上有限射影である。
この有限射影加群は局所環 $T_z/(u)$ 上で階数三十六の自由加群である。
その基底を持ち上げると、Nakayama により全射 $T_z^{36}\to B_z$ を得る。
核を $L$ とすると、$u$ の $B_z$-正則性による
$\operatorname{Tor}_1^{T_z}(B_z,T_z/(u))=0$ から $L/uL=0$ が従う。
$L$ は有限生成なので再び Nakayama により $L=0$。
従って $B_z$ は $T_z$ 上有限自由である。
自然な射
\[
 B_z\otimes_{T_z}B_z^{\mathrm{op}}
       \longrightarrow\operatorname{End}_{T_z}(B_z)
\]
は、$u$ で割れば同型である。両辺は同じ階数の有限自由加群なので、
Nakayama の補題、または行列式の単元性からこの射自身が同型となる。
従って $B_z$ は Azumaya である。

$\mathcal N\subset\operatorname{Spec}T$ を $B$ の非 Azumaya locus とする。
これは閉集合で、正次数の dilation 作用で不変である。
Walton--Wang--Yakimov
\cite[Theorem~1.3(1), Lemma~5.4]{CMOWaltonWangYakimov}によれば、
位数六の $S$ の非 Azumaya locus は、中心の特異 locus に等しく、
三本の曲線の和である。従って、前段の持上げから
\[
 \dim\bigl(\mathcal N\cap V(u)\bigr)\leq1.
\]
もし $\mathcal N$ に次元三以上の既約成分 $Y$ があれば、
$\mathbb G_m$ の連結性により $Y$ 自身が dilation 不変であり、
斉次閉錐として原点を含む。
$u|_Y=0$ なら上の次元上限に反する。
$u|_Y\ne0$ なら、非空の主切断 $Y\cap V(u)$ は
有限型 domain の次元公式から次元 $\dim Y-1\geq2$ であり、やはり矛盾する。
従って $\dim\mathcal N\leq2$ である。

$u\ne0$ 上の上記 Laurent 表示では、Azumaya 性の忠実平坦降下により
\[
 \mathcal N\cap D(u)=
       \bigl(\operatorname{Spec}Z_\alpha\setminus
                     \operatorname{Az}(D_\alpha)\bigr)\times\mathbb G_m.
\]
従って $D_\alpha$ の非 Azumaya locus は次元一以下であり、
$D_\alpha$ は中心の余次元一で Azumaya である。

以上で、prime、中心上有限、有限大域次元、正規中心、中心上 Cohen--Macaulay、
余次元一の Azumaya 性が揃った。
reduced trace を選べば Brown--Yakimov の Main Theorem の trace 仮定も満たされ、
同論文\cite[Corollary~5.2]{CMOBrownYakimov}から
Azumaya locus と正則 locus の一致が従う。
完備局所環は剰余体 $\C$ を持つ Hensel 環なので、その上の Azumaya 代数は分裂する。
PI 次数が六であることから、最後の行列代数の大きさも従う。
\end{proof}

\begin{remark}
この持上げは、中心の全生成元が持ち上がることと、
斉次化の正次数構造を使っている。
PBW 性、有限中心性、有限大域次元の三条件だけから
Azumaya locus と正則 locus の一致を推論してはいない。
\end{remark}

% !TEX root = ../../mixed-sklyanin-cubics.tex
\section{中心節点の局所代数と中央切断}
\label{mix:sec:local-orders}

\subsection{十四節点における完備局所代数}

三次元節点上の完備局所代数を決定するため、まず一般的な補題を示す。
ここで NCCR とは、正規 Gorenstein 環 $R$ 上の非零反射的加群 $M$ に対し、
$\operatorname{End}_R(M)$ が Cohen--Macaulay な $R$-加群であり、
全ての素点で局所大域次元が局所中心の次元に等しいことをいう。
この定義と有限大域次元を用いる定義の一致は
\cite[Definition~1.7 and Lemma~2.23]{CMOIyamaWemyssMM} に従う。

\begin{lemma}[三次元完備節点の穿孔スペクトラムの Brauer 群]
\label{cmo:lem:formal-node-brauer}
\[
 R=\mathbb C[[x,y,z,w]]/(xy-zw),\qquad
 X=\operatorname{Spec}R,\qquad
 U=X\setminus\{\mathfrak m\}
\]
とする。このとき
\[
 \operatorname{Br}'(U):=
 H^2_{\text{ét}}(U,\mathbb G_m)_{\mathrm{tors}}=0.
\]
従って $U$ 上の任意の Azumaya 代数は分裂する。
特に次数 $n$ の Azumaya 代数 $\mathcal B$ は、階数 $n$ の
ベクトル束 $\mathcal E$ を用いて
$\mathcal B\simeq\mathcal E nd_{\mathcal O_U}(\mathcal E)$ と書ける。
\end{lemma}

\begin{proof}
$X\times\mathbb P^1$ の中の閉部分スキーム
\[
 Y=\{xt=zs,\ wt=ys\},
 \qquad [s:t]\in\mathbb P^1,
\]
と第一射影 $f:Y\to X$ を取る。
この射は projective であり、閉ファイバーは
$E=f^{-1}(\mathfrak m)=\mathbb P^1_{\mathbb C}$ である。
$U$ 上では行列
$\left(\begin{smallmatrix}x&z\\w&y\end{smallmatrix}\right)$ の
階数は一なので $[s:t]$ が一意に定まり、
$Y\setminus E\xrightarrow{\sim}U$ となる。

純性の仮定を、完備化した基底に対して直接確認する。
$E$ 上の閉点のうち $s\ne0$ の chart で $t/s=a$ となる点を取り、
$h=t/s-a$ と置く。その点での完備局所環は
\[
 \begin{split}
 \widehat{\mathcal O}_{Y,p}
 &\simeq
 \mathbb C[[x,y,z,w,h]]/
 (xy-zw,\ x(a+h)-z,\ w(a+h)-y)\\
 &\simeq \mathbb C[[x,w,h]].
 \end{split}
\]
この表示で $E$ のイデアルは $(x,w)$ である。
$t\ne0$ の chart も同様なので、$E$ 上の閉点で $Y$ は正則、
かつ $E$ の余次元は二である。
$E$ の一般点の局所環も、これらの正則局所環の局所化なので正則で、
その次元は二である。
従って $E$ の全ての点で $\mathcal O_{Y,p}$ は次元二以上の正則局所環である。
$Y$ は Noetherian であるから、
\v{C}esnavi\v{c}ius の純性定理
\cite[Theorem~1.1 (=Theorem~6.1)]{CMOCesnaviciusBrauer}
の全仮定を満たし、制限により
\begin{equation}
 H^2_{\text{ét}}(Y,\mathbb G_m)
 \xrightarrow{\sim}
 H^2_{\text{ét}}(Y\setminus E,\mathbb G_m)
 =
 H^2_{\text{ét}}(U,\mathbb G_m)
 \label{cmo:eq:node-brauer-purity}
\end{equation}
を得る。

任意の $n\ge1$ を取る。$R$ は剰余体 $\mathbb C$ を持つ完備局所環なので
strictly Henselian である。
strictly Henselian 局所環のスペクトラム上では大域切断関手は
閉点の stalk 関手と一致し exact である
（\cite[Lemma~59.55.1, Tag~03QN]{CMOStacksNode}）。
従って、proper base change
（\cite[Theorem~59.91.11, Tag~095T]{CMOStacksNode}）と Leray spectral sequence により、
自然な制限写像は同型
\[
 H^2_{\text{ét}}(Y,\mu_n)
 \xrightarrow{\sim}
 H^2_{\text{ét}}(E,\mu_n)
 \simeq\mathbb Z/n\mathbb Z
\]
となる。最後の同型は $E=\mathbb P^1_{\mathbb C}$ に対する
曲線の Kummer 計算
（\cite[Lemma~59.69.1, Tag~03RQ]{CMOStacksNode}）である。
ここでは $f$ の smoothness を仮定しておらず、
proper base change を使っている。

第二射影からの引戻しで得られる $\mathcal L=\mathcal O_Y(1)$ は
$E$ に制限すると $\mathcal O_{\mathbb P^1}(1)$ となる。
その Kummer Chern class は
$H^2_{\text{ét}}(E,\mu_n)$ の生成元なので、
上の同型と自然性により
\[
 \operatorname{Pic}(Y)/n
 \longrightarrow H^2_{\text{ét}}(Y,\mu_n)
\]
は全射である。
$n$ は $\mathbb C$ 上で可逆だから、Kummer 完全列
（\cite[Lemma~59.28.1, Tag~03PK]{NFMStacksKummer}）は
\[
 \operatorname{Pic}(Y)/n
 \longrightarrow H^2_{\text{ét}}(Y,\mu_n)
 \longrightarrow H^2_{\text{ét}}(Y,\mathbb G_m)[n]
 \longrightarrow0
\]
を与える。従って最後の群は零である。
$n$ は任意なので $\operatorname{Br}'(Y)=0$、
式 \eqref{cmo:eq:node-brauer-purity} から
$\operatorname{Br}'(U)=0$ を得る。

Azumaya Brauer 群は cohomological Brauer 群に単射で入り、
次数 $n$ の Azumaya 代数の類は $n$ で消える
（\cite[Section~59.62, Lemma~59.62.2, Tag~0A2J]{CMOStacksNode}）。
従って $\mathcal B$ の Brauer 類は零である。
中心拡大
$1\to\mathbb G_m\to\operatorname{GL}_n\to
\operatorname{PGL}_n\to1$
の非可換 cohomology 完全列で、その
$\operatorname{PGL}_n$-torsor は $\operatorname{GL}_n$-torsor に持ち上がる。
対応する階数 $n$ のベクトル束 $\mathcal E$ が主張を与える。
\end{proof}

\begin{lemma}[完備化した有限代数の大域次元]
\label{cmo:lem:completion-global-dimension}
$B$ を両側 Noetherian 代数、$Z\subset Z(B)$ を可換 Noetherian 部分環とし、
$B$ は $Z$ 上有限とする。
$\mathfrak m\subset Z$ を剰余体 $\mathbb C$ の極大イデアルとし、
\[
 R=\widehat{Z_{\mathfrak m}},\qquad
 A=B\otimes_Z R
\]
と置く。$B$ の左および右大域次元が $d$ 以下なら、$A$ の左右の大域次元も
$d$ 以下である。
\end{lemma}
\begin{proof}
$A$ は完備局所環 $R$ 上有限なので半完全環であり、
有限生成加群は有限生成射影加群による極小射影解消を持つ。
これには、$\mathfrak m R A\subset\operatorname{rad}A$、
$A/\mathfrak m R A$ が有限次元代数であること、および
完備性による冪等元の持上げを用いる。
任意の左単純 $A$-加群 $S$ は $\mathfrak m R$ で消されるので、
$B/\mathfrak m B$ の単純加群から来る。
$B$ 上の長さ $d$ 以下の有限生成射影解消を取り、
$Z\to R$ の平坦性で底変更すると、$S$ の $A$ 上の射影解消を得る。
従って $\operatorname{pd}_A S\le d$ である。

有限生成右 $A$-加群 $N$ の極小射影解消 $P_\bullet$ を考える。
もし $P_{d+1}\ne0$ なら、ある左単純加群 $S$ に対し
$P_{d+1}\otimes_A S\ne0$ となる。
極小性により $P_\bullet\otimes_A S$ の微分は全て零なので、
$\operatorname{Tor}_{d+1}^A(N,S)\ne0$ である。
これは先ほどの左射影次元の評価に反する。
従って全ての有限生成右加群の射影次元は $d$ 以下である。
Noetherian 環の大域次元は有限生成加群の射影次元の上限に等しいので、
右大域次元は $d$ 以下となる。左右を交換して同じ結論を得る。
\end{proof}

\begin{proposition}[節点上の分裂 NCCR の局所形]
\label{cmo:prop:conifold-order}
\[
 R=\mathbb C[[x,y,z,w]]/(xy-zw),\quad
 U=\operatorname{Spec}R\setminus\{\mathfrak m\},\quad I=(x,z)
\]
とする。中心が $R$ である有限 $R$-代数 $A$ に対し、次を仮定する。
\begin{enumerate}
 \item $A$ は $R$ 上 Cohen--Macaulay である。
 \item $A|_U$ は次数 $n$ の Azumaya 代数である。
 \item $A$ の大域次元は有限である。
\end{enumerate}
このとき正整数 $a,b$ が存在して
\begin{equation}
 A\simeq \operatorname{End}_R(R^{\oplus a}\oplus I^{\oplus b}),\qquad a+b=n
 \label{cmo:eq:conifold-order}
\end{equation}
となる。これは $R$-代数としての同型であり、従って特に $A$ は $R$ の
NCCR、かつ大域次元三である。
$I$ を双対に取り換えると $a,b$ が交換されるため、二つの ruling の選択を
固定しない場合には $\{a,b\}$ を考える。
また $A$ の単純加群はちょうど二種類で、その複素次元は $a,b$ である。
\end{proposition}

\begin{proof}
補題\ref{cmo:lem:formal-node-brauer}より
$\operatorname{Br}'(U)=0$ なので、$U$ 上の階数 $n$ のベクトル束
$\mathcal E$ を用いて
$A|_U\simeq\mathcal E nd(\mathcal E)$ と書ける。
ここで分裂を一般中心単純環に対する仮定として置く必要はない。

$\mathcal E$ を $\operatorname{Spec}R$ 上の連接加群 $N$ に延長し、
$M=N^{**}$ と置く。連接延長および以下で使う反射的延長は
\cite[Tag~0EBJ]{CMOStacksNode} の証明と主張による。
その仮定は $R$ が整域かつ Noetherian で、除いた閉点で
$\operatorname{depth}R=3\ge2$ であることから満たされる。
$M$ は有限生成反射的加群で、
$M|_U\simeq\mathcal E$、$\operatorname{rank}_R M=n$ である。
$R$ は正規で、$A$ は Cohen--Macaulay なので反射的である。
さらに反射的加群への $\operatorname{Hom}$ は反射的なので
（\cite[Tag~0AV3]{CMOStacksNode}）、
$\operatorname{End}_R(M)$ も反射的である。
余次元二以上の補集合を持つ開集合への制限は反射的加群上で完全忠実であるから、
$U$ 上の同型は一意に
\[
 A\simeq\operatorname{End}_R(M)
\]
へ延長する。積との整合性も $U$ 上で確認すればよい。
この段階では $M$ が Cohen--Macaulay であることも、
$R\in\operatorname{add}M$ であることも仮定していない。

$R$ は三次元の正規 Gorenstein 完備局所整域であり、特に
\cite{CMOIyamaWemyssMM} の意味で $3$-sCY である。
同論文の Lemma~2.23 により、Cohen--Macaulay 性と有限大域次元から
$\operatorname{End}_R(M)$ は NCCR である。
同論文の Proposition~4.5 により $M$ は maximal modifying (MM) 加群である。

次に $R$ を
$\mathbb C[[z,w,x,y]]/(zw-xy)$ と書き、
\cite[Theorems~5.1(2), 5.20]{CMOIyamaWemyssCA}
を適用する。ここでは $k=\mathbb C$、$f=zw$ の二つの既約因子は
$z,w$ であり、いずれも一次式なので定理の全仮定を満たす。
同定理は自由直和因子を仮定しない MM 加群も分類している。
従って $M$ の basic module は、ある階数一反射的加群 $J$ に対し
\[
 (J\otimes_R(R\oplus I_\varepsilon))^{**},\qquad
 I_\varepsilon\in\{(x,z),(x,w)\}
\]
の形である。二つのイデアルの類は互いに逆であり、
$(x,w)\simeq I^*$ である。
完備局所環上の Krull--Schmidt 分解を用いると、正整数 $a,b$ に対し
\[
 M\simeq J^{\oplus a}\oplus\bigl((J\otimes_R I_\varepsilon)^{**}\bigr)^{\oplus b}
\]
となる。

階数一反射的加群による反射的 tensor は、反射的加群の自己同値であり、
$\operatorname{Hom}$ と合成を保存する。
これは余次元一で線束による tensor となることを確認し、反射的延長を取ればよい。
$J^*$ による tensor により共通因子を除くと
\[
 A\simeq\operatorname{End}_R(R^{\oplus a}\oplus I_\varepsilon^{\oplus b}).
\]
$I_\varepsilon=I^*$ の場合にはさらに $I$ で tensor し、二つの直和因子を
入れ換えれば式\eqref{cmo:eq:conifold-order}になる。
階数から $(a+b)^2=\operatorname{rank}_R A=n^2$、従って $a+b=n$ である。

$R$ と $I$ は非同型の indecomposable 加群であり、
$\operatorname{End}_R(R)=\operatorname{End}_R(I)=R$ である。
従って Krull--Schmidt 分解に対応する Jacobson 根基の商は
\[
 A/\operatorname{rad}A\simeq
 \operatorname{Mat}_a(\mathbb C)\times
 \operatorname{Mat}_b(\mathbb C).
\]
これが単純加群についての主張を与える。
\end{proof}

\begin{theorem}[指定二因子の全非 Azumaya 単純加群]
\label{cmo:thm:fourteen-node-orders}
$\alpha=(1,4,9)$ または $(16,25,49)$ とする。
$X_\alpha=\operatorname{Spec}Z(D_\alpha)$ の十四個の節点 $p$ において、
$R=\widehat{\mathcal O}_{X_\alpha,p}$ を標準節点環と同定すると
\[
 \widehat{D_{\alpha,p}}
 \simeq\operatorname{End}_R(R^{\oplus a}\oplus I^{\oplus b}),\qquad a+b=6
\]
である。二つの ruling の交換を許した重複度は次の通りである。
\[
\begin{array}{c|c|c|c}
 \text{中心点の種類}&\text{中心点の個数}&\{a,b\}
 &\text{各点上の単純加群の次元}\\\hline
 \text{一次元単純加群の中心}&8&\{1,5\}&1,5\\
 \text{二次元単純加群の中心}&2&\{2,4\}&2,4\\
 \Omega=0\text{ 上の節点}&4&\{3,3\}&3,3
\end{array}
\]
特に各因子の非 Azumaya 中心上の単純加群は全部で二十八同型類であり、
次元 $1,2,3,4,5$ の個数は順に $8,2,8,2,8$ である。
正則中心の各閉点上には、六次元単純加群が一つある。
\end{theorem}

\begin{proof}
命題\ref{cmo:prop:center-critical-polynomials}より十四点は三次元通常二重点である。
形式的 Morse 補題により、その完備局所環は標準節点環 $R$ に同型である。
定理\ref{cmo:thm:azumaya-regular}より $D_\alpha$ は中心上
Cohen--Macaulay、一般階数三十六であり、正則中心上で次数六の Azumaya 代数である。
これらの性質は中心での局所化と完備化で保たれる。
また中心は有限個の代数生成元との交換子写像の共通の核として表されるので、
平坦な局所化と完備化により、完備化した代数の中心は正確に $R$ である。
局所中心の特異集合は閉点だけで、その完備化における逆像も閉点だけなので、
完備化した代数は穿孔スペクトラム上で Azumaya である。
$D_\alpha$ の左右の大域次元は三であり、
補題\ref{cmo:lem:completion-global-dimension}により完備化の大域次元も有限である。
従って命題\ref{cmo:prop:conifold-order}を $n=6$ に適用できる。

八つの一次元単純加群は全て異なる非零の $\Omega$ 値を持つ。
命題\ref{cmo:prop:two-dimensional-classification}および
命題\ref{cmo:prop:two-dimensional-center-rigidity}より、
二次元単純加群は各因子について正確に二つであり、その中心値も互いに異なり、
零および一次元加群の中心値と異なる。
これら十個の中心点は $H_i(\Omega^2)=0$ に対応する十節点を尽くす。
局所代数の二つの単純加群の次元和は六なので、
一次元加群のある点では $\{a,b\}=\{1,5\}$、
二次元加群のある点では $\{a,b\}=\{2,4\}$ となる。
残る四点では全一次元・二次元分類により、そのいずれも存在しない。
従って $a,b\ge3$、かつ $a+b=6$ から $a=b=3$ である。

元の代数の中心 $p$ 上の単純加群と完備化の単純加群は、
共通の有限次元商 $D_\alpha/\mathfrak m_pD_\alpha$ を通じて一致する。
これで元の $D_\alpha$ の単純加群に関する表と個数が従う。
正則中心上の主張は Azumaya 性による。
\end{proof}

\begin{remark}[用いた既知結果と本例に固有の入力]
完全局所 $cA_n$ 環上の MM 分類と NCCR の一般論は
Iyama--Wemyss の既知結果である。
本例でこれらを適用するための入力は、中心の十四節点、
中心上の Cohen--Macaulay 性、正則中心上の Azumaya 性、PI 次数六、
および既に得た一次元・二次元単純加群の完全分類である。
穿孔完備節点の Brauer 群、反射的 endomorphism 表示、完備化後の有限大域次元を
上の補題で確認したため、generic CSA の分裂や自由直和因子の存在を
追加仮定として残していない。
この局所同定のみから任意の中心切断の有限大域次元は結論せず、
指定中心関数の線形項と節点の二次形式の関係は別に検証する。
\end{remark}


\subsection{通常二重点上の非可換解消の中央切断}
\label{cmo:subsec:conifold-slice}

以下の完備局所模型について、加群の大域次元と形式的対角の perfect 性を区別する。
形式的対角には完備テンソル積を用いる。
形式冪級数環を、通常の $\C$ 上のテンソル積に関して有限型の滑らかな代数と扱うことはしない。

\begin{proposition}[非等方な中央切断]
\label{cmo:prop:conifold-nonisotropic-slice}
\[
 R=\C[[x,y,z,w]]/(xy-zw),\qquad
 I=(x,z),\qquad \Lambda=\operatorname{End}_R(R\oplus I)
\]
とする。$f\in\mathfrak m_R$ の一次部分を
$\ell=a x+b y+c z+d w$ と書き、
\begin{equation}\label{cmo:eq:conifold-dual-nonisotropic}
 ab-cd\ne0
\end{equation}
を仮定する。
このとき、$f$ の二次以上の項にかかわらず
\[
 \Lambda/(f)\simeq\C[[p,q]]\rtimes\mu_2,
 \qquad -1\cdot(p,q)=(-p,-q)
\]
という非標準的な同型があり、その左右大域次元は二である。
さらにこの代数の形式的対角は、完備包絡環上で長さ二の有限生成射影分解を持つ。
$\operatorname{End}_R(R^{\oplus r}\oplus I^{\oplus s})$、$r,s>0$ についても
その中央切断は同じ代数と Morita 同値であり、同じ結論が成立する。
\end{proposition}

\begin{proof}
二次形式 $xy-zw$ を $\ker\ell$ に制限した三変数二次形式が非退化であることは
式\eqref{cmo:eq:conifold-dual-nonisotropic}と同値である。
$\ell\ne0$ なので $f$ を一つの形式座標として使い、$f=0$ で消去する。
残る方程式の二次部分は上の非退化二次形式である。
形式的 Morse 補題により
\[
 S=R/(f)\simeq\C[[X,Y,Z]]/(XY-Z^2)
\]
を得る。この議論は $f$ の高次項を任意に許す。

$R\oplus I$ と $\Lambda$ は $R$ 上 Cohen--Macaulay である。
実際、$I$ とその双対は通常二重点の階数一 Cohen--Macaulay 加群であり、
$\Lambda$ の四ブロックは $R,R,I,I^*$ である。
$R$ は整域で $f\ne0$ なので、$f$ はこれら全てに正則に作用する。
$J=I/fI$ と置くと、$J$ は $S$ 上の階数一 Cohen--Macaulay 加群である。
自然な射
\begin{equation}\label{cmo:eq:conifold-end-base-change}
 \Lambda/f\Lambda\longrightarrow
     \operatorname{End}_S(S\oplus J)
\end{equation}
は同型である。これを Hom と商の交換として無条件に使わず、次のように確認する。
両辺は正規二次元環 $S$ 上で反射的である。
左辺の反射性は深さ二、右辺の反射性は $S\oplus J$ の反射性による。
閉点を除けば $I$ は可逆なので、上の射は行列環の底変更として同型になる。
正規環上の反射的加群は余次元一の制限から回復されるので、射全体が同型となる。

$f\in\mathfrak m_R$ より
\[
 \dim_\C J/\mathfrak m_SJ
 =\dim_\C I/\mathfrak m_RI=2.
\]
従って $J$ は自由でない。
$A=\C[[p,q]]$、$G=\mu_2$ とすると $S=A^G$ であり、
$S$ の階数一反射的加群の同型類は $\operatorname{Cl}(S)\simeq\mathbb Z/2$ で分類される。
この最後の記述も本例では直接確認できる。
閉点を除いた $\operatorname{Spec}A$ は Picard 群が零であり、
$G$ はその上で自由に作用する。
従って降下する直線束は $G$ の二つの指標で与えられる。
より具体的には $A^\times=\C^\times\times(1+\mathfrak m_A)$ と形式対数を使うと、
平均化により $H^1(G,1+\mathfrak m_A)=0$ である。
よって $J$ は奇部分 $A^-$ と同型であり、$S\oplus J\simeq A$ となる。

自然な射 $A\rtimes G\to\operatorname{End}_{A^G}(A)$ も、
閉点の外で有限エタール Galois 降下による同型である。
両辺の反射性により閉点を越えて同型となる。
これで式\eqref{cmo:eq:conifold-end-base-change}から主張した表示を得る。
この同型は本例における Auslander の対応の直接証明でもある。

$|G|=2$ は可逆である。
任意の $A\rtimes G$ 加群 $N$ は、平均化によって
$(A\rtimes G)\otimes_A N$ の直和因子になる。
$A$ の大域次元は二であり、$A\rtimes G$ は左右 $A$ 自由なので、
$\operatorname{pd}_{A\rtimes G}N\le2$ を得る。
原点における単純加群を $A$ に制限すると剰余体 $\C$ になるので、
逆の下限も成立し、大域次元は正確に二である。

形式的対角の主張には、二変数の完備 Koszul 分解を $G$ 同変化する。
各項は $\C G$ を介した誘導双加群であり、$\C G$ の分離性により
完備包絡環上有限生成射影となる。
長さは二で、随伴次数付き複体は通常の二変数 Koszul 分解だから完全である。
最後の Morita 同値は $R\oplus I$ と $R^{\oplus r}\oplus I^{\oplus s}$ が同じ加法閉包を持つことから従う。
\end{proof}

\begin{remark}[一次項が等方な場合]
\label{cmo:rem:conifold-isotropic-slice}
非等方性は省略できない。
特に $f=\ell\ne0$ が線形で $ab-cd=0$ なら、$\Lambda/(f)$ の大域次元は無限である。
これを中央切断の可換特異性だけから結論するのではなく、quiver で確かめる。
標準 conifold 代数は、二頂点、矢
$a_1,a_2:0\to1$、$b_1,b_2:1\to0$ と
\[
 a_i b_j a_k=a_k b_j a_i,\qquad
 b_i a_j b_k=b_k a_j b_i
\]
の関係で与えられる。
中心の四つの生成元は $z_{ij}=a_i b_j+b_j a_i$ であり、
$\det(z_{ij})=0$ である。
この表示は $a_1=x,a_2=z,b_1=1,b_2=w/x$ という
$R\oplus I$ の間の乗法写像で得られる。
三次関係によって交互語中の $a$ 同士、$b$ 同士を並べ替えられ、
長さ $2n$ の各対角ブロックの次元は $(n+1)^2$、
長さ $2n+1$ の各非対角ブロックの次元は $(n+1)(n+2)$ になる。
これは上記 End 環の Hilbert 級数と一致するので、表示には余分な関係はない。

矢の基底変換により、階数一の線形中央元を $z_{11}$ にできる。
商では $a_1b_1=b_1a_1=0$ である。
他の二矢を零にする射と $a_1,b_1$ を残す包含によって、
二頂点二周期の radical-square-zero 代数が代数 retract になる。
その頂点半単純代数上の bar 複体は全次数に非零項を持ち、微分は零である。
retract によりその Tor は商代数の Tor の直和因子となるから、大域次元は無限である。

高次項を許しても、一次部分が非等方なら前命題がそのまま適用できる。
一次部分が等方、または零の場合には、前命題からの結論は出さない。
例えば等方な一次部分に高次項を加えると、可換切断は $A_2$ 以上の型になり得るため、
切断後の加群と End 環を別に調べる必要がある。
\end{remark}

\begin{lemma}[指定中心の全節点における非等方性]
\label{cmo:lem:all-nodes-nonisotropic}
二つの指定中心 $X_i=V(P_i(s^2,U,V,T))$ の各十四節点で、
$s-s_0$ の一次部分は通常二重点の双対二次形式に関して非等方である。
\end{lemma}
\begin{proof}
$H_3=\det\operatorname{Hess}_{U,V,T}P_i$、
$H_4=\det\operatorname{Hess}_{W,U,V,T}P_i$ と置く。
再現コード \path{computations/verify_mixed_center_geometry.py} は、既存の検証済み関係式を入力して
\[
 J_i=(P_i,P_{i,W},P_{i,U},P_{i,V},P_{i,T})
\]
上で $H_3,H_4$ が単元であることを有理数上の Gr\"obner 消去で検査する。
$W$ も単元であり、これは消去多項式 $H_i$ が零を根に持たないことから従う。
ここで $\dim_\mathbb Q\mathbb Q[W,U,V,T]/J_i=5$ である。
従来の検査に追加される条件は $H_3$ の単元性である。
$s_0\ne0$ の十節点では Hessian は
\[
 \operatorname{Hess}_{s,U,V,T}P_i(s^2,U,V,T)
 =\operatorname{diag}(2s_0,1,1,1)
   \operatorname{Hess}_{W,U,V,T}P_i
   \operatorname{diag}(2s_0,1,1,1)
\]
となる。その逆行列の $(s,s)$ 成分は
$H_3/(4s_0^2H_4)\ne0$ である。

$s_0=0$ の四節点では、既存の長さ四の被約環上で
$P_{i,W}$ と $H_3|_{W=0}$ がともに単元である。
Hessian は $2P_{i,W}$ と三変数 Hessian の直和となり、
逆行列の $(s,s)$ 成分は $1/(2P_{i,W})\ne0$ である。
これらはまさに共変ベクトル $ds$ の双対二次形式による値である。
\end{proof}

% !TEX root = ../../mixed-sklyanin-cubics.tex
\section{有限台加群の perfect 性と残余圏の非滑らかさ}
\label{mix:sec:perfect-modules}

\subsection{有限台の chart 加群と線束 perfect 性}

以下でいう線束 perfect とは
\[
 \mathcal C_A=\operatorname{thick}\{\mathcal O(j):j\in\mathbb Z\}
 \subset D^b(\operatorname{qgr}A)
\]
に属することをいう。この部分圏の proper 性は、全 coherent 圏の
smoothness を仮定せずに成立する。

\begin{lemma}[有限中心上の有限台 chart 加群]
\label{cmo:lem:finite-support-perfect}
$A$ を次数一で生成される連結局所有限 noetherian 次数付き代数とし、
$u\in Z(A)_1$ を正則元とする。
$u$ を含む可換有限型次数付き中心部分環 $C\subset A$ があり、
$A$ は $C$ 上有限であると仮定する。
$D=(A[u^{-1}])_0$ とし、有限次元右 $D$-加群 $M$ に対して
補題\ref{cmo:lem:chart-ext}の $\mathcal P(M)=\pi(M[u])$ を取る。
このとき
\begin{equation}\label{cmo:eq:chart-perfect-iff}
 \mathcal P(M)\in\mathcal C_A
 \quad\Longleftrightarrow\quad
 \operatorname{pd}_{D}M<\infty.
\end{equation}
\end{lemma}
\begin{proof}
$\ell>0$ を十分割り切れる整数として、$C^{(\ell)}$ が次数一で生成されるようにする。
$Z=\operatorname{Proj}C^{(\ell)}$、
$\mathcal A=\widetilde{A^{(\ell)}}$ と置く。
定理\ref{cmo:thm:relative-smoothness}の証明の有限中心模型と同じく、
\[
 \operatorname{QGr}A\simeq\operatorname{Qcoh}(Z,\mathcal A),\qquad
 \operatorname{qgr}A\simeq\operatorname{coh}(Z,\mathcal A)
\]
がある。これは smoothness の結論に先立つ同値である。
また、次数 $\ell$ の中央元 $f$ を反転した $A[f^{-1}]$ は強次数付きである。
実際、$f$ を次数一の元の長さ $\ell$ の積の和として書き、
先頭または末尾の元を分ければ
$1\in A[f^{-1}]_1A[f^{-1}]_{-1}$ および
$1\in A[f^{-1}]_{-1}A[f^{-1}]_1$ を得る。
従って全ての $\mathcal O(j)$ は局所有限生成射影 $\mathcal A$-加群である。
中央の ample twist $\mathcal A\otimes\mathcal O_Z(m)$ は
$\mathcal O(\ell m)$ に対応し、これらは
$D_{\mathrm{qc}}(Z,\mathcal A)$ を compact に生成する。
有限アファイン被覆上での局所有限射影分解と同じ生成論から、
\begin{equation}\label{cmo:eq:line-is-local-perfect}
 \mathcal C_A=\operatorname{Perf}(Z,\mathcal A)
\end{equation}
である。混合族では、命題\ref{cmo:prop:proper-contraction}の有限線束和 $E$
がこれら全 twist を厚生成するため、同じ等式はその compact generator
の記述からも従う。

ここで $P=M[u]$ 上の中心作用を調べる。
各 $c\in C_e$ は、$M$ 上の $cu^{-e}$ の作用行列に $u^e$ を掛けたものとして作用する。
従って中心像は、$u$ を含む次数付き可換部分環
\[
 S=C/\operatorname{Ann}_C P
 \ \subset\ \operatorname{End}_{\mathbb C}(M)[u]
\]
である。右辺は $\mathbb C[u]$ 上有限であり、
$S$ はその $\mathbb C[u]$-部分加群だから、$S$ も有限である。
従って $S/(u)$ は有限次元であり、
\[
 \operatorname{Proj}S^{(\ell)}\subset Z
\]
は $u\ne0$ 内の有限閉部分スキームである。
より具体的には、$S[u^{-1}]_0\subset\operatorname{End}_{\mathbb C}(M)$ は
有限次元可換代数で、上の Proj はその Spec に等しい。
$\mathcal P(M)$ の中央台はこの有限閉集合に含まれる。

$U=D_+(u^\ell)$ では $\mathcal A|_U$ のアファイン代数は $D$ であり、
$\mathcal P(M)|_U$ は $M$ に対応する。
もし $\operatorname{pd}_D M<\infty$ なら $U$ 上で perfect である。
中央台の補集合では零なので、$Z$ 全体で局所 perfect となる。
式\eqref{cmo:eq:line-is-local-perfect}から順方向を得る。

逆に $j^*:\operatorname{QGr}A\to\operatorname{Mod}D$ は exact で、
$j^*\mathcal O(j)\simeq D$、$j^*\mathcal P(M)=M$ である。
従って $\mathcal P(M)\in\mathcal C_A$ なら
$M\in\operatorname{Perf}(D)$、すなわち有限生成加群 $M$ の射影次元は有限である。
これは、任意の perfect 加群を開埋込みの $j_*$ が perfect に保つという主張ではない。
順方向で不可欠なのは、有限次元性から得た、開 chart 内で閉じた中央台である。
\end{proof}

\begin{proposition}[定常点の線束 perfect 性の完全判定]
\label{cmo:prop:point-perfect}
命題\ref{cmo:prop:character-pd}の $k$ 因子の混合型を取り、
$p_i=(a_i,b_i,c_i)\in(\mathbb C^\times)^3$ とする。
中心商の係数を $c_0=-\sum_i\beta_i\lambda_{p_i}$ とし、
対応する次数付き一点加群 $P=\mathbb C[u]$ と
$\mathcal P=\pi P$ を取る。このとき
\begin{equation}\label{cmo:eq:point-perfect-criterion}
 \mathcal P\in\mathcal C_A
 \quad\Longleftrightarrow\quad
 \text{少なくとも一つの }\kappa(p_i)\ne0.
\end{equation}
有限の場合、対応する右 affine character の射影次元は正確に $3k-1$ である。
特に例\ref{cmo:ex:sixty-four}の全六十四個の定常点は線束 perfect である。
\end{proposition}
\begin{proof}
最初に左右の規約を揃える。各因子に対して、語の順序を逆転させるとともに
\[
 x\longmapsto x,\qquad y\longmapsto z,\qquad
 z\longmapsto y,\qquad u\longmapsto u
\]
とする同型
\[
 \mathcal B_{(a^2,b^2,c^2)}
 \xrightarrow{\sim}
 \mathcal B_{(a^2,c^2,b^2)}^{\mathrm{op}}
\]
がある。三つの関係式を直接代入すれば確認でき、
中心元 $\Omega$ と character はそれぞれ対応する中心元と $(a,c,b)$ へ移る。
$\lambda$ と $\kappa$ は三座標の置換で不変だから、
命題\ref{cmo:prop:character-pd}の左射影次元の結論は右加群にも移る。
有限次元体双対による自己 Ext の比較だけから右射影次元を推論しているのではない。

定理\ref{cmo:thm:finite-center}の普遍中心元の像が生成する中心部分環を使うと、
$A$ は $u$ を含む有限型次数付き中心上有限である。
従って前補題と右 affine character の射影次元判定が適用できる。
一点の場合は中心像がちょうど $\mathbb C[u]$ なので、
中央台は $u\ne0$ 内の一つの閉点である。
六十四点では系\ref{cmo:cor:sixty-four-ext}の二つの不等式により
両因子の $\kappa$ がともに非零である。
\end{proof}

\begin{corollary}[低次元単純加群の coherent な perfect 延長]
\label{cmo:cor:hundred-perfect}
命題\ref{cmo:prop:one-two-dimensional-pd}の百個の左単純加群を体双対して
右加群 $M$ とし、各出現係数の $A_c$ 上で $\mathcal P(M)=\pi(M[u])$ と置く。
これらは全て線束 perfect である。
\end{corollary}
\begin{proof}
体双対は有限次元端点の自己 Ext を、反対向きの自己 Ext に同型に移す。
例えば有限生成射影分解を体双対した単射分解と Hom の随伴からこの同型を得る。
右側の周囲 tensor 代数も有限大域次元六を持ち、$M$ は有限生成なので、
長さ六の有限生成射影分解がある。
従って、双対した後に系\ref{cmo:cor:hypersurface-pd}を右加群へ適用すれば
有限自己 Ext から $\operatorname{pd}M=5$ を得る。
体双対だけから射影次元を移してはいない。
補題\ref{cmo:lem:finite-support-perfect}を適用して結論を得る。
\end{proof}

\begin{remark}[判定の範囲]
この結果は定常点そのものが有限線束複体の直和因子になることを示す。
全 coherent 対象の perfect 性や、全カテゴリー・対角の smoothness を意味しない。
全 $\kappa=0$ の場合には、既存の非有界自己 Ext と同じく
$\mathcal P\notin\mathcal C_A$ である。
また $\operatorname{RHom}(\mathcal O,\mathcal P)=\mathbb C$ なので、
$\mathcal P$ 自体は残余圏の対象ではない。
\end{remark}

\subsection{二ブロックの定常点を残余圏へ射影する}

次の Serre 入力は全 coherent 圏の smoothness を要求しない。

\begin{lemma}[混合型の線束圏上の Serre 双対]
\label{cmo:lem:mixed-line-serre}
$N=3k+1$、$n=N-2$、$a=N-3$ とする。
本節の混合中心三次商の proper 圏 $\mathcal C_A$ は
\[
 S_{\mathcal C_A}=(-a)[n]
\]
を Serre 関手として持つ。特に線束 perfect な定常点では
$S_{\mathcal C_A}(\mathcal P)\simeq\mathcal P[n]$ である。
\end{lemma}
\begin{proof}
命題\ref{cmo:prop:proper-contraction}の証明は、
各因子の potential 解消の tensor と中心差 $u_L-u_R$ の Koszul 解消から
\[
 R\underline{\operatorname{Hom}}_{B\otimes B^{\mathrm{op}}}
     (B,B\otimes B^{\mathrm{op}})\simeq B(N)[-N]
\]
を真正な次数付き双加群として構成している。
この構成は中心商の係数 $c_0$ に依存しない。
正則中心三次による change of rings と両側局所双対により
\[
 R\Gamma_A(A)\simeq A^\vee(a)[-(N-1)]
\]
を得る。
これは、定理\paperref{moduli}{thm:nfm-relative-serre}の証明で使う
局所双対式\papereqref{moduli}{eq:nfm-relative-local-dual}と同じ入力であり、
一ブロックという表示はこの入力を得た後には使わない。
同証明の導来局所化の三角を双対化すると、tails の双対化双加群は
内部 shift $-a$ と cohomological shift $N-2$ を持つ。
従って全 compact 対象 $K,L$ について
\[
 R\operatorname{Hom}(K,L)^\vee
 \simeq R\operatorname{Hom}(L,K(-a)[N-2])
\]
となる。有限中心模型では compact 部分は $\mathcal C_A$ であり、
この圏は proper であるから、これが主張した Serre 双対である。
$u:P(-1)\to P$ の余核は有限次元なので
$\mathcal P(-1)\simeq\mathcal P$ であり、最後の結論を得る。
\end{proof}

\begin{proposition}[残余射影の自己 Ext]
\label{cmo:prop:projected-point-ext}
$k=2$ とし、少なくとも一つの $\kappa(p_i)$ が非零である定常点
$\mathcal P$ を取る。
$\mathcal E=\langle\mathcal O,\mathcal O(1),\mathcal O(2),\mathcal O(3)\rangle$ と置き、
\[
 F\longrightarrow\mathcal P\longrightarrow K
\]
を $F\in\mathcal E$、$K\in\mathcal E^\perp\cap\mathcal C_A
=\mathcal K_A^{\mathrm{line}}$ となる射影三角とする。
このとき
\begin{equation}\label{cmo:eq:projected-point-ext}
 \operatorname{Ext}^j(K,K)=
 \begin{cases}
  \mathbb C,&j=0,\\
  \mathbb C^{10},&j=1,\\
  0,&j\ne0,1.
 \end{cases}
\end{equation}
特に $\chi(K,K)=-9$ であり、$K$ は exceptional 対象でも spherical 対象でもない。
この結論は全六十四個の定常点に適用できる。
\end{proposition}
\begin{proof}
\textbf{第一段階：斉次分解と例外射影。}
一因子では $X=x-au,Y=y-bu,Z=z-cu$ と置き、
式\eqref{cmo:eq:character-resolution}の行列の定数 $a,b,c$ を
$au,bu,cu$ へ置換する。これにより
\[
 0\to\mathcal B_\alpha(-3)\to\mathcal B_\alpha(-2)^3
 \to\mathcal B_\alpha(-1)^3\to\mathcal B_\alpha
 \to\mathbb C[u]\to0
\]
を得る。$u=0$ では $S(1,-1,-1)$ の Koszul 解消となる。
増大複体の全項が $\mathbb C[u]$ 上平坦であり、内部次数は下に有界なので、
$u$ による帰納法で完全性が持ち上がる。
二因子の $\mathbb C[u]$ 上の tensor は、右加群の規約へ移せば
\[
 B(-i)^{\binom6i}\qquad(0\leq i\leq6)
\]
を第 $i$ 項とする $P$ の有限自由 $B$-分解となる。

斉次な higher homotopy を使う Shamash 分解 $Q_\bullet\to P$ の項は
\[
 Q_m=\bigoplus_{q\ge0}A(-m-q)^{\binom6{m-2q}},
\]
ただし範囲外の二項係数は零とする。
$Q_\bullet(3)$ のうち、自由生成元の line twist が非負の項を取ると、
微分で閉じた有限部分複体
\begin{equation}\label{cmo:eq:point-exceptional-complex}
 F=\left[
 \mathcal O^{20}\longrightarrow
 \mathcal O(1)^{15}\oplus\mathcal O\longrightarrow
 \mathcal O(2)^6\longrightarrow\mathcal O(3)
 \right]
\end{equation}
を得る。表示した項の cohomological 次数は順に $-3,-2,-1,0$ である。
増大写像と $\mathcal P(3)\simeq\mathcal P$ により $F\to\mathcal P$ がある。

$0\leq i,j\leq3$ では線束間の高次 Ext は零である。
従って $R\operatorname{Hom}(\mathcal O(i),F)$ は項ごとの Hom で計算できる。
これは自由分解 $Q_\bullet$ の内部次数 $3-i$ の成分そのものである。
実際、$F$ から除いた項の line twist は負なので、この成分に寄与しない。
従ってこの Hom 複体は次数零の $P_{3-i}=\mathbb C$ のみを持つ。
また exact な $j^*\dashv j_*$ により
\[
 R\operatorname{Hom}(\mathcal O(i),\mathcal P)
 =R\operatorname{Hom}_D(D,M)=\mathbb C
\]
であり、上の増大写像はこの同型を与える。
よって全ての $i=0,1,2,3$ について $R\operatorname{Hom}(\mathcal O(i),K)=0$。
命題\ref{cmo:prop:point-perfect}により $\mathcal P$ は線束 perfect なので、
$K$ も線束 perfect である。
従って式\eqref{cmo:eq:point-exceptional-complex}は所要の例外射影を実際に表す。

\textbf{第二段階：$R\operatorname{Hom}(F,\mathcal P)$ と制限写像。}
二因子の affine ambient character 分解に Hom を施した複体を
$(C^\bullet,\delta)$ と書く。
各因子は次元 $(1,3,3,1)$、唯一の非零微分が二つの三次元空間の間の
$\operatorname{diag}(-2a_i,-2b_i,-2c_i)$ である。
従って tensor 複体では
\[
 \dim(C^0,C^1,C^2,C^3)=(1,6,15,20),\qquad
 \delta:C^1\hookrightarrow C^2,\qquad
 \operatorname{rank}(\delta:C^2\to C^3)=9.
\]
全 Shamash cochain 複体は
$C^\bullet[t]$、$|t|=2$、微分
$\delta+t h_1^*+t^2h_2^*+\cdots$ である。
式\eqref{cmo:eq:point-exceptional-complex}に対応する Hom 複体は
\[
 C^0\longrightarrow C^1\longrightarrow C^2\oplus tC^0
 \longrightarrow C^3.
\]
最初の微分は零、次の微分は $\delta:C^1\to C^2$ の成分が単射なので階数六、
最後の微分は $\delta:C^2\to C^3$ なので階数九である。
従って
\begin{equation}\label{cmo:eq:exceptional-to-point-ext}
 \operatorname{Ext}^j(F,\mathcal P)=
 \begin{cases}
 \mathbb C,&j=0,2,\\
 \mathbb C^{11},&j=3,\\
 0,&\text{その他}.
 \end{cases}
\end{equation}

自然な写像
\[
 \operatorname{Ext}^j(\mathcal P,\mathcal P)
 \longrightarrow\operatorname{Ext}^j(F,\mathcal P)
\]
は、全 Shamash cochain 複体から上の有限複体への制限で計算できる。
これは $u\ne0$ への制限と $j^*\dashv j_*$ の随伴により確認できる。
次数零では恒等写像である。
次数二では両複体の次数一・二の項が同じなので、cohomology 上の写像は単射である。
両側の次元は一だから同型となる。
次数三の単射性も直接確認できる。
制限が境界となる全 cocycle から、その境界の持上げを引くと、
残りは $t v$、$v\in C^1$ の形である。
その微分の $tC^2$ 成分は $t\delta v$ である。
$\delta:C^1\to C^2$ が単射だから $v=0$ となる。
従って次数三でも制限写像は単射である。

\textbf{第三段階：proper 線束圏での双対と cone。}
前補題を $k=2$ に適用すると $S=(-4)[5]$ であり、
$\mathcal P(-4)\simeq\mathcal P$ である。
従って式\eqref{cmo:eq:exceptional-to-point-ext}を双対化して
\[
 \operatorname{Ext}^j(\mathcal P,F)=
 \begin{cases}
 \mathbb C^{11},&j=2,\\
 \mathbb C,&j=3,5,\\
 0,&\text{その他}
 \end{cases}
\]
を得る。Serre 双対の自然性により
$\operatorname{Ext}^j(\mathcal P,F)\to
\operatorname{Ext}^j(\mathcal P,\mathcal P)$ は、
第二段階の次数 $5-j$ の写像の双対である。
従って次数二では全射、次数三・五では同型となる。
この双対を適用している圏は proper な $\mathcal C_A$ であり、
全 $D^b(\operatorname{qgr}A)$ の proper 性や smoothness は使わない。

$R\operatorname{Hom}(F,K)=0$ なので射影三角から
\[
 R\operatorname{Hom}(K,K)\simeq R\operatorname{Hom}(\mathcal P,K)
 \simeq\operatorname{Cone}\bigl(
 R\operatorname{Hom}(\mathcal P,F)\to
 R\operatorname{Hom}(\mathcal P,\mathcal P)\bigr).
\]
既存の自己 Ext の次元 $(1,1,1,1)$（次数 $0,2,3,5$）と、
今得た写像の階数を長完全列へ代入すると、
次数零に一、次数一に $11-1=10$ のみが残る。
これが式\eqref{cmo:eq:projected-point-ext}である。
\end{proof}

\begin{remark}[数値的な独立検算と積について]
例外列の Euler 行列は
\[
 U=\begin{pmatrix}
 1&7&28&83\\0&1&7&28\\0&0&1&7\\0&0&0&1
 \end{pmatrix}.
\]
$\chi(\mathcal O(i),\mathcal P)=1$ から
$U(c_0,c_1,c_2,c_3)^{\mathsf t}=(1,1,1,1)^{\mathsf t}$ を解くと
$(c_0,c_1,c_2,c_3)=(-19,15,-6,1)$ となり、
\[
 [K]=[\mathcal P]+19[\mathcal O]-15[\mathcal O(1)]
                 +6[\mathcal O(2)]-[\mathcal O(3)]
\]
を得る。$\chi(\mathcal P,\mathcal P)=0$、
$\chi(\mathcal P,\mathcal O(i))=-1$ と合わせると $\chi(K,K)=-9$ であり、
上の全次数計算と一致する。
この証明は、自己 Ext の計算に使った外積表示を Yoneda 積と同定していない。
また $\mathcal P$ の自己 Ext と残余射影 $K$ の自己 Ext は異なる。
\end{remark}


\subsection{定常 norm とは異なる六十四個の非滑らかな係数}

\begin{theorem}[Azumaya 部分の Morse 点から生じる中心商の特異性]
\label{cmo:thm:morse-singular-sixty-four}
例\ref{cmo:ex:sixty-four}の二因子を取る。
命題\ref{cmo:prop:center-critical-polynomials}で定める四次多項式を $Q_1(W),Q_2(W)$ とし、
\begin{equation}\label{cmo:eq:morse-resultant}
 P_{64}(c)=\operatorname{pp}_+
   \operatorname{Res}_{z}\bigl(Q_1((-c-100z)^2),Q_2(z^2)\bigr)
       \in\mathbb Z[c]
\end{equation}
と置く。$\operatorname{pp}_+$ は整数係数の最大公約数を除き、
最高係数を正とする規格化である。
この多項式は次数六十四の平方因子を持たない偶多項式であり、
\[
 P_{64}(0)\ne0,\qquad \gcd(P_{64},N)=1
\]
を満たす。ここで $N$ は定常 character の norm である。

$P_{64}(c)=0$ の各係数に対して、$A_c$ の中心 chart は
三十六次元単純右加群 $M$ を持ち、その完備中央近傍は
\begin{equation}\label{cmo:eq:morse-matrix-node}
 \operatorname{Mat}_{36}\!\left(
   \C[[v_1,\ldots,v_6]]/(v_1^2+\cdots+v_6^2)\right)
\end{equation}
に同型である。その coherent な延長 $\mathcal P(M)=\pi(M[u])$ は
\begin{equation}\label{cmo:eq:morse-self-ext}
 \sum_{j\ge0}\dim_\C\operatorname{Ext}_{\operatorname{qgr}A_c}^j
       (\mathcal P(M),\mathcal P(M))s^j
       =\frac{(1+s)^6}{1-s^2}
\end{equation}
を満たす。
従って、これら全六十四係数で qgr の homological dimension は無限、
$D^b(\operatorname{qgr}A_c)$ は非 proper、線束生成圏とその残余部分は非 smooth であり、全 coherent 残余圏も非 proper である。
特に、この六十四係数は滑らかな可換 cubic への全 coherent Ku 比較が
成立する係数には含まれない。
\end{theorem}
\begin{proof}
$Q_i$ の根は零でなく相異なり、対応する $\Omega_i$ の八個の値
\[
 \mathcal M_i=\{\lambda\in\C:Q_i(\lambda^2)=0\}
\]
は、中心の正則部分における非退化臨界点の値である
（命題\ref{cmo:prop:center-critical-polynomials}）。
終結式の定義より、$P_{64}$ の根集合は
$-\mathcal M_1-100\mathcal M_2$ である。
次数、平方因子の不存在、零および norm 根との不交性は、
二つの四次多項式から有理数体上の終結式と最大公約数を計算して確認する。
再現コードは \texttt{verify\_mixed\_center\_geometry.py} であり、
係数の浮動小数点近似や根の数値分離は使わない。
次数六十四で平方因子がないため、八行八列の和の六十四値は全て相異なる。

$c=-\lambda-100\mu$ に対応する二つの臨界点 $z_1,z_2$ を取る。
定理\ref{cmo:thm:azumaya-regular}から、各 $D_{\alpha^{(i)}}$ は
$z_i$ の近傍で次数六の Azumaya 代数である。
その外積は、滑らかな六次元中心の $z=(z_1,z_2)$ の近傍で
次数三十六の Azumaya 代数になる。
$\Omega_1+100\Omega_2+c$ の Hessian は二つの非退化な三次元 Hessian の
block 対角和であり、係数 $100\ne0$ なので非退化である。
標数零の形式的 Morse 補題により、その方程式は形式座標変更の下で
$v_1^2+\cdots+v_6^2$ となる。
Azumaya 代数を剰余体 $\C$ の完備局所環上で分裂させ、中心商を取れば
式\eqref{cmo:eq:morse-matrix-node}を得る。

この行列代数の単純剰余加群は次元三十六である。
完備化前にも、対応する中心点上の唯一の単純加群 $M$ として定義される。
局所有限自由分解と完備化の平坦性により、$M$ の自己 Ext は完備化後と一致する。
実際、中心の極大イデアルは Hom 複体の標的 $M$ を零化し、
全ての自己 Ext も同じ剰余体上のベクトル空間なので、完備化で変わらない。

Morita 同値の下でこれは五次元通常二重点の剰余体の自己 Ext である。
六変数正則局所環の Koszul 分解に、極大イデアルの二乗に属する一つの二次式の
Shamash 変数を加えた分解は極小であり、その Betti 級数は
$(1+s)^6/(1-s^2)$ となる。
補題\ref{cmo:lem:chart-ext}で全次数の qgr 自己 Ext へ移すと
式\eqref{cmo:eq:morse-self-ext}を得る。

\paragraph{有限中央台に対する両方向の Ext の局所性.}
補題\ref{cmo:lem:finite-support-perfect}の有限中心模型
$Z=\operatorname{Proj}C^{(\ell)}$、$\mathcal A=\widetilde{A^{(\ell)}}$ を使う。
有限次元 chart 加群 $M$ の延長 $\mathcal P=\mathcal P(M)$ の中央台は、
$U=D_+(u^\ell)$ に含まれる有限閉集合である。
従って任意の coherent $\mathcal A$-加群 $\mathcal F$ について
\begin{equation}\label{cmo:eq:finite-support-local-ext}
 \operatorname{Ext}^j_{\operatorname{qgr}A}(\mathcal P,\mathcal F)
 \simeq\operatorname{Ext}^j_D(M,j^*\mathcal F)
 \qquad(j\ge0)
\end{equation}
となる。実際、各層 $\mathcal E xt^j_{\mathcal A}(\mathcal P,\mathcal F)$ は
coherent で、その台はこの有限閉集合に含まれる。
したがって高次層 cohomology は零であり、局所-to-大域 Ext spectral sequence は
次数ごとに退化する。その大域切断は $U$ 上の切断に等しく、
$U$ は affine なので右辺となる。
この議論は $\mathcal P$ の perfect 性を仮定しない。
有限中心模型と coherent 部分の導来完全忠実性は既存の補題と同じものを使っている。

\paragraph{節点上の単純加群から coherent 残余圏の非 proper 性へ.}
式\eqref{cmo:eq:morse-matrix-node}の単純加群 $M$ に対し、
\[
 R=\mathbb C[[v_1,\ldots,v_6]]/(v_1^2+\cdots+v_6^2)
\]
と置く。六変数正則局所環の Koszul 双対と中心正則元の Rees 同型から
\[
 R\operatorname{Hom}_R(\mathbb C,R)\simeq\mathbb C[-5]
\]
である。$\operatorname{Mat}_{36}(R)$ の右正則加群は Morita 同値で
$R^{36}$ に対応し、その単純加群は $\mathbb C$ に対応する。
したがって局所化・完備化の忠実平坦性により
\[
 \operatorname{Ext}^j_D(M,D)=
 \begin{cases}
  \mathbb C^{36},&j=5,\\
  0,&j\ne5.
 \end{cases}
\]
有限生成加群の有限階数自由分解を用いれば、各次数の Ext は完備化と可換である。
またその中央台は $M$ の一点の台に含まれるため、
この一点での局所化・完備化の計算が元の Ext を決める。
式\eqref{cmo:eq:finite-support-local-ext}と $j^*\mathcal O(i)\simeq D$ から、
全ての $i$ について
\begin{equation}\label{cmo:eq:node-point-to-line}
 \operatorname{Ext}^j(\mathcal P,\mathcal O(i))=
 \begin{cases}
  \mathbb C^{36},&j=5,\\
  0,&j\ne5
 \end{cases}
\end{equation}
を得る。
逆方向は exact な chart の随伴から
$R\operatorname{Hom}(\mathcal O(i),\mathcal P)\simeq M$ である。

有限例外列
$\mathcal E=\langle\mathcal O,\mathcal O(1),\mathcal O(2),\mathcal O(3)\rangle$
に沿って有限 mutation を行い、
\[
 F\longrightarrow\mathcal P\longrightarrow K,
 \qquad F\in\mathcal E,\quad
 K\in\mathcal K_A^{\mathrm{coh}}=\mathcal E^\perp
\]
を取る。線束から coherent 対象への $R\operatorname{Hom}$ は有限なので、
この三角は $D^b(\operatorname{qgr}A)$ の中で構成される。
式\eqref{cmo:eq:node-point-to-line}と有限 cone・直和因子に対する閉性により
$R\operatorname{Hom}(\mathcal P,F)$ は有界で有限次元である。
一方 $R\operatorname{Hom}(F,K)=0$ だから、射影三角より
\begin{equation}\label{cmo:eq:node-residual-self-ext}
 R\operatorname{Hom}(K,K)
 \simeq
 \operatorname{Cone}\bigl(
  R\operatorname{Hom}(\mathcal P,F)\longrightarrow
  R\operatorname{Hom}(\mathcal P,\mathcal P)\bigr).
\end{equation}
従って十分高い全ての次数で
\[
 \operatorname{Ext}^j(K,K)
 \simeq\operatorname{Ext}^j(\mathcal P,\mathcal P)
 \simeq\mathbb C^{32}.
\]
最後の次元は
$(1+s)^6/(1-s^2)=(1+s)^5/(1-s)$ の係数から従う。
このように全 coherent 残余圏自体が非 proper である。
また線束圏は proper なので、この $K$ は残余線束圏に属さない。
滑らかな可換 cubic の Ku は proper なので、両者は同値にならない。
この結論は、全 $D^b(\operatorname{qgr}A)$ の dg smoothness 単独を否定する主張ではない。

\paragraph{線束圏とその残余部分について.}
式\eqref{cmo:eq:morse-self-ext}から
$D^b(\operatorname{qgr}A)$ が非 proper であることは直ちに従う。
線束生成圏 $\mathcal C_A$ は常に proper であり、
もし smooth なら定理\paperref{main}{crs:thm:coherent}により全 coherent 圏へ飽和して矛盾する。
従って $\mathcal C_A$ は非 smooth である。
また命題\paperref{main}{crs:prop:residual-smooth}による有限例外 gluing から、
残余線束圏が smooth なら $\mathcal C_A$ も smooth になるので、
残余線束圏も非 smooth である。

\end{proof}

\begin{remark}[二種類の六十四係数]
定常 norm の六十四根は有理数であり、そこで生じる一次元定常点は
命題\ref{cmo:prop:point-perfect}により全て線束 perfect である。
これに対し $P_{64}$ の六十四根では、三十六次元単純加群が
非有界自己 Ext を持つ。この二集合は交わらない。
従って定常点の出現値だけでは非 smooth locus を捉えられないことが、
同じ固定混合族の中で具体的に示された。
定理\ref{cmo:thm:exact-sixty-four}で節点切断と射影境界を調べると、
非 smooth 係数は正確に $P_{64}$ の六十四根である。
系\ref{cmo:cor:fixed-zero-norm-smooth}により、
norm の六十四係数と $c=0$ は全て smooth proper である。
それらでの可換 Ku 比較は定理\ref{cmo:thm:all-central-commutative-kernel}と
系\ref{cmo:cor:norm-commutative-kernel}で構成され、主予想への反例にはならない。
\end{remark}

% !TEX root = ../../mixed-sklyanin-cubics.tex
\section{収縮族の相対滑らかさと余有限可換実現}
\label{mix:sec:relative-smoothness}

\subsection{分離表示への代数的収縮と有限線束模型}

本節の混合表示には、中心係数の無限遠を分離表示へ戻す代数的収縮がある。
この収縮は実際の代数の同型を与えるが、それだけで
smoothness や可換 Ku 実現の開性を結論しない。

\begin{proposition}[混合係数の収縮と中心係数の三乗被覆]
\label{cmo:prop:contraction}
$k\geq2$ とし、$\alpha^{(i)}\in\C^3$、$a_i\in\C^\times$ を固定する。
例\ref{cmo:ex:sixty-four}の表示と同じ
$\mathcal B_{\alpha^{(i)}}$ を用い、
定理\ref{cmo:thm:finite-center}の普遍中心元
$\Omega_{\alpha^{(i)}}$ を用いる。
$R=\C[t]$、$\deg t=0$ と置き、
\[
 \mathscr B=
   \bigotimes_{R[u],\,i=1}^{k}\mathcal B_{t^2\alpha^{(i)}},
 \qquad
 w_t=\sum_i a_i\Omega_{t^2\alpha^{(i)}}+u^3,
 \qquad \mathscr A=\mathscr B/(w_t)
\]
とする。ここでは $x_i,y_i,z_i,u$ は内部次数一である。
各 $\mathcal B_{t^2\alpha^{(i)}}$ の式は
\[
 [y_i,z_i]=x_i^2-t^2\alpha^{(i)}_1u^2
\]
とその巡回置換、および $u$ の中央性で与える。

この族は以下を満たす。
\begin{enumerate}[label=(\arabic*)]
\item $\mathscr B$ は $R[u]$ 上自由であり、$w_t$ は正則中心三次である。
      $\mathscr A$ の各内部次数は有限自由な $R$-加群であり、
      各ファイバーの Hilbert 級数は
      \[
        \frac{1-s^3}{(1-s)^{3k+1}}
      \]
      である。
\item $S=\C\langle x,y,z\rangle/
      ([y,z]-x^2,[z,x]-y^2,[x,y]-z^2)$ と書くと、
      $t=0$ は分離表示
      \[
       \mathscr A_0=
       \bigl(S^{\otimes k}\otimes\C[u]\bigr)/
          \left(\sum_i a_i\Omega_0+u^3\right)
      \]
      である。従って定理\paperref{main}{csc:thm:comparison}の下で、
      これは滑らかな可換三次 $(3k-1)$-fold の全 coherent Ku 成分を持つ。
\item $t_0\ne0$ では、元の中心変数を $v$ と書くことにより
      \begin{equation}\label{cmo:eq:contraction-isomorphism}
       \mathscr A_{t_0}\simeq
       \left(\bigotimes_{\C[v],\,i=1}^{k}
                           \mathcal B_{\alpha^{(i)}}\right)
       \bigg/\left(\sum_i a_i\Omega_{\alpha^{(i)}}+t_0^{-3}v^3\right).
      \end{equation}
      同型は全 active 変数を保ち、$v\mapsto t_0u$ とする次数付き代数の同型である。
\end{enumerate}
\end{proposition}

\begin{proof}
元の混合周囲代数を $\mathcal B^{\mathrm{orig}}$ とし、
その中心変数を $v$ と書く。中心環の準同型
\[
 \C[v]\longrightarrow\C[t,u],\qquad v\longmapsto tu
\]
による底変更が、表示した $\mathscr B$ そのものである。
例\ref{cmo:ex:sixty-four}で使った Braverman--Gaitsgory の判定は、
三次重なりの定数補正について $\beta\otimes1-1\otimes\beta=0$ を使うだけなので、
任意の固定した $\alpha^{(i)}$ に対して PBW 性を与える。
その Rees 代数は $\C[v]$ 上自由であり、$k$ 因子の tensor 積も同様である。
従って $\mathscr B$ は $R[u]$ 上自由である。
各 $\Omega_{t^2\alpha^{(i)}}$ は指定した中心元の底変更なので中心であり、
$u$ も中心である。

$u=0$ では $\mathscr B/(u)=S^{\otimes k}\otimes R$ で、
$w_t$ は $\sum_i a_i\Omega_0$ となる。
各 $S$ は $\C[\Omega_0]$ 上自由なので、この和は正則である。
従って $w_t f=0$ なら、まず $u=0$ へ移して $f\in u\mathscr B$ を得る。
$u$ の正則性により同じ議論を繰り返せる。
固定内部次数では $u$ で無限回割れないので $f=0$ である。
同じ証明は任意の $t=t_0$ のファイバーで成立する。

各内部次数における $w_t$ の乗法は有限自由 $R$-加群間の写像であり、
全ての閉ファイバーで単射である。
最大小行列式による階数判定から、その余核は局所自由である。
$R=\C[t]$ は PID なので余核は自由であり、
正則中心三次の二項完全列から Hilbert 級数を得る。

$t=0$ の式は定義から従う。
$S$ の点曲線は $E=V(x^3+y^3+z^3+xyz)$、
$\Omega_0$ はその THCR 商を与える標準中心三次の非零倍である。
分離表示の中心式は $u^3$ で、その Jacobian 環は有限次元であるため、
定理\paperref{main}{csc:thm:comparison}が適用できる。

最後に $t_0\ne0$ なら、$v=t_0u$ によって全二次関係と
$\Omega_{\alpha^{(i)}}$ がそれぞれ
$\mathcal B_{t_0^2\alpha^{(i)}}$ と $\Omega_{t_0^2\alpha^{(i)}}$ の式へ移る。
また $t_0^{-3}v^3=u^3$ である。
逆変換 $u=t_0^{-1}v$ も次数を保つため、
式\eqref{cmo:eq:contraction-isomorphism}は実際の代数同型となる。
\end{proof}

\begin{proposition}[収縮族の proper な有限線束模型]
\label{cmo:prop:proper-contraction}
上の族で $N=3k+1$、$n=N-2$、$a=N-3$ と置く。
$\mathscr B$ は左右 noetherian で、$R$ に相対的な長さ $N$ の
Koszul 双加群解消とその片側 augmentation 解消を持つ。
それらは任意の係数変更と両立し、相対 Gorenstein 型は
\[
 R\underline{\operatorname{Hom}}_{\mathscr B}(R,\mathscr B)
       \simeq R(N)[-N],\qquad
 R\underline{\operatorname{Hom}}_{\mathscr A}(R,\mathscr A)
       \simeq R(a)[-(N-1)]
\]
である。

$\mathcal O(j)=\pi\mathscr A(j)$ とすると、有限線束部分
$\mathscr C^{\mathrm{line}}$ とその残余部分は
\[
 E=\bigoplus_{j=0}^{N-1}\mathcal O(j),\qquad
 \mathscr C^{\mathrm{line}}=\operatorname{thick}_R(E),\qquad
 G=L_{\mathcal O}\cdots L_{\mathcal O(a-1)}(E)
\]
で生成される。
特に
\[
 \Gamma_R=R\operatorname{End}(G)\in\operatorname{Perf}(R),\qquad
 \mathscr K^{\mathrm{line}}\simeq\operatorname{Perf}(\Gamma_R)
\]
という proper な相対 dg 模型を持つ。
 $R$ の局所化および閉ファイバーへの通常の底変更は、
この線束模型と残余模型を対応する代数の模型へ移す。
相対 smoothness は、定理\ref{cmo:thm:relative-smoothness}で有限中心模型から導く。
\end{proposition}

\begin{proof}
各因子は $Q=R[u]$ 上の cyclic potential
\[
 \Phi_i=x_iy_iz_i-y_ix_iz_i
       -\frac{x_i^3+y_i^3+z_i^3}{3}
       +t^2u^2\left(\alpha^{(i)}_1x_i+
                     \alpha^{(i)}_2y_i+\alpha^{(i)}_3z_i\right)
\]
の三つの cyclic 微分で定まる。
active 濾過の随伴代数は $S\otimes Q$ であり、PBW 分裂は上の底変更と両立する。
定理\paperref{moduli}{thm:nfm-relative-ambient-cy}の証明中の、
中心環に相対的な長さ三の potential 解消を各因子で作る。
その $Q$ 上の tensor 積は長さ $3k$ の相対解消であり、
完全性は分裂した PBW 濾過の随伴複体で確認できる。
$C=\mathscr B\otimes_R\mathscr B^{\mathrm{op}}$、
$D=\mathscr B\otimes_Q\mathscr B^{\mathrm{op}}$ と置く。
共有する中心変数 $u$ の左右差 $z=u_L-u_R$ は $C$ 上正則で、$D=C/(z)$ である。
その Koszul 双対は
$R\underline{\operatorname{Hom}}_C(D,C)=D(1)[-1]$ を与える。
従って potential 解消の自己双対性と change of rings により
\[
 R\underline{\operatorname{Hom}}_C(\mathscr B,C)
 \simeq R\underline{\operatorname{Hom}}_D
              (\mathscr B,D(1)[-1])
 \simeq\mathscr B(N)[-N]
\]
を得る。ここでは中心差の因子を一度だけ使っており、
異なるブロックで共有する $u$ を重複して数えていない。
また $D$ は $C$-perfect、$\mathscr B$ は $D$-perfect なので、
$\mathscr B$ は $C$-perfect である。

二次 Koszul 模型はこの change of rings とは別に指定できる。
各体値で active 濾過の随伴代数は
$S^{\otimes k}\otimes\C[u]$ であり、有限な内部次数の bar spectral sequence により
その Koszul 性が元の代数へ移る。
Hilbert 級数 $(1-s)^{-N}$ と合わせると、Koszul 空間の階数は
$(1+s)^{3k}(1+s)=(1+s)^N$ の係数である。
被約基底 $R$ 上の一定階数判定によりこれらは有限射影加群で、
標準 Koszul 双加群複体は係数変更と両立して完全となる。
augmentation を施せば最初の Gorenstein 式となり、
正則中心三次の二項解消から第二式を得る。

$S^{\otimes k}$ の noetherian 性は補題\paperref{main}{mbnp:lem:noetherian}による。
$\mathscr B/(u)=S^{\otimes k}\otimes R$ は中心多項式拡大なので noetherian である。
正次数中心元 $u$ による初項の持上げは、固定内部次数では有限回で止まるので、
$\mathscr B$ も左右 noetherian となる。$\mathscr A$ もその商として同様である。

$R=\C[t]$ の大域次元は一であり、上の有限双加群解消を片側加群へ tensor すると
$\operatorname{gldim}\mathscr B\leq N+1$ を得る。
正則中心元の change of rings により、$\mathscr A$ の左右 injective dimension も有限である。
従って noetherian 性と相対 Gorenstein 式を合わせると、
Burke--Stevenson の $\chi$ と有限 cohomological dimension の入力も満たされる。
以上の有限解消と Gorenstein 型を用いると、
補題\paperref{main}{lp:lem:local-cohomology}と式\papereqref{main}{lp:eq:H}の証明が適用できる。
実際、ここで使う入力は上の noetherian 性・有限次元性・相対 augmentation 双対と次数別自由性であり、
特殊ファイバーの可換性ではない。
従って
\[
 \operatorname{Ext}^p(\mathcal O(i),\mathcal O(j))=
 \begin{cases}
  \mathscr A_{j-i},&p=0,\\
  (\mathscr A_{i-j-a})^\vee,&p=n,\\
  0,&\text{その他}
 \end{cases}
\]
となる。右辺の各次数は有限自由 $R$-加群である。
$R$ の ambient augmentation Koszul 解消を $\mathscr A$ へ導来 tensor する。
中心三次の二項解消から $\operatorname{Tor}^{\mathscr B}_i(R,\mathscr A)$ は
$i=0,1$ にのみあり、それぞれ $R,R(-3)$ である。
これらは内部次数について torsion なので、得られる有限複体は tails 内で完全となる。
両端は一つの線束で、中間項の shift は $1,\ldots,N-1$ にあるため、
任意の連続した $N$ 個の twist からその前後を生成できる。
これで $E$ の生成性と properness を得る。
命題\paperref{main}{lp:prop:residual}と同じ有限 mutation によって $G$ を作ると、
その自己 Hom は $R$-perfect である。
底変更についても命題\paperref{main}{lp:prop:basechange}の
次数別に安定する局所 cohomology と有限 Hom の議論を使う。
指定した局所化と閉ファイバーでは、$\mathscr B,\mathscr A$ の noetherian 性は
中央局所化または商から従う。底も有限大域次元の正則環であり、
有限解消と次数別自由性は保たれるので、その適用条件を満たす。
\end{proof}

\subsection{有限中心上の通常の加群に対する局所持上げ}

\begin{lemma}[中心的な一変数に沿う射影次元の持上げ]
\label{cmo:lem:projective-lift}
$(C,\mathfrak m)$ を noetherian な可換局所環、
$t\in\mathfrak m$ とする。$B$ は中央 $C$-代数で $C$ 上有限、
$M$ は有限 $B$-加群、$d\ge0$ とする。
$t$ は $B$ と $M$ 上正則であると仮定する。
このとき
\[
 \operatorname{pd}_{B/tB}(M/tM)\le d
 \quad\Longrightarrow\quad
 \operatorname{pd}_B M\le d.
\]
\end{lemma}
\begin{proof}
$B$ は左右 noetherian である。$M$ の有限階数自由分解の最初の
$d$ 段を取り、その第 $d$ syzygy を $K$ とする。
全ての syzygy は $t$-torsionfree であり、$t$ で割っても
分解の完全性が保たれる。従って $K/tK$ は $B/tB$ 上有限射影である。

$C$ の $(t)$-adic 完備化 $\widehat C$ は忠実平坦である
（\cite[Tag~00MC]{CMOStacksCompletion}）。
$\widehat B=B\otimes_C\widehat C$ と
$\widehat K=K\otimes_C\widehat C$ を取る。
有限性によりこれらは $(t)$-adically complete で、
$t$ は $\widehat B$ の Jacobson radical に属する。
$K/tK$ を表示する有限行列の idempotent は $\widehat B$ へ持ち上がる。
この非可換行列に対する事実は、$e^2-e\in(t)$ から出発し、
\[
 e\longmapsto e-(2e-1)^{-1}(e^2-e)
\]
で誤差を二乗化して極限を取れば直接分かる。
逆元は $(2e-1)^2=1+4(e^2-e)$ から存在する。
従って $K/tK$ を持ち上げる有限射影 $\widehat B$-加群 $Q$ がある。

$Q$ の射影性により、指定した $Q/tQ\simeq\widehat K/t\widehat K$ を
射 $Q\to\widehat K$ へ持ち上げる。
余核は $t$ で割ると零なので Nakayama によりこの射は全射である。
核 $L$ は有限であり、$\widehat K$ 上の $t$ の正則性から
$L/tL=0$ を得る。再び Nakayama により $L=0$、
従って $\widehat K$ は射影である。

有限射影性はこの中央の忠実平坦底変更から降下する。
具体的には、有限自由表示 $0\to J\to B^r\to K\to0$ の分裂障害を考える。
$K$ の有限階数自由分解と平坦性により
\[
 \operatorname{Ext}^1_B(K,J)\otimes_C\widehat C
 \simeq
 \operatorname{Ext}^1_{\widehat B}(\widehat K,\widehat J)=0.
\]
忠実平坦性から障害は零で、$K$ は有限射影である。
\end{proof}

\begin{lemma}[中央スキーム上の perfect locus]
\label{cmo:lem:perfect-locus}
$Y$ を noetherian スキーム、$\mathcal B$ を coherent な
$\mathcal O_Y$-代数、$M$ を coherent $\mathcal B$-加群とする。
$M_y$ が $\mathcal B_y$ 上有限射影次元を持つ点の集合は開である。
この開集合上では $M$ は局所的に有限生成射影分解を持つ。
\end{lemma}
\begin{proof}
アファイン近傍上で $M$ の有限自由分解を必要な段階まで取り、
射影になった syzygy を $K$ とする。
$K_y$ の有限自由加群からの射影表示とその splitting は有限個の元で記述される。
$K$ の有限表示性により、射と splitting の等式の分母を有限個払えば
$y$ の中央開近傍上で $K$ 自身が射影となる。
\end{proof}

\begin{proposition}[proper な中央台からの相対 perfect 性]
\label{cmo:prop:proper-perfect-spread}
$R=\mathbf C[t]$、$Y\to\operatorname{Spec}R$ を proper な
noetherian スキームとする。
$\mathcal B$ を coherent $\mathcal O_Y$-代数、$M$ を coherent
$\mathcal B$-加群とし、両者が $R$-flat であると仮定する。
$M_0$ が $\mathcal B_0$ 上 perfect なら、
$0$ を含む Zariski 開近傍 $U\subset\operatorname{Spec}R$ が存在し、
$M_U$ は $\mathcal B_U$ 上 perfect となる。
\end{proposition}
\begin{proof}
$Y_0$ の各点で前二補題を適用する。
非 perfect locus は閉集合で $Y_0$ と交わらない。
proper な射によるその像も閉であり $0$ を含まないから、
この像を除けばよい。
\end{proof}


\subsection{有限中心模型による収縮族の相対 smoothness}

\begin{theorem}[分離ファイバーを含む相対 smooth proper 近傍]
\label{cmo:thm:relative-smoothness}
命題\ref{cmo:prop:contraction}の任意の
$\alpha^{(i)}\in\C^3$、$a_i\in\C^\times$、$k\ge2$ を取る。
$0$ を含む Zariski 開近傍 $U\subset\mathbb A^1_t$ 上で、
命題\ref{cmo:prop:proper-contraction}の線束模型と残余模型は
相対 smooth proper である。特に
\begin{equation}\label{cmo:eq:relative-diagonal}
 \Gamma_U\in\operatorname{Perf}\!\left(
       \Gamma_U^{\mathrm{op}}\otimes_U^{\mathbf L}\Gamma_U\right).
\end{equation}
ここで追加の対角 perfect 性の仮定は課さない。
\end{theorem}
\begin{proof}
系\ref{cmo:cor:relative-finite-center}によって、$\mathscr A$ は
$u$、各 $\Omega_{t^2\alpha^{(i)}}$、各六次中心元の像が生成する
可換有限型次数付き中心部分環 $C_R$ 上有限である。
その生成元の内部次数は $1,3,6$ である。
$C_R\subset\mathscr A$ は $R=\C[t]$ 上 torsionfree であり、
$R$ が PID なので $R$-flat である。
六次 Veronese $C_R^{(6)}$ は次数一で生成される。
実際、重み $1,3,6$ の単項式で重みが六の倍数なら、重み六の単項式へ分割できる。

\[
 Z=\operatorname{Proj}C_R^{(6)},\qquad
 \mathcal A=\widetilde{\mathscr A^{(6)}}
\]
と置く。$Z\to\operatorname{Spec}R$ は projective かつ flat であり、
$\mathcal A$ は coherent な $\mathcal O_Z$-代数で $R$-flat である。
flatness は、中央元を反転した次数零部分が $R$-torsionfree であることから従う。
Veronese の tails 同値と、可換標準次数環上の Serre 同値を
有限 $C_R^{(6)}$-代数の加群作用とともに用いれば
\[
 \operatorname{QGr}\mathscr A\simeq\operatorname{Qcoh}(Z,\mathcal A),\qquad
 \operatorname{qgr}\mathscr A\simeq\operatorname{coh}(Z,\mathcal A)
\]
を得る。ここでは $\mathscr A$ が内部次数一で生成されることを使う。
$C_R$ のファイバーが全中心と一致するという仮定は不要である。

全ての元の twist $\mathcal O(j)$ が $\mathcal A$-perfect であることを確認する。
$Z$ を六次の中央元 $f$ の標準開集合で覆い、
$L=\mathscr A[f^{-1}]$ とする。
$f$ を長さ六の語の和として表示し、先頭または末尾の一文字を分けると
\[
 1\in L_1L_{-1},\qquad 1\in L_{-1}L_1
\]
となる。これを反復すれば $L$ は強 $\mathbb Z$-次数付きである。
従って $L_j$ は $L_0$ 上有限生成射影で、各 twist は局所射影である。
中央の ample twist $\mathcal A\otimes\mathcal O_Z(m)$ は
元の $\mathcal O(6m)$ に対応する。
これらは $D_{\mathrm{qc}}(Z,\mathcal A)$ を compact に生成する。
具体的には、忘却関手との随伴と ample twist の検出性
\cite[Tag~0BQR]{CMOStacksDerived}を用いる。
生成性は ample な切断によるアファイン局所化、compact 性は有限アファイン被覆の
\v Cech 複体と局所射影性から従う。
命題\ref{cmo:prop:proper-contraction}の有限線束和 $E$ は全 twist を厚生成するので、
$E$ はこの非有界導来圏の compact generator である。
この記述は閉ファイバーと $R$ の局所化に両立する。

$Y=Z\times_R Z$ 上で
\[
 \Lambda=\mathcal A^{\mathrm{op}}\boxtimes_R\mathcal A,
 \qquad M=\Delta_*\mathcal A,
 \qquad F=E^\vee\boxtimes_R E,
 \qquad E^\vee=R\mathcal Hom_{\mathcal A}(E,\mathcal A)
\]
と置く。$\Lambda,M$ は coherent かつ $R$-flat であり、
$M_0=\Delta_{0*}\mathcal A_0$ である。
双対は左右の perfect 圏の反同値なので、$E^\vee$ は反対側を生成する。
$F$ の厚生成部分は全ての中央 twist $\Lambda(i,j)$ を含み、
特に ample な $\Lambda(m,m)$ を含む。
従って $F$ は $D_{\mathrm{qc}}(Y,\Lambda)$ の compact generator である。

$\Sigma_R=R\operatorname{End}_{\mathcal A}(E)$ と書く。
局所有限射影加群の Hom は $R$-flat であり、有限アファイン被覆の
\v Cech 模型を tensor すると K\"unneth 同型を得る
\cite[Tag~0FU4]{CMOStacksDerived}。
局所的な有限射影加群の評価同型と合わせて
\begin{equation}\label{cmo:eq:kernel-diagonal}
 \begin{split}
 R\operatorname{End}_\Lambda(F)
      &\simeq\Sigma_R^{\mathrm{op}}\otimes_R^{\mathbf L}\Sigma_R,\\
 R\operatorname{Hom}_\Lambda(F,M)&\simeq\Sigma_R
 \end{split}
\end{equation}
となる。第二式は局所的に $T\mapsto[\varphi\otimes e\mapsto\varphi(T(e))]$ で与えられる。
$S^\vee\boxtimes U$ による前合成は $T\mapsto STU$ に対応するため、
左右の作用を保持し、$\Sigma_R$ の対角双加群そのものを与える。
compact generator による dg Morita 同値
\cite[Theorem~3.8(c)]{NFMKellerDg}から、
$M$ の $\Lambda$-perfect 性は線束圏の相対 smoothness と同値である。
同じ説明が $t=0$ でも成立する。

特殊ファイバーは分離表示であり、定理\paperref{main}{csc:thm:comparison}によって
全 coherent 圏が smooth proper である。
その有限線束和は生成対象なので、式\eqref{cmo:eq:kernel-diagonal}の特殊化により
$M_0$ は $\Lambda_0$-perfect である。
$Y\to\operatorname{Spec}R$ は proper なので、
命題\ref{cmo:prop:proper-perfect-spread}を $\Lambda,M$ に適用できる。
従って $0$ を含む開近傍 $U$ で $M_U$ は perfect であり、
線束圏は $U$ 上 smooth となる。
proper 性は既に命題\ref{cmo:prop:proper-contraction}で確認した。
最後に、有限例外列による admissible な残余成分には smooth proper 性が移るため、
$\Gamma_U$ について式\eqref{cmo:eq:relative-diagonal}を得る。
\end{proof}

\subsection{一般定数型混合係数に対する余有限可換実現}

\begin{corollary}[定数型混合族の余有限可換実現]
\label{cmo:cor:cofinite}
任意の $k\ge2$、$\alpha^{(i)}\in\C^3$、$a_i\in\C^\times$ を固定し、
命題\ref{cmo:prop:contraction}の定数型混合族を考える。
このとき元の $c$-族について有限集合 $F\subset\C$ が存在し、
$c\notin F$ では $D^b(\operatorname{qgr}A_c)$ は smooth proper であり、
ある滑らかな可換三次 $(3k-1)$-fold $X_c$ と
\[
 \operatorname{Ku}_{\mathrm{coh}}(A_c)
             \simeq\operatorname{Ku}(X_c)
\]
が存在する。
これは通常の dg Morita 同値であり、rotation・周期・point scheme の
標識の保存を追加しない。
\end{corollary}

\begin{proof}
定理\ref{cmo:thm:relative-smoothness}によって $\operatorname{Perf}(\Gamma_U)$ は
実際の相対 smooth proper な圏族である。
$n=3k-1\geq5$ なので $n\ne4$ であり、
$t=0$ の可換実現に定理\paperref{main}{grl:thm:open-locus}を適用できる。
従って $0$ を含む開近傍 $U'\subset U$ 上で
各残余線束圏は可換 Ku 成分を持つ。
各閉ファイバーの周囲代数は、上の有限 Koszul 双加群解消により大域次元が
高々 $N$、augmentation 双対の非零な最高 Ext により丁度 $N$ であり、
Gorenstein parameter $N$ の両側 noetherian Koszul AS-regular 代数である。
従って中心三次商に対する命題\paperref{main}{crs:prop:residual-smooth}と
定理\paperref{main}{crs:thm:coherent}の仮定を満たす。
残余圏と有限例外列の proper な gluing から線束部分の smoothness を得て、
その後の coherent 飽和により線束部分と全 coherent 圏が一致する。

$\mathbb A^1\setminus U'$ は有限集合である。
$t\mapsto t^{-3}$ は $\mathbb G_m$ の有限全射であり、
式\eqref{cmo:eq:contraction-isomorphism}により
$U'\cap\mathbb G_m$ の像に属する全 $c$ で結論を得る。
例えば $F$ は $0$ と $\mathbb A^1\setminus U'$ の非零点の像を含む
有限集合として取れる。
\end{proof}

\begin{remark}[有限中心性の役割と残る個別判定]
\label{cmo:rem:cofinite-scope}
この証明は、明示的な六次中心元による有限中心性、通常の coherent 対角加群、
proper な中央台を使う。一般の proper dg 代数について、smooth な特殊ファイバーや
全 Artin 切断での互換同値だけから相対 smoothness を推論しているのではない。
その推論への反例と、通常の完備底で生成性が失われる機構は研究ノートに記す。

可換実現を保証するための有限除外集合 $F$ はまだ有効に決定していない。
一方、例\ref{cmo:ex:sixty-four}の指定二因子については、
定理\ref{cmo:thm:exact-sixty-four}が非 smooth 係数を $V(P_{64})$ と完全に決定した。
norm の全六十四根と $c=0$ は smooth proper locus に属する。
この結果と余有限可換実現を合わせても、各滑らかな係数での比較核は自動的には得られない。
norm は定常点の出現を検出し、$P_{64}$ は指定族の smoothness を判定する。
これらを区別した上で、定理\ref{cmo:thm:all-central-commutative-kernel}は
指定二ブロック族の全 smooth locus に実際の比較核を構成する。
一般のブロック数に対する上の収縮・開性の証明の適用範囲は、その直接構成とは別である。
また得られた可換標的について、point scheme と Poisson 零スキームの
同時比較や、元のブロック分解の保存はまだ主張していない。
\end{remark}

% !TEX root = ../../mixed-sklyanin-cubics.tex
\section{指定混合族の smooth locus と Hochschild homology}
\label{mix:sec:smooth-locus}

\subsection{中央切断から混合中心和の対角へ}
\label{cmo:subsec:slice-mixed-diagonal}

\begin{lemma}[形式的中央ファイバーから相対対角への持上げ]
\label{cmo:lem:relative-diagonal-from-slice}
$A$ を有限型可換中心上有限な noetherian $\C$ 代数、
$s\in Z(A)$ とし、$A$ は $C=\C[s]$ 上平坦とする。
$T=A^{\mathrm{op}}\otimes_C A$ と置く。
中心の閉点 $z$、$s(z)=\lambda$ において、
\[
 \widehat A_z/(s-\lambda)
\]
の形式的対角が、完備包絡環上で長さ $d$ の有限生成射影分解を持つと仮定する。
このとき $A$ の相対対角は、$T$ の中央点 $(z,z)$ の近傍で perfect であり、
射影次元は高々 $d$ である。
\end{lemma}

\begin{proof}
$T$ の可換中心部分環を $(z,z)$ で局所化して完備化し、
得られる代数を $\widehat T$、$T$ 加群 $A$ の完備化を $\widehat M$ と書く。
$T$ が有限な可換中心上有限であるため、この完備化は中央の忠実平坦底変更である。
$t=s-\lambda$ は $\widehat T$ と $\widehat M$ に正則に作用し、
$t\in\operatorname{Jac}(\widehat T)$ である。
さらに $\widehat T/(t)$ は切断代数の完備包絡環であり、
$\widehat M/t\widehat M$ はその形式的対角である。

$\widehat M$ の有限生成自由分解を最初の $d$ 段まで取り、最後の syzygy を $K$ とする。
$t$ の正則性から $\operatorname{Tor}^{\widehat T}_{>0}(\widehat T/(t),\widehat M)=0$ であり、
この部分分解を $t$ で割っても完全である。
仮定により $K/tK$ は射影である。
完備環上でその冪等行列を持ち上げ、有限生成射影加群 $P$ と
$P/tP\simeq K/tK$ を取る。
$P$ の射影性によりこの同型を $P\to K$ に持ち上げる。
Nakayama により全射となり、$t$ の $K$-正則性によりその核を $t$ で割ると零になる。
再び Nakayama から核も零なので、$K$ は射影である。
従って $\widehat M$ は射影次元高々 $d$ である。

有限表示加群の射影性は中央忠実平坦底変更から降下するため、
同じ有限自由部分分解の syzygy は完備化前でも局所的に射影である。
有限表示性によりこの結論は $(z,z)$ の Zariski 近傍へ広がる。
この証明は、非 Hensel 環 $\C[s]_{(s-\lambda)}$ 上での冪等元の持上げを仮定しない。
\end{proof}

\begin{lemma}[一方の相対対角が perfect な混合中心和]
\label{cmo:lem:mixed-relative-diagonal}
$C=\C[s]$、$A,B$ を中央の $C$-作用を持つ $C$-平坦代数とし、
\[
 T=A^{\mathrm{op}}\otimes_C A,\qquad R=A\otimes_C B
\]
と置く。
$A$ の相対対角が $T$ 上 perfect、$B$ の対角が
$B^e=B^{\mathrm{op}}\otimes_\C B$ 上 perfect なら、
$R$ の対角は $R^e$ 上 perfect である。
それぞれの分解の長さが $d,e$ なら、得られる長さは高々 $d+e$ である。
この主張は有限可換中心上の中央点の近傍へ局所化しても成立する。
\end{lemma}

\begin{proof}
包絡環の左作用を
$(a^{\mathrm{op}}\otimes b)\cdot x=bxa$ に統一する。
$(R^e,T)$ 双加群
\[
 N=R^e\otimes_{B^e}B
\]
を考える。$A$ の左右二つの中央パラメータは $B$ の対角作用によって一致するので、
\[
 N\simeq T\otimes_C B
\]
である。右 $T$-作用は $T$ 因子の右乗法で与えられ、
左 $R^e$-作用と可換である。
ここでは $B^e\to B$ を環準同型として使っていない。

また
\[
 R^e\simeq(A^{\mathrm{op}}\otimes_\C A)
                 \otimes_{C\otimes_\C C}B^e
\]
であり、$A$ の $C$-平坦性から $R^e$ は右 $B^e$-平坦である。
従って $B$ の有限射影双加群分解を誘導すると、
$N$ は $R^e$ 上 perfect で、その射影次元は高々 $e$ となる。
一方、$B$ の $C$-平坦性から $N=T\otimes_C B$ は右 $T$-平坦である。
$A$ の有限 $T$-射影分解 $P_\bullet\to A$ をテンソルすると、
\[
 N\otimes_TP_\bullet\longrightarrow
 N\otimes_TA\simeq R
\]
は完全である。各項は有限個の $N$ の和の直和因子なので、
$R^e$ 上 perfect であり、全複体から長さ $d+e$ の上限を得る。

中央局所化後も同じテンソル積の表示と平坦性が成立する。
相対包絡環の中央基底を $Z(A)\otimes_C Z(A)$ に選ぶ。
その近傍を定める元は $Z(A)\otimes_\C Z(A)$ へ持ち上がるので、
絶対包絡環の中央元として $R^e$ へ写す。
その局所化上では $N$ の左右の作用で当該中央元が一致するため、
局所的な有限 $T$-射影解消にも同じ議論を適用できる。
相対対角の情報を完備局所模型から得る場合は、前補題により先に
代数的な相対包絡環の近傍へ降下し、その後でこの議論を適用する。
従って形式冪級数環の通常の包絡環に関する smoothness を仮定する必要はない。
\end{proof}

\begin{corollary}[節点を含む混合中心和への適用]
\label{cmo:cor:node-mixed-smooth}
混合因子の完備局所代数が
$\operatorname{End}_R(R^{\oplus r}\oplus I^{\oplus s})$、$r,s>0$ と Morita 同値である点を考える。
その点で $d\Omega_1$ が通常二重点の双対 Hessian に関して非等方なら、
他方の三次元 Calabi--Yau 因子と作る中心和
\[
 (D_1\otimes_\C D_2)/(\Omega_1+\beta\Omega_2+c),\qquad \beta\ne0,
\]
の対角は、その中央点の近傍で perfect であり、長さは高々五である。
\end{corollary}
\begin{proof}
$C=\C[s]$ は第一因子に $s\mapsto\Omega_1$、
第二因子に $s\mapsto-c-\beta\Omega_2$ と作用させる。
各因子は domain で、指定中心元 $\Omega_i$ は $\C$ 上超越的なので、$C$ 上 torsion-free、従って平坦である。
命題\ref{cmo:prop:conifold-nonisotropic-slice}と
補題\ref{cmo:lem:relative-diagonal-from-slice}が第一因子の相対長さ二を与える。
第二因子の既存の三次元 Calabi--Yau 双加群分解は長さ三である。
補題\ref{cmo:lem:mixed-relative-diagonal}から結論を得る。
この系の標準 conifold への Morita 同値は独立の入力であり、
中心が通常二重点であることだけからは仮定しない。
\end{proof}

\subsection{射影境界の対角 perfect 性}

アファイン chart $u\ne0$ の判定を全射影圏へ移すため、
ここでは $u=0$ の境界を任意の $c$ で扱う。
内部次数一の $u$ を、六次 Veronese の chart 上でそのまま
次数零の中心元と扱うことはできない。その点を有限巡回次数付けで解消する。

\begin{lemma}[中心一次元に沿う smoothness の局所持上げ]
\label{cmo:lem:central-smooth-lift}
$Q$ を有限生成可換 $\mathbb C$-代数上有限な代数、
$u\in Z(Q)$ を正則元とする。
$\overline Q=Q/uQ$ の対角が、指定した閉中心点の近傍で perfect であるとする。
このとき $u=0$ 上の対応する点の近傍で、$Q$ の対角も perfect である。
\end{lemma}
\begin{proof}
底 $C=\mathbb C[u]$ を $(u)$ で局所化してよい。
$u$ の正則性から $Q$ はこの離散付値環上 torsionfree、従って flat である。
相対包絡環と対角加群を
\[
 T=Q^{\mathrm{op}}\otimes_C Q,\qquad M=Q
\]
と置く。両者は $C$-flat であり、
$T/uT=\overline Q^{\mathrm{op}}\otimes_{\mathbb C}\overline Q$、
$M/uM=\overline Q$ である。
$T$ は有限型可換中心部分環上有限である。
補題\ref{cmo:lem:projective-lift}をその中心点で適用し、
同補題の忠実平坦完備化と射影性の降下を用い、
補題\ref{cmo:lem:perfect-locus}で近傍へ広げると、$M$ は局所的に $T$-perfect となる。

$Q^e=Q^{\mathrm{op}}\otimes_{\mathbb C}Q$ 上の中心元
$\delta=u\otimes1-1\otimes u$ は正則である。
実際、$Q^e$ は $C\otimes_{\mathbb C}C$ 上 flat であり、
そこで二つの座標の差は正則である。
従って $T=Q^e/(\delta)$ は $Q^e$ 上長さ一の有限自由解消を持つ。
$M$ の有限 $T$-射影解消をこの解消と合成すれば、
$M$ は $Q^e$-perfect となる。
局所化後の結論は有限表示の分母を払って元の中央近傍へ戻る。
\end{proof}

\begin{proposition}[定数型混合族の全係数における射影境界]
\label{cmo:prop:projective-boundary-smooth}
任意の $k\ge2$、$\alpha^{(i)}\in\mathbb C^3$、
$a_i\in\mathbb C^\times$ と $c\in\mathbb C$ に対し、
\[
 A_c=\left(\bigotimes_{\mathbb C[u]}B_{\alpha^{(i)}}\right)
       \Big/\left(\sum_i a_i\Omega_i+cu^3\right)
\]
の有限中心射影模型の対角は、$u=0$ の近傍で perfect である。
従ってこの境界から非 smooth 係数は生じない。
\end{proposition}
\begin{proof}
既存の正則列と有限中心性により $u$ は $A_c$ 上正則であり、
\[
 \overline A=A_c/uA_c
   =S^{\otimes k}/\left(\sum_i a_i\Omega_0^{(i)}\right)
\]
は $\alpha^{(i)},c$ に依存しない。
定理\paperref{main}{pps:thm:elliptic-existence}と、その入力である
系\paperref{main}{mbnp:cor:coherent}により、この分離表示の全 coherent 圏は
smooth proper である。この直接比較では次元四を除外しない。
有限中心模型の対角と compact generator の関係は
式\eqref{cmo:eq:kernel-diagonal}と同じなので、
$\overline A$ の各中央射影 chart の対角は perfect である。

$A_c$ の次数 $1,3,6$ の既知の中心元が生成する可換環を $C_A$ とする。
$\operatorname{Proj}C_A^{(6)}$ を六次中心元 $f$ の標準開集合で覆い、
\[
 L=A_c[f^{-1}],\qquad Q=L/(f-1)
\]
と置く。$A_c$ は次数一で生成されるので、
定理\ref{cmo:thm:relative-smoothness}の証明と同じく $L$ は強 $\mathbb Z$-次数付きである。
$Q$ は強 $\mathbb Z/6$-次数付きで、
その零成分は $Q_{\bar0}\simeq L_0$ である。
具体的には $L_{6j}$ の元を $f^{-j}$ 倍して $L_0$ と同定する。
各巡回次数成分は $L_0$ 上可逆な双加群であり、
$Q$ は左右とも有限生成射影な generator である。

双対群 $\mu_6$ が次数に従って $Q$ に作用する。
有限群次数付けと smash product の End 表示は既知の一般論である
\cite[Theorems~1.3, 2.1]{CMONastasescuRodino}。
本例では強次数付けから
$\operatorname{Hom}_{L_0}(Q_{\bar s},Q_{\bar r})\simeq Q_{\overline{r-s}}$
を左乗法で得る。各次数成分への射影を群代数の冪等元で表せば、自然な作用写像は
\[
 Q\rtimes\mu_6\ \simeq\ \operatorname{End}_{L_0}(Q_{L_0})
\]
となる。従って $Q\rtimes\mu_6$ と $L_0$ は Morita 同値である。
位数六が可逆なので、この事実から
$Q$ と $L_0$ の対角 perfect 性は同値となる。
実際、有限群の skew group への移行では平均化による分裂を用い、
逆向きには $(Q\rtimes\mu_6)$ が $Q$ 上有限自由であり、
その対角の制限が $Q$ を直和因子に持つことを使う。
同じ議論は中央開集合上で成立する。

$u$ は $L$ 上正則である。
$L/uL$ 上でも $f$ は可逆な六次の元なので、有限個の次数成分しか持たない元の
最高次数を見ると $f-1$ は正則である。
ここで $ua=(f-1)b$ なら、$u$ で割ってから $f-1$ の正則性を使うと
$b=ud$ と書ける。$L$ 上で $u$ を消去すれば $a=(f-1)d$ である。
従って $u$ は $Q$ 上正則となる。
また $\overline Q=Q/uQ$ は $\overline A[f^{-1}]$ に対する同じ巡回次数付き模型である。
先ほどの Morita 同値と平均化を $\overline Q$ に適用すると、
分離境界の対角 perfect 性から $\overline Q$ の対角 perfect 性を得る。
補題\ref{cmo:lem:central-smooth-lift}により、$Q$ は $u=0$ の近傍で smooth である。

必要なら有限群 $\mu_6$ による有限個の近傍の共通部分を取り、
近傍を群作用で不変にする。$Q$ は中央 chart の可換環
$(C_A[f^{-1}])_0$ 上有限であり、その境界 $u^6/f=0$ の上では
全ての閉点で $u=0$ である。
従って有限群の skew group と Morita 同値を通して、
$L_0$ の対角は境界の近傍で perfect となる。
これらの chart は射影境界を覆うので結論を得る。
\end{proof}

\subsection{指定二因子の smooth locus の完全決定}

\begin{theorem}[非 smooth 係数の正確な六十四点集合]
\label{cmo:thm:exact-sixty-four}
例\ref{cmo:ex:sixty-four}の
$\alpha^{(1)}=(1,4,9)$、$\alpha^{(2)}=(16,25,49)$ と
$a_1=1,a_2=100$ を固定する。
式\eqref{cmo:eq:morse-resultant}の $P_{64}$ に対し、
次の条件は同値である。
\begin{enumerate}
 \item $P_{64}(c)\ne0$。
 \item アファイン中心 chart
 \[
  R_c=(D_{\alpha^{(1)}}\otimes_{\mathbb C}D_{\alpha^{(2)}})
           /(\Omega_1+100\Omega_2+c)
 \]
 の対角は perfect である。
 \item $A_c$ の線束生成圏は smooth である。
 \item $D^b(\operatorname{qgr}A_c)$ は smooth proper である。
 \item $\operatorname{Ku}_{\mathrm{coh}}(A_c)$ は smooth proper である。
\end{enumerate}
さらに $P_{64}(c)\ne0$ では $\operatorname{qgr}A_c$ の
homological dimension は五であり、$P_{64}(c)=0$ では無限である。
従って中心の非 Azumaya 点と射影境界を含めても、
定理\ref{cmo:thm:morse-singular-sixty-four}の六十四点以外に
非 smooth 係数は存在しない。
\end{theorem}
\begin{proof}
$X_i=\operatorname{Spec}Z(D_{\alpha^{(i)}})$ とする。
$R_c$ は可換有限型環
\[
 C_c=(Z(D_{\alpha^{(1)}})\otimes_{\mathbb C}Z(D_{\alpha^{(2)}}))
          /(s_1+100s_2+c)
\]
の中央作用に関して有限である。その閉点は条件 $s_1(z_1)+100s_2(z_2)+c=0$ を満たす
$(z_1,z_2)\in X_1\times X_2$ として表される。
中心をさらに拡大した場合にも、この有限中央模型で対角の perfect 性を判定できる。

各因子について、正則中心上の非臨界点では
中心写像 $s_i$ が smooth であり、Azumaya 性から相対対角は perfect である。
十四節点では、定理\ref{cmo:thm:fourteen-node-orders}と
補題\ref{cmo:lem:all-nodes-nonisotropic}、
命題\ref{cmo:prop:conifold-nonisotropic-slice}により、切断代数は形式的に
$\mathbb C[[p,q]]\rtimes\mu_2$ と Morita 同値である。
その完備対角の有限解消を代数的な中央局所環へ降下し、
中心変数に沿う射影次元持上げを用いると、そこでも相対対角は perfect である。
この二つの操作は補題\ref{cmo:lem:relative-diagonal-from-slice}による。
従って各因子で相対対角が perfect でない可能性がある点は、
$Q_i(s_i^2)=0$ に対応する八個の Morse 点だけである。

$P_{64}(c)\ne0$ とする。
終結式の定義から、任意の $(z_1,z_2)\in\operatorname{Spec}C_c$ で
少なくとも一方の $z_i$ はこの八個の集合に属さない。
その因子の相対対角 perfect 性と、他因子の三次元 Calabi--Yau 双加群解消を
補題\ref{cmo:lem:mixed-relative-diagonal}で合成する。
すると $R_c$ の対角は当該閉中心点で perfect である。
対角の台の外では対角加群自体が零である。
有限型中央スキームの全閉点をこのように覆い、perfect locus の開性と
準コンパクト性を使えば、$R_c$ の対角は大域的に perfect となる。
これで (1)$\Rightarrow$(2) を得る。

命題\ref{cmo:prop:projective-boundary-smooth}により、
$u=0$ の射影境界の対角は任意の $c$ で perfect である。
従って (2) から有限中心射影模型全体の対角 perfect 性を得る。
式\eqref{cmo:eq:kernel-diagonal}の compact generator による Morita 比較で
線束生成圏は smooth となる。その proper 性は既存の
命題\ref{cmo:prop:proper-contraction}と同じ次数別 Hom の議論、
または補題\ref{cmo:lem:mixed-line-serre}の構成による。
定理\paperref{main}{crs:thm:coherent}の coherent 飽和を適用すると、
この線束圏は全 $D^b(\operatorname{qgr}A_c)$ に等しい。
有限例外列を除く残余圏にも smooth proper 性が移る。
従って (2)$\Rightarrow$(3),(4),(5) である。

逆に $P_{64}(c)=0$ なら、
定理\ref{cmo:thm:morse-singular-sixty-four}により
$R_c$ は無限射影次元の三十六次元単純加群を持ち、対角は perfect でない。
同定理は線束生成圏の非 smooth 性、
全 coherent 圏とその残余圏の非 proper 性も与える。
従って (2)--(5) はいずれも成立しない。

最後に (1) の場合、coherent 飽和と
補題\ref{cmo:lem:mixed-line-serre}により全 coherent 圏の Serre 関手は
$(-4)[5]$ である。
coherent 加群間の負次数 Ext の消滅と Serre 双対から五を超える Ext は零となり、
$\operatorname{Ext}^5(\mathcal O,\mathcal O(-4))\simeq\mathbb C$ である。
従って homological dimension は五である。
(1) が成立しない場合の無限性は、同じ定理の無限自己 Ext による。
\end{proof}

\begin{corollary}[零係数と定常 norm の全六十四係数]
\label{cmo:cor:fixed-zero-norm-smooth}
上の指定混合族では、$c=0$ および定常 norm $N(c)$ の全六十四根で
$D^b(\operatorname{qgr}A_c)$ と
$\operatorname{Ku}_{\mathrm{coh}}(A_c)$ は smooth proper である。
norm の各根では二つの楕円曲線に一つの被約定常点が加わり、
$c=0$ では point scheme は二つの楕円曲線だけである。
\end{corollary}
\begin{proof}
定理\ref{cmo:thm:morse-singular-sixty-four}の
$P_{64}(0)\ne0$ と $\gcd(P_{64},N)=1$ を前定理へ適用する。
point scheme の記述は定理\ref{cmo:thm:point-norm}と
例\ref{cmo:ex:sixty-four}による。
\end{proof}

\begin{remark}[smooth locus と可換実現 locus]
\label{cmo:rem:smooth-versus-realization}
この定理は指定二因子について source の smooth locus を完全に決定する。
滑らかな各 $c$ の残余圏が可換 cubic の Ku と同値であることまで
証明したものではない。
系\ref{cmo:cor:cofinite}による余有限可換実現を保持するが、
特に $c=0$ と norm の根での実際の比較核はなお構成されていない。
既知の可換実現 locus が開であることと、
$\mathbb A^1\setminus V(P_{64})$ が連結であることだけから
その locus が全体に等しいとは結論できない。
従って次の核心は、既に滑らかと判定した個別係数での可換実現と、
比較核の境界延長である。
\end{remark}

\subsection{全 smooth locus 上の相対模型と通常の Hochschild 群}

\begin{proposition}[指定混合族の全 smooth locus 上の相対模型と Hochschild homology]
\label{cmo:prop:full-relative-hh}
定理\ref{cmo:thm:exact-sixty-four}の指定係数を固定し、
\[
 R=\C[c],\qquad R^\circ=\C[c,P_{64}(c)^{-1}],
 \qquad U=\operatorname{Spec}R^\circ
\]
と置く。族
\[
 \mathscr B=(\mathcal B_{\alpha^{(1)}}
       \otimes_{\C[u]}\mathcal B_{\alpha^{(2)}})\otimes_\C R,
 \qquad
 \mathscr A=\mathscr B/
       (\Omega_1+100\Omega_2+cu^3)
\]
の有限線束模型と残余模型は、$U$ 上相対 smooth proper である。
残余模型の閉ファイバーは
$\operatorname{Ku}_{\mathrm{coh}}(A_c)$ と同値であり、
全ての $c\in U(\C)$ について
\[
 \dim_\C HH_i(\operatorname{Ku}_{\mathrm{coh}}(A_c))
 =
 \begin{cases}
  2,&i=0,\\
  21,&i=\pm1,\\
  0,&\text{その他}
 \end{cases}
\]
となる。ここで $HH_i$ は Hochschild \emph{homology} である。
この結論は明示可換候補への同値を仮定しない。
\end{proposition}

\begin{proof}
まず、相対 smoothness に必要な通常の有限中心模型を構成する。
$\mathscr B$ は固定した ambient 代数の $R$ への底変更なので、
命題\ref{cmo:prop:proper-contraction}の有限 Koszul 双加群解消、
augmentation 双対および noetherian 性を持つ。
中心元 $w_c=\Omega_1+100\Omega_2+cu^3$ は各閉ファイバーで正則である。
実際、$u=0$ では正則な中心和 $\Omega_0^{(1)}+100\Omega_0^{(2)}$
となるから、$w_cf=0$ なら $f$ は $u$ で割れる。
$u$ は ambient 代数上正則であり、固定内部次数でこの議論を繰り返せば $f=0$ となる。
同じ議論は $R$ 上でも成立する。
従って各内部次数での $w_c$ の乗法は全閉ファイバーで単射となり、
有限自由 $R$-加群間の一定階数判定によって
$\mathscr A_d$ は有限自由である。

命題\ref{cmo:prop:proper-contraction}の証明のうち、
有限解消、正則中心三次、次数別自由性および相対 Gorenstein 型だけを使う部分を
この族に適用する。$N=7,n=5,a=4$ とすると
\[
 E=\bigoplus_{j=0}^{6}\mathcal O(j)
\]
は全 twist を厚生成し、
\[
 \operatorname{Ext}^p(\mathcal O(i),\mathcal O(j))
 =
 \begin{cases}
  \mathscr A_{j-i},&p=0,\\
  (\mathscr A_{i-j-4})^\vee,&p=5,\\
  0,&\text{その他}
 \end{cases}
\]
である。従って $\Sigma_R=R\operatorname{End}(E)$ は $R$-perfect で、
局所化および閉ファイバーへの導来底変更と両立する。
四つの相対例外線束に沿う有限 mutation で得る残余生成対象 $G$
についても $\Gamma_R=R\operatorname{End}(G)$ は $R$-perfect である。
この時点では proper な線束模型だけを用い、相対 smoothness はまだ結論しない。

系\ref{cmo:cor:relative-finite-center}により、$\mathscr A$ は
$u,\Omega_1,\Omega_2$ と各因子の三つの六次中心元の像が生成する
可換有限型中心部分環 $C_R$ 上有限である。
$C_R\subset\mathscr A$ は $R$-torsionfree であり、$R$ は PID なので $R$-flat である。
中心生成元の次数は $1,3,6$ だから $C_R^{(6)}$ は次数一で生成され、
\[
 Z=\operatorname{Proj}C_R^{(6)},\qquad
 \mathcal A=\widetilde{\mathscr A^{(6)}}
\]
は $R$ 上 projective な中央スキームと、その上の coherent な有限代数を与える。
各中央標準開上では対応する次数零局所化が $R$-torsionfree なので、
$Z$ と $\mathcal A$ はいずれも $R$-flat である。
中心部分環のファイバーが各 $A_c$ の全中心と一致することは要求しない。

各六次中央元 $f$ を反転した $L=\mathscr A[f^{-1}]$ は強 $\mathbb Z$-次数付きである。
実際、$\mathscr A$ の次数一生成性から $f$ を長さ六の語の和で書け、
語の最初または最後の一文字を取り出すと
$1\in L_1L_{-1}$ および $1\in L_{-1}L_1$ を得る。
従って全 twist は $\mathcal A$-perfect であり、$E$ は
$D_{\mathrm{qc}}(Z,\mathcal A)$ の compact generator である。
生成性には、$E$ が全 twist を厚生成することと、
中央の ample twist $\mathcal A(m)$ が全準連接導来圏を生成することを用いる。
この記述は各閉ファイバーでも成立する。

$Y=Z\times_R Z$ とし、
\[
 \Lambda=\mathcal A^{\mathrm{op}}\boxtimes_R\mathcal A,
 \qquad M=\Delta_*\mathcal A,\qquad
 F=E^\vee\boxtimes_R E
\]
と置く。$Y$ は $R$ 上 proper である。
$\Lambda,M$ は $Y$ 上 coherent で $R$-flat であり、
全ての $c_0$ について
\[
 \Lambda\otimes_R\C_{c_0}
   =\mathcal A_{c_0}^{\mathrm{op}}\boxtimes_\C\mathcal A_{c_0},
 \qquad
 M\otimes_R\C_{c_0}=\Delta_{c_0*}\mathcal A_{c_0}
\]
である。ここでは $\mathcal A$ の $R$-flatness により、これらの通常の
底変更が導来底変更にも一致する。
$F$ は $D_{\mathrm{qc}}(Y,\Lambda)$ の compact generator となり、
式\eqref{cmo:eq:kernel-diagonal}と同じ局所有限射影評価および
\v Cech--K\"unneth 計算が
\[
 \begin{aligned}
 R\operatorname{End}_\Lambda(F)
 &\simeq\Sigma_R^{\mathrm{op}}\otimes_R^{\mathbf L}\Sigma_R,\\
 R\operatorname{Hom}_\Lambda(F,M)&\simeq\Sigma_R
 \end{aligned}
\]
を左右作用込みで与える。

任意の $c_0\in U(\C)$ を取る。
定理\ref{cmo:thm:exact-sixty-four}によって $\Sigma_{c_0}$ は smooth である。
前の Morita 比較をこのファイバーに適用すると、
$M_{c_0}$ は $\Lambda_{c_0}$ 上 perfect である。
そこで底の変数を $c-c_0$ として、
命題\ref{cmo:prop:proper-perfect-spread}を $\Lambda,M$ に適用する。
必要な有限性・properness と両者の $R$-flatness は上で確認済みなので、
$c_0$ の Zariski 開近傍で $M$ は $\Lambda$-perfect となる。
全ての $c_0\in U(\C)$ についてこの議論を行えば、$U$ はそのような開集合で覆われる。
$U$ は Jacobson かつ準コンパクトなので、有限部分被覆を取ることができ、
$M_U$ は大域的に perfect である。
従って $\Sigma_{R^\circ}$ は相対 smooth proper である。
有限例外列を除く相対 admissible 成分 $\operatorname{Perf}(\Gamma_{R^\circ})$
も相対 smooth proper となる。
閉ファイバーでの coherent 飽和は
定理\ref{cmo:thm:exact-sixty-four}による。

Mathew の相対退化定理 \cite[Theorem~1.3]{CMOMathewDegeneration}により、
各 $HH_i(\operatorname{Perf}(\Gamma_{R^\circ})/R^\circ)$ は有限射影 $R^\circ$-加群である。
Hochschild 複体の導来底変更と、この射影性から通常のファイバーも
$HH_i(\operatorname{Ku}_{\mathrm{coh}}(A_c))$ に一致する。
$U$ は連結なので各階数は一定である。
系\ref{cmo:cor:cofinite}の可換実現 locus は $U$ と交わり、
そこでの標的は滑らかな三次五様体である。
その Hodge 数は $h^{p,p}=1$ ($0\le p\le5$) と
$h^{3,2}=h^{2,3}=21$ で、これら以外は零である。
通常の HKR と四つの例外対象を除く Hochschild homology の加法性から、
その Ku の次数零成分は二次元、次数 $\pm1$ は各二十一次元となる。
階数一定性により同じ公式が $U$ 全体に成立する。
\end{proof}

\begin{remark}[相対化と対合付き比較の区別]
上の相対 smoothness は、fiberwise smooth からの一般的な開性を仮定していない。
通常の coherent な有限中心代数 $\Lambda$ とその対角加群 $M$ を作り、
両者の底上の flatness と中央台の properness を確認した上で、
命題\ref{cmo:prop:proper-perfect-spread}を適用した。
この入力を持たない一般の proper dg 族へ同じ結論を流用しない。

また、$u\mapsto-u$ は $c\ne0$ で $A_c$ と $A_{-c}$ を交換する。
従って上の通常の Hochschild homology の階数一定性は、
$c=0$ における対合の固有空間次元を決めない。
非可換側の対合付き不変量には独立の計算または同変比較が必要である。
\end{remark}

% !TEX root = ../../mixed-sklyanin-cubics.tex
\section{明示可換候補とアファイン MF 比較の限界}
\label{mix:sec:commutative-candidate}

\subsection{同じ判別集合を持つ明示的な可換三次族}

指定族について、可換 Ku 比較の標的候補を具体的に与える。
以下の判別集合と局所特異点の一致は証明済みであるが、
候補への実際の Ku 同値は別の未解決問題である。

\begin{proposition}[Jacobi 環による臨界値の一致]
\label{cmo:prop:commutative-critical-spectrum}
$\alpha^{(1)}=(1,4,9)$、$\alpha^{(2)}=(16,25,49)$ とし、
\[
 f_\alpha(x,y,z)=\frac{x^3+y^3+z^3+xyz}{3}
                 -\alpha_1x-\alpha_2y-\alpha_3z,\qquad
 J_\alpha=\C[x,y,z]/(\partial_xf_\alpha,\partial_yf_\alpha,\partial_zf_\alpha)
\]
と置く。両方の $\alpha$ について $J_\alpha$ は階数八の被約代数であり、
八個の臨界点は全て Morse 型で、その臨界値も相異なる。
$m_{f_\alpha}$ を乗法作用とし、$\lambda^2=28/27$、$\lambda\ne0$ と選ぶと、
命題\ref{cmo:prop:center-critical-polynomials}の規格化の下で
\begin{equation}\label{cmo:eq:commutative-critical-spectrum}
 \det(t-\lambda m_{f_{\alpha^{(i)}}})
    =\frac{Q_i(t^2)}{3^{17}}\qquad(i=1,2).
\end{equation}
\end{proposition}
\begin{proof}
三つの偏微分は
$x^2+yz/3-\alpha_1$ とその巡回置換である。
有理数上の Gröbner 基底から得る標準単項式は八個であり、
この基底での乗法行列を直接計算する。
未規格化の特性多項式を $h_i(t)=\det(t-m_{f_{\alpha^{(i)}}})$ とすると、
\begin{align*}
 h_1(t)={}&t^8-\frac{88490}{63}t^6
 +\frac{2638222849}{3969}t^4\\
 &{}-\frac{1087802003896}{9261}t^2
 +\frac{488560998869056}{83349},\\
 h_2(t)={}&t^8-\frac{5343910}{21}t^6
 +\frac{27957473572387}{1323}t^4\\
 &{}-\frac{145061299763739531640}{250047}t^2
 +\frac{92587433591992310910016}{83349}.
\end{align*}
両多項式は平方因子を持たず、Hessian 行列式の乗法行列も可逆である。
これにより臨界点・値の相異性、被約性、Morse 性を得る。
係数を比較すると
$\lambda^8h_i(t/\lambda)=Q_i(t^2)/3^{17}$ である。
全乗法行列・Hessian・恒等式は
\texttt{computations/verify\_mixed\_commutative\_candidate.py}
が元の三つの偏微分から厳密に再生成する。
標準単項式の範囲は Gröbner 基底の純冪先頭項から定め、
有限次数の探索を商全体の次元判定に代用しない。
\end{proof}

\begin{theorem}[判別集合と局所節点が一致する可換候補]
\label{cmo:thm:commutative-discriminant-candidate}
$G_\alpha(x,y,z,u)$ を $f_\alpha$ の次数三斉次化、すなわち
\[
 G_\alpha=\frac{x^3+y^3+z^3+xyz}{3}
              -u^2(\alpha_1x+\alpha_2y+\alpha_3z)
\]
とし、$\lambda=2\sqrt{21}/9$ を固定する。
\begin{equation}\label{cmo:eq:commutative-candidate}
 X_c=V\!\left(\lambda\bigl(G_{\alpha^{(1)}}(x_1,y_1,z_1,u)
            +100G_{\alpha^{(2)}}(x_2,y_2,z_2,u)\bigr)+cu^3\right)
       \subset\PP^6_\C
\end{equation}
と置くと、次が成立する。
\begin{enumerate}
\item $X_c$ が滑らかであることと $P_{64}(c)\ne0$ は同値である。
      従って $c=0$ と定常ノルムの全六十四根でも $X_c$ は滑らかである。
\item $P_{64}(c)=0$ の各係数では、$X_c$ はただ一つの特異点 $p_c$ を持ち、
      それは五次元通常二重点である。対応する非可換側の中央点 $z_c$ について、
      非標準的な同型
      \[
       \widehat{\mathcal O}_{X_c,p_c}
          \simeq\C[[v_1,\ldots,v_6]]/(v_1^2+\cdots+v_6^2),\qquad
       \widehat{(R_c)}_{z_c}
          \simeq\operatorname{Mat}_{36}(\widehat{\mathcal O}_{X_c,p_c})
      \]
      がある。
\item 両ブロックの Jacobian Poisson tensor の和が定める $\pi_c$ に対して、
      $X_c$ が滑らかなら
      \[
       Z_{X_c}(\pi_c)=E_1\sqcup E_2,\qquad
       E_i=V(x_i^3+y_i^3+z_i^3+x_iy_iz_i)\subset\PP^2
      \]
      がスキームとして成立する。
\end{enumerate}
これらの結論は $\Ku_{\mathrm{coh}}(A_c)\simeq\Ku(X_c)$ を含まない。
\end{theorem}
\begin{proof}
$u=0$ 上で六つの active 偏微分が同時に零になるなら、
滑らかな Hesse 三次の gradient が各ブロックで零になる。
各ブロックベクトルは零となるので、射影点は存在しない。
従って $X_c$ は射影境界で全ての $c$ に対して滑らかである。

$u=1$ 上では、特異点は $f_{\alpha^{(1)}}$ と
$f_{\alpha^{(2)}}$ の臨界点の組で、
$\lambda f_{\alpha^{(1)}}+100\lambda f_{\alpha^{(2)}}+c=0$
を満たすものに等しい。
前命題の臨界値多項式から、その存在条件は正確に
式\eqref{cmo:eq:morse-resultant}の $P_{64}(c)=0$ である。
$P_{64}$ が平方因子を持たず次数六十四であることから、
八掛ける八の臨界値の和は全て相異なり、各特異ファイバーの特異点は一つである。
六変数の Hessian は二つの可逆な三変数 Hessian の直和なので、
形式的 Morse 補題が表示した通常二重点を与える。
非可換側の同型は定理\ref{cmo:thm:morse-singular-sixty-four}の
行列代数模型と合わせて得る。
これは各特異点の完備局所代数の比較であり、族全体の大域同値ではない。

$P_{64}(0)\ne0$ および $\gcd(P_{64},N)=1$ は既存の厳密検算による。
再現コードは可換 Jacobi 環の特性多項式だけから終結式を作り、
その六十五係数の正準 hash が既存の $P_{64}$ と一致することも確認する。
最後の零スキームは命題\ref{cmo:prop:poisson-zero}を適用する。
\end{proof}

\begin{corollary}[有理数値 Euler 形式と零次 Hochschild 対合]
\label{cmo:cor:rational-euler-comparison}
$P_{64}(c)\ne0$ とする。
$\mathcal K_c=\operatorname{Ku}_{\mathrm{coh}}(A_c)$ と
$\mathcal K_c^{\mathrm{cl}}=\operatorname{Ku}(X_c)$ のそれぞれで、
四つの例外線束の右直交への射影を $\operatorname{pr}$ とし、
\[
 e_4=[\operatorname{pr}(\mathcal O(4))],\qquad
 e_5=[\operatorname{pr}(\mathcal O(5))]
\]
と置く。この二類は各圏の $K_{\mathrm{num}}\otimes_\mathbb Z\mathbb Q$
の基底であり、Euler 形式は両側で
\[
 (\chi(e_i,e_j))_{i,j=4,5}
 =
 \begin{pmatrix}
  -27&-189\\
  -162&-1161
 \end{pmatrix},
 \qquad \det=729
\]
となる。従って $e_i\mapsto e_i^{\mathrm{cl}}$ は有理 Euler 空間の同型を与える。
基底 $(e_4,e_5-7e_4)$ での行列は
\[
 27\begin{pmatrix}-1&0\\1&-1\end{pmatrix}
\]
である。ここでは二類が積分数値格子の基底であるとは主張しない。

$c=0$ では $u\mapsto-u$ は両側の
$HH_0$ と $K_{\mathrm{num}}\otimes\mathbb Q$ 上で恒等となり、trace は二である。
非可換側の $HH_{\pm1}$ 上の trace はこの系では決まらない。
\end{corollary}

\begin{proof}
両側の線束の Euler 数は、次数別 Hom と Serre 双対から
\[
 \chi(\mathcal O(i),\mathcal O(j))
 =
 \binom{j-i+6}{6}-\binom{j-i+3}{6}
\]
である。負の上添字の二項係数は通常の多項式として解釈する。
例外列 $\mathcal O,\ldots,\mathcal O(3)$ の Euler 行列に対し
$\chi(\mathcal O(i),e_j)=0$ ($0\le i\le3$) を解くと、
$\mathcal O,\ldots,\mathcal O(5)$ に関する二類の列は
\[
 P=
 \begin{pmatrix}
 28&196\\
 -34&-210\\
 21&113\\
 -7&-28\\
 1&0\\
 0&1
 \end{pmatrix}
\]
となる。六本の線束の Euler 行列を $H$ とすれば、
$P^{\mathsf T}HP$ が主張した二次行列を与える。

dg Hirzebruch--Riemann--Roch
\cite[Theorem~3.5 and \S3.2]{CMOShklyarovHRR}によって、
Euler 形式は $HH_0$ に値を取る Chern character の双線形 pairing を通じて因子化する。
表示行列が非退化なので $\operatorname{ch}(e_4),\operatorname{ch}(e_5)$ は一次独立である。
命題\ref{cmo:prop:full-relative-hh}と可換側の HKR によって、
両側で $\dim HH_0=2$ だから、この二つの character は $HH_0$ の基底となる。
したがって、二次行列として表示された Hochschild pairing も非退化である。

任意の $x\in K_0\otimes\mathbb Q$ について、
表示行列の逆行列を用いて一意な $a,b\in\mathbb Q$ を選べば、
$x-ae_4-be_5$ は $e_4,e_5$ との右側 Euler pairing がともに零となる。
character の基底性と前段の非退化性から
$\operatorname{ch}(x-ae_4-be_5)=0$ であり、
dg HRR によってこの差は全対象と両方向に Euler 直交する。
よって $e_4,e_5$ は有理数値 Grothendieck 群を生成する。
行列の非退化性から数値的に独立でもある。

$c=0$ の対合は各 $\mathcal O(j)$ を保ち、その射影とも可換なので、
$e_4,e_5$ とその character を固定する。
基底性から最後の主張を得る。
\end{proof}

\subsection{固定したアファイン曲率の実際の比較とその限界}

ここで $\operatorname{MF}_{\mathrm{pr}}(T,w)$ は、有限生成射影
$\mathbb Z/2$ 次数付き $T$-加群 $E$ と奇自己準同型 $d$、
$d^2=w\operatorname{id}$ からなる dg 圏を表す。
内部次数は忘れ、上付き $\pi$ は pretriangulated hull の冪等完備化を表す。
以下の比較はこの意味の dg Morita 同値であり、内部 shift、
指定周期 $(3)\simeq[2]$、または係数方向の比較核を含まない。

\begin{lemma}[孤立した Azumaya 曲率特異点への完備化]
\label{cmo:lem:isolated-azumaya-mf-completion}
$T$ は可換有限型 $\mathbb C$-代数 $Z\subset Z(T)$ 上有限な代数であり、
左右の大域次元は有限とする。$w\in Z$ は $T$ の非零因子とし、
$\Lambda=T/(w)$ と置く。
$\mathfrak p\in\operatorname{Spec}Z$ は $w\in\mathfrak p$ を満たす閉点で、
$\Lambda_{\mathfrak q}$ は $\mathfrak q\ne\mathfrak p$ の各閉点で有限大域次元を持つとする。
さらに、$\mathfrak p$ における完備化について
\[
 (\widehat T_{\mathfrak p},w)
 \simeq
 \left(\operatorname{Mat}_r(\mathbb C[[v_1,\ldots,v_n]]),
             \sum_{j=1}^n v_j^2\right)
\]
と仮定する。右辺の曲率はスカラー行列として作用する。
このとき完備化関手は
\[
 \begin{aligned}
 \operatorname{MF}_{\mathrm{pr}}(T,w)^\pi
 &\simeq
 \operatorname{MF}_{\mathrm{pr}}
   (\mathbb C[[v_1,\ldots,v_n]],\sum_jv_j^2)^\pi\\
 &\simeq \operatorname{Perf}_{\mathbb Z/2}(\mathrm{Cl}_n)
 \end{aligned}
\]
という dg Morita 同値を与える。
ここで $\mathrm{Cl}_n$ は、奇生成元 $e_1,\ldots,e_n$ と
$e_ie_j+e_je_i=-2\delta_{ij}$ を持ち、微分が零の Clifford 超代数である。
また、全ての $\Lambda_{\mathfrak q}$ が有限大域次元なら
$\operatorname{MF}_{\mathrm{pr}}(T,w)^\pi=0$ である。
\end{lemma}

\begin{proof}
最初に、中央非零因子に対する cokernel の基本事実を確認する。
行列分解を
\[
 P_1\xrightarrow{\phi}P_0\xrightarrow{\psi}P_1,
 \qquad \phi\psi=w,\quad\psi\phi=w
\]
と書く。$\phi,\psi$ は単射であり、法 $w$ で得られる二周期複体と
その双対は完全である。従って $M=\operatorname{coker}\phi$ は
$\Lambda$ 上の totally acyclic な射影解消を持つ。
二つの行列分解の間の $H^0\operatorname{Hom}$ は、対応する cokernel の
射を射影 $\Lambda$-加群を経由する射で割った stable Hom に等しい。
実際、cokernel の射は $P_0$ の射影性で持ち上がり、$\phi$ の単射性により
$P_1$ の射へ延長する。cokernel 上零である射は $\phi$ を経由し、
奇 homotopy を与える。射影 $\Lambda$-加群を経由する射についても、
法 $w$ の双対二周期複体の完全性を使うと同じ homotopy を構成できる。
この議論は中心の局所化および平坦な完備化と両立する。
特に有限大域次元の局所商上では $M$ は射影となり、その行列分解は零対象である。

二つの有限射影行列分解 $E,F$ に対し、
$\operatorname{Hom}_{T}(E,F)$ の各 parity の成分は有限 $Z$-加群である。
従ってその cohomology も有限 $Z$-加群である。
上の局所観察から、その台は $\{\mathfrak p\}$ に含まれるため有限長である。
$Z\to\widehat{Z_{\mathfrak p}}$ は平坦であり、有限射影加群間の Hom と
平坦底変更は可換する。従って完備化は各 Hom 複体の擬同型
(quasi-isomorphism) を誘導する。
これは完備化関手の quasi-fully-faithfulness を示す。
候補点がない場合は同じ議論で全 Hom cohomology が零となる。

本質的全射性を冪等完備後に示す。
$L$ を $\mathfrak p$ 上の単純 $T$-加群とする。
完備化が行列環なので、この単純加群は一意で、完備化後の Morita 比較で
剰余体 $\mathbb C$ に対応する。$L$ は $w$ で消されるため
$\Lambda$-加群として扱える。
$d=\operatorname{gldim}T$ とし、$L$ の十分高い $\Lambda$-syzygy を $M$ とする。
各射影 $\Lambda$-加群の $T$ 上の射影次元は高々一であるから、
次元移動により、$j\ge2$ に対して
\[
 \operatorname{Ext}_T^j(M,N)
 \simeq \operatorname{Ext}_T^{j+k}(L,N)=0
\]
となるよう $k\ge d-1$ を取れる。従って $\operatorname{pd}_T M\le1$ である。
有限生成射影解消 $0\to P_1\xrightarrow\phi P_0\to M\to0$ を選ぶ。
$wM=0$ から $w:P_0\to P_0$ は $\phi$ を経由し、
$\phi\psi=w$ を満たす $\psi:P_0\to P_1$ が存在する。
$\phi$ の単射性と $w$ の中心性から $\psi\phi=w$ でもある。
これにより $M$ は大域的な有限射影行列分解の cokernel になる。

その完備化は、剰余体の stabilization の shift に対応する。
孤立した可換超曲面では剰余体の stabilization が split-generator である
\cite[Theorems~4.1 and~5.2]{ROGDyckerhoff}。
従って完備化関手の像は右辺を split-generate し、
quasi-fully-faithfulness と合わせて冪等完備後の同値を得る。
最後の Clifford 模型は同論文 Section~5.5 による。
この証明では、可換スキームについての完備化不変性を
非可換 order にそのまま適用していない。
\end{proof}

\begin{remark}[Clifford 模型の具体式]
$Q=\mathbb C[[v_1,\ldots,v_n]]$、
$E=Q\otimes\bigwedge^\bullet\mathbb C^n$ と置き、
外積を $\varepsilon_j$、縮約を $\iota_j$ とすると
\[
 d=\sum_jv_j(\varepsilon_j+\iota_j),\qquad d^2=\sum_jv_j^2.
\]
奇作用素 $\iota_j-\varepsilon_j$ は $\operatorname{End}(E)$ の cycle であり、
その反交換子は $-2\delta_{ij}$ となる。
これらから得る Clifford 代数の包含は擬同型 (quasi-isomorphism) である。
従って局所比較は、形式的 Morse 座標、行列環 Morita、
この有限 rank Koszul 行列分解の三つで具体化できる。
形式的局所行列を大域多項式行列として書いたという主張は含まない。
\end{remark}

\begin{proposition}[一因子と混合族の非次数付き affine MF 比較]
\label{cmo:prop:pointwise-affine-mf-comparison}
指定された二つの $\alpha^{(i)}$ と
$\lambda^2=28/27$ を用い、
\[
 T=D_{\alpha^{(1)}}\otimes D_{\alpha^{(2)}},\quad
 w_c=\Omega_1+100\Omega_2+c,\quad
 F_c=\lambda f_{\alpha^{(1)}}+100\lambda f_{\alpha^{(2)}}+c
\]
と置く。このとき任意の固定した $s,c\in\mathbb C$ について
\begin{align*}
 \operatorname{MF}_{\mathrm{pr}}(D_{\alpha^{(i)}},\Omega_i-s)^\pi
 &\simeq
 \operatorname{MF}_{\mathrm{pr}}(\mathbb C[x,y,z],\lambda f_{\alpha^{(i)}}-s)^\pi,\\
 \operatorname{MF}_{\mathrm{pr}}(T,w_c)^\pi
 &\simeq
 \operatorname{MF}_{\mathrm{pr}}(\mathbb C[x_1,y_1,z_1,x_2,y_2,z_2],F_c)^\pi.
\end{align*}
一因子の双方は $Q_i(s^2)\ne0$ なら零、$Q_i(s^2)=0$ なら
$\operatorname{Perf}_{\mathbb Z/2}(\mathrm{Cl}_3)$ と Morita 同値である。
混合族の双方は $P_{64}(c)\ne0$ なら零、$P_{64}(c)=0$ なら
$\operatorname{Perf}_{\mathbb Z/2}(\mathrm{Cl}_6)$ と Morita 同値である。
\end{proposition}

\begin{proof}
まず完備化補題の ambient に関する仮定を確認する。
定理\ref{cmo:thm:azumaya-regular}の証明により、各 $D_{\alpha^{(i)}}$ は
有限型中心上有限、左右 noetherian、三次元 Calabi--Yau で、
左右の大域次元は三である。二つの長さ三の有限射影双加群解消を
$\mathbb C$ 上で tensor すると、$T$ の長さ六の有限射影双加群解消を得る。
従って $T$ の左右の大域次元も有限であり、$T$ は中心の tensor 積上有限なので
左右 noetherian である。
例\ref{cmo:ex:sixty-four}の PBW 濾過では
$\operatorname{gr}D_{\alpha^{(i)}}=S(1,-1,-1)$ であり、
$\Omega_i-s$ の最高部分はその正則中心元 $\Omega_0$ である。
従って $\Omega_i-s$ は全 $s$ で非零因子である。
また $S$ は $\mathbb C[\Omega_0]$ 上自由なので、
$S\otimes S$ は $\mathbb C[\Omega_0,\Omega'_0]$ 上自由である。
従って $w_c$ の最高部分 $\Omega_0+100\Omega'_0$ も正則であり、
$w_c$ は全 $c$ で非零因子である。
この中心和の正則性の議論は、命題\ref{cmo:prop:rational-singular}の証明と同じである。

一因子では、正則中心の非臨界点での商は正則曲面上の Azumaya 代数である。
十四個の中心節点でも、既存の非等方切断の計算により、完備商は
$\mathbb C[[p,q]]\rtimes\mu_2$ と Morita 同値で、有限大域次元を持つ。
この有限射影次元は忠実平坦完備化から降下する。
従って圏的特異点は八つの正則中心上の Morse 点だけである。
臨界値は相異なるので、指定した $s$ に属する圏的特異点は高々一つであり、
存在するときの ambient 完備化は
$(\operatorname{Mat}_6(\mathbb C[[v_1,v_2,v_3]]),\sum v_j^2)$ である。
前補題が適用できる。
命題\ref{cmo:prop:commutative-critical-spectrum}は、
可換側の八つの相異なる Morse 臨界値が厳密に同じ集合であることを与える。
従って可換側にも同じ前補題を適用できる。

混合族では、定理\ref{cmo:thm:exact-sixty-four}の局所的な証明により、
両因子が同時に上の Morse 点となる点を除き、商は有限大域次元を持つ。
$P_{64}$ の squarefree 性と次数六十四から、各悪い $c$ では該当点が唯一である。
そこでは ambient 完備化が三十六次行列環であり、曲率は
六変数の非退化二次式となる。
可換側についても定理\ref{cmo:thm:commutative-discriminant-candidate}が
同じ唯一の Morse 点を与える。再び前補題を両側へ適用する。
\end{proof}

\begin{remark}[射影残余圏を忘れる具体的な障害]
\label{cmo:rem:affine-mf-forgets-residual}
上の同値は実際の固定係数で成立するが、smooth locus では常に
$0\simeq0$ である。一方、命題\ref{cmo:prop:full-relative-hh}から
全ての $c\in U$ で射影残余圏の $HH_0$ は二次元なので、その圏は非零である。
例えば $c=-44828$ は定常 norm の根であり、
$P_{64}(-44828)\ne0$ なので、その affine MF 圏は零である。
一方、命題\ref{cmo:prop:projected-point-ext}の残余射影 $K$ は
\[
 \operatorname{Ext}^0(K,K)=\mathbb C,
 \qquad \operatorname{Ext}^1(K,K)=\mathbb C^{10}
\]
を持ち、$\operatorname{Ku}_{\mathrm{coh}}(A_{-44828})$ は非零である。
従ってこの affine MF 圏を介する関手では、その射影残余圏との同値は得られない。
graded MF の原点と内部 shift、または射影境界との貼合せを保持する
別の比較が必要である。

各 $s$ の比較を単純に族へ貼り合わせることもできない。
$D_\alpha[s]$ 上の曲率 $\Omega-s$ の商は $D_\alpha$ 自身なので、
全空間の通常の有限射影 MF 圏は零である。
その特異な特殊ファイバーの MF 圏は非零であり、
全空間の通常の MF 圏からの底変更では生成されない。
これは相対特異圏などの別の族概念の否定ではなく、
通常の全空間 MF とファイバーの MF を同一視できないという注意である。
\end{remark}

% !TEX root = ../../mixed-sklyanin-cubics.tex
\section{Eckardt 候補の根スタックと同変残余圏}
\label{mix:sec:eckardt}

\subsection{零係数の Eckardt 点と二重被覆}

$V=\C^6$、$x=(x^{(1)},x^{(2)})$ とし、
\[
 J(x_1,x_2,x_3)=\frac{x_1^3+x_2^3+x_3^3+x_1x_2x_3}{3},\qquad
 G(x)=J(x^{(1)})+100J(x^{(2)}),
\]
\[
 L(x)=x^{(1)}_1+4x^{(1)}_2+9x^{(1)}_3
       +100(16x^{(2)}_1+25x^{(2)}_2+49x^{(2)}_3)
\]
と置く。零係数の明示可換候補は
\[
 X_0=V(G-u^2L)\subset\mathbb P(V\oplus\C)
\]
であり、$\lambda$ の選択に依存せず $\mathbb Q$ 上で定義される。

\begin{proposition}[零係数の具体的射影模型]
\label{cmo:prop:zero-eckardt-resolution}
$X_0$ は滑らかであり、$p=[0:\cdots:0:1]$ は Eckardt 型の点である。
その接超平面は $L=0$、接超平面切断は
滑らかな三次三様体
\[
 Y=V(G,L)\subset\mathbb P(V)=\mathbb P^5
\]
を底とし $p$ を頂点とする cone である。

$\mathbb P^5$ 上の四次二重被覆
\[
 S=\operatorname{Spec}_{\mathbb P^5}
      \bigl(\mathcal O\oplus\mathcal O(-2)\bigr),
 \qquad w^2=LG
\]
を考える。分岐因子は超平面 $H=V(L)$ と滑らかな分離型三次四様体
$Z=V(G)$ の和であり、両者は $Y$ に沿って横断的に交わる。
$p$ からの射影は、滑らかな blow-up $\widetilde X_0=\operatorname{Bl}_pX_0$ 上で
\[
 \widetilde X_0\longrightarrow S\longrightarrow\mathbb P^5
\]
と因子化する。第一の射は $S$ の crepant resolution であり、
$S\setminus Y$ 上では同型、$Y$ 上の幾何ファイバーは $\mathbb P^1$ である。
$S$ の $Y$ に沿う完備局所型は $A_1$ surface singularity と滑らかな三次元因子の積である。
\end{proposition}
\begin{proof}
$X_0$ の滑らかさは定理\ref{cmo:thm:commutative-discriminant-candidate}と $P_{64}(0)\ne0$ による。
$G$ の勾配の共通零点は affine 原点だけであるため、$Z$ は滑らかである。
$p$ では $d(G-u^2L)=-dL\ne0$ であり、接超平面は $L=0$ である。
その切断の方程式 $G=L=0$ は $u$ に依存しないので cone の記述を得る。

もし $Y$ が特異なら、ある $x\ne0$ と $a\in\C$ で
$dG(x)=a\,dL$、$L(x)=0$ となる。
$Z$ の滑らかさから $a\ne0$ であり、$u^2=a$ を選べば
$[x:u]\in X_0$ は特異点となって矛盾する。
従って $Y$ は滑らかで、$H$ と $Z$ は横断的に交わる。
二重被覆 $S$ は $Y$ の近傍で $w^2=ab$ と書けるため、
記した $A_1$ 局所型を持ち、その他では滑らかである。

射影の incidence 模型では、基底 $[x]\in\mathbb P(V)$ 上の直線を
$[tx:u]$ と書く。$X_0$ の strict transform は
\[
 t^2G(x)-u^2L(x)=0
\]
である。$t\ne0$ では $w=L(x)u/t$、$u\ne0$ では $w=G(x)t/u$ と置く。
二式は strict transform の方程式により一致し、
基底座標 $x$ の変更に対して次数二で変換するので、二重被覆への射を定める。
$G=L=0$ 上のファイバーは全 $\mathbb P^1$ となり、それ以外では $S$ への射は同型である。
従ってこれは $\widetilde X_0$ からの解消である。

$\pi:\widetilde X_0\to X_0$ の例外因子を $E$、
$X_0$ の超平面類を $H_X$ とすると
\[
 K_{\widetilde X_0}=-4\pi^*H_X+4E
   =-4q^*\mathcal O_{\mathbb P^5}(1).
\]
一方、四次分岐二重被覆の随伴公式は
$K_S=q_S^*\mathcal O_{\mathbb P^5}(-4)$ を与える。
従って解消は crepant である。
\end{proof}

\subsection{Eckardt 候補の根スタック模型と同変残余圏}
\label{cmo:subsec:eckardt-root-model}

この小節では $k=\mathbb C$ とする。
$G\in H^0(\mathbb P^5,\mathcal O(3))$ と
$L\in H^0(\mathbb P^5,\mathcal O(1))$ に対して
\[
 X=V(G-u^2L)\subset\mathbb P^6,\qquad
 Z=V(G)\subset\mathbb P^5,\qquad H=V(L),\qquad Y=H\cap Z
\]
とおき、$X,Z,Y$ が滑らかであると仮定する。
特に $H,Z$ は $Y$ に沿って横断する。
命題\ref{cmo:prop:zero-eckardt-resolution}の明示 $c=0$ 候補は
これらの仮定を満たす。
$\iota(x,u)=(x,-u)$、$\Gamma=\langle\iota\rangle\cong\mu_2$ とし、
$p=[0:\cdots:0:1]$、$b:\widetilde X=\operatorname{Bl}_pX\to X$ とする。

\begin{proposition}[滑らかな商と二つの分離された分岐成分]
\label{cmo:prop:eckardt-smooth-quotient}
自然な $\Gamma$ 作用に関する商は
\[
 T=\widetilde X/\Gamma\cong\operatorname{Bl}_Y\mathbb P^5
\]
である。$q:T\to\mathbb P^5$、$h=q^*\mathcal O_{\mathbb P^5}(1)$、
$e$ を $q$ の例外因子とする。
$H,Z$ の狭義変換 $H',Z'$ は交わらず、
\[
 H'\cong\mathbb P^4,\qquad Z'\cong Z,\qquad
 [H']=h-e,\qquad[Z']=3h-e.
\]
写像 $f:\widetilde X\to T$ は
$M=\mathcal O_T(2h-e)$ と
$s\in H^0(T,M^{\otimes2})$、$\operatorname{div}(s)=H'+Z'$ による
滑らかな二重被覆
\[
 \widetilde X\cong\operatorname{Spec}_T
 \bigl(\mathcal O_T\oplus M^{-1}\bigr)
\]
である。したがって
\[
 \mathcal R=[\widetilde X/\Gamma]
 \cong\sqrt[2]{(T,H'+Z')}
\]
となる。
\end{proposition}

\begin{proof}
点 $p$ からの射影の解消を
$\mathbb P_{\mathbb P^5}(\mathcal O(-1)\oplus\mathcal O)$ 内の
\[
 t^2G(x)=u^2L(x)
\]
と書く。ここで $[tx:u]$ が $X$ 上の点であり、$\Gamma$ は
$[t:u]\mapsto[t:-u]$ と作用する。
二乗写像 $[t:u]\mapsto[a:b]=[t^2:u^2]$ による商は
$\mathbb P_{\mathbb P^5}(\mathcal O(-2)\oplus\mathcal O)$ 内の
$aG=bL$ である。
$L,G$ は正則列なので、その Rees 代数の唯一の一次関係はこの式であり、
これは $\operatorname{Bl}_{(L,G)}\mathbb P^5$ を与える。
$L=0$ と $G=0$ の狭義変換は、それぞれ $a=0$ と $b=0$ であり、
互いに交わらない。各中心 $Y$ は $H,Z$ 内の Cartier 因子なので
その blow-up は元の因子に同型であり、表示した同型と因子類を得る。

四次二重被覆 $w^2=LG$ を $T$ に引き戻すと、右辺は例外因子を
重複度二で含む。例外因子の局所方程式で $w$ を一回割って正規化すると
$M=q^*\mathcal O(2)\otimes\mathcal O_T(-e)$ に値を持つ平方根となり、
残る分岐因子は $H'+Z'$ である。
たとえば $L=a,G=av$ の chart では $z=w/a$ に対して $z^2=v$ となる。
これが上の二乗商と一致する。
分岐因子は滑らかな互いに素な二成分なので被覆は滑らかであり、
二重被覆の局所商 $[\operatorname{Spec}A[z]/(z^2-s)/\mu_2]$ は
対応する二乗根スタックである。
$H'$ と $Z'$ が交わらないため、この単一の根スタックは
二成分に沿う独立した二乗根を取ったスタックにも自然に同型である。
実際、$T\setminus H'$ と $T\setminus Z'$ は $T$ を覆い、
各開集合では一方の根の section が可逆なので、その根スタックは底自身となる。
共通部分では両 section が可逆であるから、この局所同型は一意に貼り合う。
\end{proof}

\paragraph{二つの明示線束.}
$\pi:\mathcal R\to T$ とし、$\mathcal H,\mathcal Z$ を
$H',Z'$ の被約逆像である根因子とする。
\[
 \begin{gathered}
 2\mathcal H=\pi^*H',\qquad
 2\mathcal Z=\pi^*Z',\\
 \mathcal N=\mathcal O_{\mathcal R}(\mathcal H+\mathcal Z)
                   \otimes\pi^*M^{-1},\\
 \mathcal B=\pi^*\mathcal O_T(h)\otimes
                   \mathcal O_{\mathcal R}(\mathcal H)
 \end{gathered}
\]
とおく。$\mathcal N^{\otimes2}\cong\mathcal O_{\mathcal R}$ は
被覆を定義するデータから自明化され、$\mathcal N$ は
$\Gamma$ の非自明指標 $\chi$ に対応する線束である。
この自明化を $\varphi$ と書けば、作用の coherence は
$(-\otimes\mathcal N)^2\xrightarrow{\mathrm{id}\otimes\varphi}\mathrm{id}$
で指定する。線束について
$\varphi\otimes\mathrm{id}_{\mathcal N}
=\mathrm{id}_{\mathcal N}\otimes\varphi$
が成り立つので群作用の結合条件を満たす。
同じデータの有限エタール $\mu_2$-torsor
$\operatorname{Spec}_{\mathcal R}(\mathcal O_{\mathcal R}\oplus\mathcal N)$
は $\widetilde X\to[\widetilde X/\Gamma]$ そのものである。
従って以下では、二乗が恒等に同型という性質だけでなく、
この指定された coherent 作用を用いる。
忘却後の $\mathcal B$ は $b^*\mathcal O_X(1)$ である。
実際、点からの射影では
$b^*\mathcal O_X(1)=\mathcal O_{\widetilde X}(h+E_p)$ で、
$E_p$ が $\mathcal H$ に対応する。
線形化の二つの選択は $\mathcal N$ による捩りだけ異なるため、
以下の両指標を含む集合と生成圏はその選択に依存しない。
根スタックの標準束公式と blow-up の標準束公式から
\begin{equation}
\label{cmo:eq:root-canonical-sign}
 \omega_{\mathcal R}
 =\pi^*\mathcal O_T(-6h+e)
      \otimes\mathcal O_{\mathcal R}(\mathcal H+\mathcal Z)
 =\pi^*\mathcal O_T(-4h)\otimes\mathcal N.
\end{equation}
最後の指標線束を落とすと同変 Serre 関手を誤る。

\begin{theorem}[十六個の明示例外対象による実残余模型]
\label{cmo:thm:eckardt-root-residual}
$j:\mathcal H\hookrightarrow\mathcal R$ と
$\pi_H:\mathcal H\to H'\cong\mathbb P^4$ を自然な射とする。
次の十六個の対象を、表示したブロック順、各ブロック内では指数の昇順に並べる：
\[
 \begin{gathered}
 \mathcal L_{r,\epsilon}=\mathcal B^{\otimes r}\otimes\mathcal N^\epsilon,
 \\
 \mathcal E_{s,\epsilon}=
 j_*\pi_H^*\mathcal O_{\mathbb P^4}(s)\otimes\mathcal N^\epsilon,
 \\ 0\le r,s\le3,\quad\epsilon=0,1.
 \end{gathered}
\]
各同じ指数における二つの指標の順序は任意である。
これらは例外列であり、
\begin{equation}
\label{cmo:eq:eckardt-sixteen-sod}
 D^b(\mathcal R)=
 \left\langle
 \mathcal K,\,
 (\mathcal L_{r,\epsilon})_{r,\epsilon},\,
 (\mathcal E_{s,\epsilon})_{s,\epsilon}
 \right\rangle,
 \qquad
 \mathcal K\simeq\operatorname{Ku}(X)^\Gamma.
\end{equation}
特に $\mathcal K$ は十六個の表示された対象の右直交であり、
$\mathcal N$ による tensor 作用に保たれる。
指定した $\mathcal N^{\otimes2}\cong\mathcal O$ を coherence とする作用について
\begin{equation}
\label{cmo:eq:eckardt-deequivariantization}
 \operatorname{Ku}(X)\simeq\mathcal K^{\langle-\otimes\mathcal N\rangle}
\end{equation}
が成り立つ。これらは自然な dg enhancement に両立する。
\end{theorem}

\begin{proof}
Orlov の blow-up 公式\cite[Theorem~4.3]{CMOOrlovBlowup}を
$p\subset X$ に適用すると、$\operatorname{codim}_Xp=5$ なので
四つの点の成分を得る。
全ての対応、包含、射影および法束には自然な $\Gamma$ 線形化があり、
通常の fully faithful 関手と射影の構成は同変である。
有限群の不変部分を取ることは $\mathbb C$ 上で完全なので、
その Hom 比較と生成の証明は同変圏にも適用できる。
さらに $p$ での接空間作用は $-\mathrm{id}$ であり、
$\mathbb P(T_pX)=\mathbb P^4$ 上の作用は自明である。
したがって
\[
 D^b(\mathcal R)=\left\langle
 \bigl(j_*\pi_H^*\mathcal O_{\mathbb P^4}(-a)
            \otimes\mathcal N^\epsilon\bigr)_{{a=4,3,2,1}\atop{\epsilon=0,1}},
 b^*D^b(X)^\Gamma
 \right\rangle.
\]
法束の自然な線形化による指標のずれは、各 $a$ で両方の指標を
含むためこの表示に影響しない。

三次五様体の
$D^b(X)=\langle\operatorname{Ku}(X),\mathcal O_X,
\mathcal O_X(1),\mathcal O_X(2),\mathcal O_X(3)\rangle$
は $\Gamma$ 不変であり、同変化すると各線束に二つの指標が現れる。
この分解を代入して、上の先頭八対象を Serre 逆関手で最後に移す。
式\eqref{cmo:eq:root-canonical-sign}により
\[
 S_{\mathcal R}^{-1}(-)=(-)\otimes\pi^*\mathcal O_T(4h)
                         \otimes\mathcal N[-5].
\]
$h|_{H'}=\mathcal O_{\mathbb P^4}(1)$ なので、先頭の
$-4,-3,-2,-1$ は $0,1,2,3$ となる。
共通シフトを生成圏から除き、二つの指標を入れ替えることで
式\eqref{cmo:eq:eckardt-sixteen-sod}を得る。

$-\otimes\mathcal N$ は表示された十六個の対象の集合を保つので
$\mathcal K$ を保つ。
商スタックの同定の下でこの作用は $\Gamma$ の指標作用そのものである。
$\operatorname{Ku}(X)$ は冪等完備な $\mathbb C$ 線型圏であり、
$\Gamma$ は有限可換群なので、Elagin の二重同変化
\cite[Theorem~4.2, Section~8]{MEDElagin}から
式\eqref{cmo:eq:eckardt-deequivariantization}が従う。
幾何学的作用、有限群の平均化、射影および tensor の全てを
自然な enhancement 上で行えるため、ここでは単なる加法圏の同値に留まらない。
\end{proof}

\paragraph{残余への明示射影核.}
十六対象を式\eqref{cmo:eq:eckardt-sixteen-sod}の順に
$F_1,\ldots,F_{16}$ と書く。
各左 mutation は $\mathcal R\times\mathcal R$ 上の perfect 核
\[
 P_i=\operatorname{Cone}
 \left(F_i^\vee\boxtimes F_i\longrightarrow\mathcal O_\Delta\right)
\]
で与えられる。従って、対応する関手の合成
$\mathbb L_{F_1}\circ\cdots\circ\mathbb L_{F_{16}}$
（右端を先に適用する）は、$\mathcal K$ を像とする射影である。
その畳込み核は $\mathcal K$ 上で恒等となる。
これは可換標的側の実際の射影核であり、
非可換側から $\mathcal K$ への比較核をまだ与えない。

\begin{corollary}[同変 Hochschild 群]
\label{cmo:cor:eckardt-equivariant-hh}
次数付き複素ベクトル空間として、非標準的な同型
\[
 HH_*(\operatorname{Ku}(X)^\Gamma)
 \cong HH_*(\operatorname{Ku}(Z))\oplus
         HH_*(\operatorname{Ku}(Y))
\]
が存在する。具体的な次元は
\[
 \dim HH_i(\operatorname{Ku}(X)^\Gamma)=
 \begin{cases}
 24&i=0,\\
 5&i=\pm1,\\
 1&i=\pm2,\\
 0&\text{その他}.
 \end{cases}
\]
\end{corollary}

\begin{proof}
Ishii--Ueda の根スタック分解
\cite[Theorem~1.6]{CMOIshiiUedaRoot}を $H'\sqcup Z'$ に適用し、
さらに $T=\operatorname{Bl}_Y\mathbb P^5$ の余次元二の blow-up
分解を用いる。これは、順序と埋込みを省略した加法的不変量について
\[
 HH_*(\mathcal R)\cong
 HH_*(\mathbb P^4)\oplus HH_*(Z)\oplus
 HH_*(Y)\oplus HH_*(\mathbb P^5)
\]
を与える。
三次四様体 $Z$ と三次三様体 $Y$ の標準残余分解はそれぞれ
三個と二個の例外線束を含むため、右辺の次数零の例外部分は
$5+3+2+6=16$ 次元である。
定理\ref{cmo:thm:eckardt-root-residual}の十六例外対象と比較し、
有限次元の次数付きベクトル空間で同次数の直和因子を相殺すれば最初の式を得る。
HKR と滑らかな三次超曲面の Hodge 数から
\[
 \begin{array}{c|ccc}
 &HH_0&HH_{\pm1}&HH_{\pm2}\\\hline
 \operatorname{Ku}(Z)&22&0&1\\
 \operatorname{Ku}(Y)&2&5&0
 \end{array}
\]
であり、次元表を得る。
\end{proof}

\begin{proposition}[零係数の明示可換候補における対合指標]
\label{cmo:prop:zero-involution-hh}
命題\ref{cmo:prop:zero-eckardt-resolution}の記法を用い、
$X_0=V(G-u^2L)\subset\mathbb P^6$、
$Z=V(G)\subset\mathbb P^5$、$Y=V(G,L)\subset\mathbb P^5$ とする。
$\iota(x,u)=(x,-u)$ と置くと
\[
 \operatorname{Fix}(\iota)=Z\sqcup\{p\}.
\]
$\iota$ は $\operatorname{Ku}(X_0)$ を保ち、
その Hochschild homology 上の固有空間次元は
\[
\begin{array}{c|rr}
 i&\dim HH_i^+&\dim HH_i^-\\ \hline
 -1&5&16\\
 0&2&0\\
 1&5&16
\end{array}
\]
である。これら以外の次数は零である。
\end{proposition}

\begin{proof}
$\mathbb P^6$ の二つの固有部分空間は $u=0$ と一点 $p$ であり、
$X_0$ との交わりは $Z$ と $p$ である。
命題\ref{cmo:prop:zero-eckardt-resolution}により $X_0,Z,Y$ は滑らかである。

$F=G-u^2L$ とし、その七変数 Jacobian 環を $J_F$ とする。
$\operatorname{Sym}^2(\C^6\oplus\C_{\mathrm{sgn}})$ の
$\pm$ 固有次元は $(22,6)$ である。
六つの $x$ 微分は偶、$u$ 微分は奇であり、
滑らかさにより七つの二次偏微分は一次独立である。
従って $(J_F)_2$ の固有次元は $(16,5)$ となる。
Griffiths の residue 同型
\cite[Chapter~1, Theorem~4.21]{CMOHuybrechtsCubics}
は $H^{3,2}(X_0)$ を $(J_F)_2$ から与えるが、
その作用には ambient volume form の determinant character が掛かる。
ここでは $F$ は不変で $\det(\iota)=-1$ なので、
$H^{3,2}$ の固有次元は $(5,16)$ である。
実際、斉次座標 $z_0,\ldots,z_6$ の体積形式
$\sum_j(-1)^jz_j\,dz_0\wedge\cdots\wedge\widehat{dz_j}\wedge\cdots\wedge dz_6$
は線型変換の determinant で変換する。
これが $F$ の冪を分母とする residue 表示の符号を与える。
複素共役により $H^{2,3}$ も同じである。
対角 Hodge 成分は全て超平面類の冪であり、$\iota$ は恒等に作用する。
通常の HKR と $\iota$ に不変な四つの例外線束に対する加法性から、
最初の表を得る。
独立の符号確認として、固定集合の Euler 数は
$\chi(Z)+1=27+1=28$ である。
Lefschetz 数の式
$6-\operatorname{Tr}(\iota\mid H^5)=28$ も
$\operatorname{Tr}(\iota\mid H^{3,2})=-11$ と整合する。

orbifold HKR \cite[Corollary~1.17(2)]{CMOOrbifoldHKR}により
\[
 HH_*([X_0/\mu_2])
 \simeq HH_*(X_0)^{\mu_2}\oplus HH_*(Z)\oplus HH_*(p).
\]
Hochschild homology のこの公式には age shift は入らない。
右辺の次数零次元は $6+25+1=32$、
次数 $\pm1$ は各五、次数 $\pm2$ は各一である。
四つの例外線束は自然に $\mu_2$-linearize できる。
同変化後は各線束から $\operatorname{Rep}(\mu_2)$ が現れ、
二つの指標によって合計八つの例外対象を得る。
同変射影は有限群の平均化と元の射影で構成されるので、
通常の半直交分解を同変化できる。
従って次数零から八を引くと、
系\ref{cmo:cor:eckardt-equivariant-hh}の根スタックによる計算を独立に再現する。

\end{proof}

\begin{remark}[Jacobian 指標の有限検算と比較への制約]
滑らかな $F$ の七つの二次偏微分は正則列なので、
同変 Hilbert 級数は
\[
 \sum_{d\ge0}\operatorname{Tr}(\iota\mid(J_F)_d)t^d
 =
 \frac{(1-t^2)^6(1+t^2)}{(1-t)^6(1+t)}
 =(1+t)^5(1+t^2).
\]
次数 $0,\ldots,7$ の trace は
$(1,5,11,15,15,11,5,1)$ である。
特に $(J_F)_3$ の固有次元は $(25,10)$ となる。
この次数三計算は $H^1(X_0,T_{X_0})$ の対合付き変形接空間も計算する。


\texttt{computations/verify\_mixed\_comparison\_invariants.py} は指定した有理係数 $3F$ に対し、
次数二・三の Jacobian 関係行列を $\mathbb Q$ 上で構成して
固有部分の階数を計算し、上の指標と同変 Hochschild homology の算術を検算する。
通常または同変 Hochschild homology の一致から、
圏の同値や半直交分解の相殺を結論しない。
また主予想は対合を標識しないため、特定の $\iota$ の相違が見つかっても、
それだけで主予想への反例となるわけではない。
\end{remark}

\paragraph{結論の範囲.}
この計算では $\operatorname{Ku}(Y)$ と $\operatorname{Ku}(Z)$ が
$\operatorname{Ku}(X)^\Gamma$ の半直交成分となることを主張しない。
二つの異なる分解に現れる例外対象の個数の一致から、残る圏の同値や
半直交分解は一般には導けない。
また、同変 Hochschild 群の一致は、非可換 $A_0$ との
$\Gamma$ 同変比較核の必要条件に過ぎない。
通常の主予想自体は、$u\mapsto-u$ の作用を保存する同値までは要求しない。
上の根スタック模型、線束 $\mathcal N$ と十六対象は比較核の具体的な標的と
検証条件を与えるが、その非可換比較核の構成は未完了である。

\begin{remark}[比較の残る入力]
同じ判別集合、臨界値、完備局所節点を持つことだけでは、
大域的な残余圏の同値は従わない。
特に $D_\alpha$ 自体を多項式環 $\C[x,y,z]$ と通常の derived Morita 同値で
置き換える方法は使えない。
$HH^0$ である中心は derived Morita 不変
\cite[\S3.2]{CMOKellerHochschild}だが、前者の中心には十四節点があり、
後者の中心は滑らかな三次元アファイン空間だからである。
必要なのは中心三次を保持した curved/MF 比較、または全次数付き残余圏の実際の核である。
形式的な $h$ 進比較を $h=1$ へ無条件に特殊化することもしない。
ノルムの根では、可換候補の自然な零スキームと非可換点スキームは異なり、
系\ref{ppz:cor:reduced-point-poisson-counterexample}が全可換五様体の全 Poisson 構造を排除する。
\end{remark}

% !TEX root = ../../mixed-sklyanin-cubics.tex
\section{孤立 Poisson 零点と被約スキームの不一致}
\label{mix:sec:poisson-limits}

ここでは滑らかな三次超曲面上の\emph{全ての} Poisson 構造を扱う。
入力は \cite[定理\texttt{thm:pg-normal}]{NCCubicA} の全次元正規形であり、
一般三次、一般 Poisson 構造、あるいは各ブロックの楕円性を仮定しない。
以下の結論はこの正規形から直接得られ、比較核の存在や
categorical Torelli を入力としない。

退化軌跡の次元に関する既存結果との違いを確認しておく。
Gualtieri--Pym
\cite[Conjecture~1.1, Theorem~7.13]{GualtieriPymDegeneracy}と、
Druel--Pereira--Pym--Touzet
\cite[Theorem~1.2, Corollary~1.3]{DPPTBondalFive}の結果は、
十分大きな次元の既約成分の存在を保証する。
特に後者は五次元 Fano の Poisson 零集合に正次元成分を与えるが、
孤立成分の併存までは排除しない。
以下の三次超曲面に対する全孤立成分の排除とは、結論と入力が異なる。
この一般次元制約についての新規性の網羅的な確認は未了である。

\begin{theorem}[孤立 Poisson 零点の必要次元]
\label{ppz:thm:isolated-zero-dimension}
$X\subset\mathbb P^{n+1}_{\mathbb C}$ を滑らかな三次超曲面、
$n\ge3$ とし、$\pi\in H^0(X,\bigwedge^2T_X)$ を Poisson 構造とする。
零スキーム $Z_X(\pi)$ の台が孤立点を持つなら、
\[
 n=3k\quad(k\ge1),\qquad \operatorname{rank}_{\mathrm{gen}}\pi=2k.
\]
さらに各孤立点における零スキームは零次元局所完全交叉である。
特に $3\nmid n$ なら、$Z_X(\pi)_{\mathrm{red}}$ は孤立点を持たない。
\end{theorem}

\begin{proof}
$\pi=0$ なら零スキームは $X$ 全体であり、結論の仮定を満たさない。
$\pi\ne0$ とする。引用した正規形により、
$n+2=3k+m$、$k\ge1$、$m\ge0$ となる座標
$x_s=(x_{s1},x_{s2},x_{s3})$、$u=(u_1,\ldots,u_m)$ を取れ、
\[
 F=\sum_{s=1}^k f_s(x_s,u),\qquad
 q=\sum_{s=1}^k q_s,\qquad
 q_s=\iota_{d_{x_s}f_s}
 (\partial_{x_{s1}}\wedge\partial_{x_{s2}}\wedge\partial_{x_{s3}})
\]
となる。$q$ の射影化が $\pi$ で、一般階数は $2k$ である。
$f_s$ は斉次三次式であり、中心変数のみの項はいずれかの $f_s$ に配分する。
$g_s(x_s)=f_s(x_s,0)$ と置く。
この $g_s$ は特異でも零式でもよい。

Euler ベクトル場を
$E=\sum_{s,i}x_{si}\partial_{x_{si}}+\sum_\alpha u_\alpha\partial_{u_\alpha}$
と書く。原点を除く錐では $\mathbb CE$ が接束の直和因子となるので、
零スキームの理想は $E\wedge q$ の係数を $X$ へ制限したものである。
$X$ の滑らかさにより $\bigwedge^2T_X\hookrightarrow
\bigwedge^2T_{\mathbb P^{n+1}}|_X$ も局所的に直和因子であり、
この係数理想は $X$ 上の切断の零理想そのものである。

まず $u_\alpha\ne0$ の chart を考える。
$q$ は中心方向を含まないので、$E\wedge q$ の
$\partial_{u_\alpha}\wedge\partial_{x_{si}}\wedge\partial_{x_{sj}}$
の係数は、符号を除き
$u_\alpha\partial f_s/\partial x_{s\ell}$ である。
従ってこの chart の零理想は、全 $3k$ 個の active 偏微分
\[
 \left(\frac{\partial f_s}{\partial x_{si}}\ \middle|\
       1\le s\le k,\ 1\le i\le3\right)
\]
で生成される。
正則な局所環 $\mathcal O_{X,p}$ は次元 $n$ なので、
この零スキームの各局所成分は
\[
 \dim_p Z_X(\pi)\ge n-3k=m-2
\]
を満たす。孤立点があるなら $m\le2$ である。
一方、この chart があるので $m\ge1$。
$m=1$ の場合、零点では全 active 偏微分が零となり、
$F=0$ と Euler の等式から $u_1F_{u_1}=0$ も得る。
これは $X$ の滑らかさに反する。
従って $u\ne0$ に孤立点があれば $m=2$、すなわち $n=3k$ である。
この場合、その点の零理想は次元 $n$ の正則局所環の
$n=3k$ 個の元で生成される極大イデアル準素な理想である。
従ってこれらの元は正則列となり、零次元局所完全交叉を得る。

残る $u=0$ に孤立点がないことを、$g_s$ の滑らかさを使わずに示す。
点 $p=[x_1:\cdots:x_k:0]$ を取り、$\pi(p)=0$ とする。
もしある $q_s(p)\ne0$ なら、任意の $t\ne s$ について、
$E\wedge q$ のブロック外積次数 $(1\text{ in }t,2\text{ in }s)$
の成分は $x_t\wedge q_s(p)$ である。
その成分に他の項は寄与しないため、$x_t=0$ となる。
従って $p$ は一つの純ブロック平面にあり、
$g_s(x_s)=F(p)=0$ を満たす。
純ブロック平面の閉集合 $V(g_s)\subset\mathbb P^2$
全体が Poisson 零集合に入る。
実際、この平面上では他ブロックの二次偏微分は零であり、
$E\wedge q=3g_s P_s$ となる。
$g_s\ne0$ なら $V(g_s)$ は全ての点で局所次元一以上の平面三次曲線、
$g_s=0$ なら平面全体なので、$p$ は孤立しない。

次に全ての $q_s(p)=0$ の場合を考える。
非零ブロックが一つだけなら前段の純ブロック曲線に属する。
二つ以上あるなら、それぞれの非零 $x_s$ について
$d g_s(x_s)=0$、従って $g_s(x_s)=0$ である。
これらの各ブロックベクトルを独立にスカラー倍して得る射影線型空間は
$X$ と Poisson 零集合に含まれる。
実際、勾配は二次斉次、$g_s$ は三次斉次なので、
その線型空間の全点でも全 active 勾配と $F$ が零となる。
この線型空間は $p$ を通り次元一以上である。
以上により $u=0$ のどの零点も孤立せず、証明が完了する。
\end{proof}

\begin{corollary}[混合五様体による被約対応の反例]
\label{ppz:cor:reduced-point-poisson-counterexample}
例\ref{cmo:ex:sixty-four}の二因子族を取る。
定常ノルム $N(c)$ の六十四個の根のそれぞれにおいて、
$A_c$ は全 coherent 導来圏と全 coherent Kuznetsov 成分が
smooth proper である非可換三次五様体であり、
\[
 \operatorname{Pt}(A_c)=E_1\sqcup E_2\sqcup\{p_c\}
\]
は被約である。
任意の滑らかな可換三次五様体 $Y$ と任意の Poisson 構造 $\pi$ に対して
\[
 \operatorname{Pt}(A_c)_{\mathrm{red}}
   \not\cong Z_Y(\pi)_{\mathrm{red}}.
\]
従って点スキームと Poisson 零スキームの一致という素朴な期待は、
被約スキームの同型へ弱めても一般には成立しない。
\end{corollary}
\begin{proof}
source の smooth proper 性は系\ref{cmo:cor:fixed-zero-norm-smooth}、
点スキームの表示と被約性は例\ref{cmo:ex:sixty-four}による。
左辺には孤立した連結成分 $\{p_c\}$ があるが、
定理\ref{ppz:thm:isolated-zero-dimension}により
$3\nmid5$ の可換三次五様体の Poisson 零集合には孤立点がない。
孤立点の有無は抽象被約スキーム同型で保存されるので結論を得る。
この議論は同じ Ku を持つ標的の存在を仮定せず、
全ての滑らかな可換三次五様体を一括して排除している。
従って主予想の真偽については結論しない。
\end{proof}

\begin{example}[一つの有理係数の反例]
\label{ppz:ex:reduced-counterexample-rational}
$c=-44828$ を取ればよい。
定常 character は $u=1$ において
\[
 (x_1,y_1,z_1)=(-1,-2,-3),\qquad
 (x_2,y_2,z_2)=(-4,-5,-7)
\]
であり、中心値は $28$ と $448$ なので
$28+100\cdot448+c=0$ を満たす。
六十四個のノルム根が相異なることから、この係数で加わる定常点は一つだけである。
$P_{64}$ と $N$ は互いに素なので $P_{64}(-44828)\ne0$ であり、
この例は上の系の smooth proper な source を持つ。
\end{example}

\begin{remark}[既知入力と未解決点]
本節の全 Poisson 構造への適用は、未公表入力\cite{NCCubicA}の
全次元正規形に依存する。正規形の主張だけでなく、
Hessian の基点自由性、\texttt{lem:pg-basepoint-free-lie} の根基排除、
\texttt{lem:pg-quadratic-lie} の三重ベクトル分解を入力箇所として明記する。
その一般的入力と本節の孤立点の排除を区別する。
三様体の非被約局所環を用いる不一致の例は、
可換実現を扱う主稿の系\paperref{main}{ppz:cor:complex-counterexample}で
実際の Ku 同値とともに述べられている。
本節はノルム根での可換 Ku 比較核を構成せず、主予想を否定しない。
この一般次元の制約に関する文献上の新規性は未確認である。
\end{remark}


\section{残る比較問題}
\label{mix:sec:questions}

\paragraph{指定可換候補への次数付き比較核.}
$P_{64}(c)\ne0$ の全係数について、明示候補 $X_c$ への
$\Ku_{\mathrm{coh}}(A_c)\simeq\Ku(X_c)$ を構成できるか。
最初に必要なのは、一般点での実際の比較核である。
定理\ref{cmo:thm:commutative-discriminant-candidate}の判別集合の一致、
有理 Euler 形式、Hochschild 群の一致だけではこれを与えない。
命題\ref{cmo:prop:pointwise-affine-mf-comparison}の非次数付き affine MF 圏は
smooth な係数で零なので、非零の射影 Ku の比較をそこから得ることもできない。
内部次数、錐の原点、三次周期を保つ核を構成する必要がある。
その後に perfect 整核と特殊ファイバーの同値性を分けて調べる。
特殊化の unit/counit cone による判定は、主稿の
命題\paperref{main}{grl:prop:trait-kernel-test}に従う。

\paragraph{零係数の対合付き模型.}
$c=0$ では、非可換側の $u\mapsto-u$ が $HH_{\pm1}$ に及ぼす作用を求め、
可換側の固有次元 $(5,16)$ と比較する問題が残る。
この検査が整合すれば、定理\ref{cmo:thm:eckardt-root-residual}の
十六例外対象の右直交へ実際の核を構成し、指標線束の作用との両立を調べられる。
ただし通常の可換 Ku 実現は対合の保存を要求しないため、
対合付き比較の障害を通常の可換実現への反例としない。

\paragraph{指定係数の外への拡張.}
全 $\alpha$ に対する有限中心性と Azumaya locus は得られたが、
十四節点と非等方な中央切断に基づく smooth locus の完全決定は
指定した二組の係数に対する結果である。
一般の $\alpha$ で中心の特異点と中央切断が退化する locus を求め、
その外で同じ対角の議論が成立するかを調べる必要がある。
位数六の他の Sklyanin パラメータや線形補正を含む変形、
および中心変数が複数ある混合表示への拡張も別の問題として残る。
また有理数値 Euler 形式の一致は積分数値格子の同型を含まない。

\begin{thebibliography}{99}

\bibitem{NCCubicA}
\emph{滑らかな三次超曲面のPoisson構造：正規形・退化軌跡・代数的量子化},
\texttt{nc-cubic} 論文A、未公表研究稿 A-v61、2026年11月14日。
\href{https://github.com/uedakazushi/nc-cubic/blob/eabe601c1341777742964584dde5314c0bcf64a9/papers/paper_a.tex}{固定 commit の主原稿}。
支持節は \texttt{papers/support/papers/sections/poisson\_geometry.tex}
の命題 \texttt{prop:pg-represent}, \texttt{prop:pg-jacobi-extension}、
定理 \texttt{thm:pg-normal}、系 \texttt{cor:pg-generic}、および
\texttt{general\_dimension\_universal\_quantization.tex}
の定理 \texttt{thm:pg-all-dimensions}、命題 \texttt{prop:pg-ad-relative-chart}。

\bibitem{NFNCCubicB}
\emph{中心三次非可換変形の一般点における可換復元},
\texttt{nc-cubic} 論文B、未公表研究稿 B-v59、2026年11月14日。
\href{https://github.com/uedakazushi/nc-cubic/blob/eabe601c1341777742964584dde5314c0bcf64a9/papers/paper_b.tex}{固定 commit の主原稿}。
\texttt{papers/generic\_recovery/} の
\texttt{main\_result.tex}, \texttt{generic\_descent.tex},
\texttt{coherent\_saturation.tex}, \texttt{field\_model.tex} を参照。

\bibitem{Threefold}
\emph{非可換三次三様体のAbel圏モジュライ：二次包絡とKuznetsov残余射},
未公表研究稿 v22、2026年11月20日。
\href{https://github.com/uedakazushi/nc-cubic-threefold-moduli/blob/6c35d2c84671c2538e61836a6bf3a511392eacce/papers/nc-cubic-threefold-moduli.tex}{固定 commit の主原稿}。
点スキームは \texttt{papers/sections/point\_scheme\_poisson.tex}、
全 Hom の比較は \texttt{jacobian\_abel\_comparison.tex} を参照。

\bibitem{NFMEG}
P.~Etingof and V.~Ginzburg,
\emph{Noncommutative del Pezzo surfaces and Calabi--Yau algebras},
J. Eur. Math. Soc. \textbf{12} (2010), 1371--1416.
\url{https://doi.org/10.4171/JEMS/235}.
\url{https://ems.press/content/serial-article-files/31746},
Theorem~3.3.2, Section~9.

\bibitem{MoriSmith}
I.~Mori and S.~P.~Smith,
\emph{The classification of 3-Calabi--Yau algebras with 3 generators and 3 quadratic relations},
Math. Z. \textbf{287} (2017), 215--241.
\url{https://doi.org/10.1007/s00209-016-1824-5}.
\url{https://arxiv.org/abs/1502.07403v2}, Theorem~4.2.

\bibitem{LBSV}
L.~Le Bruyn, S.~P.~Smith and M.~Van den Bergh,
\emph{Central extensions of three dimensional Artin--Schelter regular algebras},
Math. Z. \textbf{222} (1996), 171--212.
\url{https://doi.org/10.1007/PL00004532}.
\url{https://sites.math.washington.edu/~smith/Research/Central-extns.pdf},
Theorem~4.2.2, Proposition~4.2.5, Section~4.3.

\bibitem{Mayanskiy}
E.~Mayanskiy,
\emph{Poisson cohomology of two Fano threefolds},
J. Algebra \textbf{424} (2015), 21--45.
\url{https://doi.org/10.1016/j.jalgebra.2014.08.049}.
\url{https://arxiv.org/abs/1303.5916v1}, Theorem~1, Section~2.2.
Theorem~1 は Loray--Pereira--Touzet の分類に帰属されている。
スキーム構造には Section~2.2 の Schouten 計算を用いる。

\bibitem{BCMStandardMonomial}
N.~C.~Bonala, O.~Clarke and F.~Mohammadi,
\emph{Standard monomial theory and toric degenerations of Richardson varieties in the Grassmannian},
\url{https://arxiv.org/abs/2103.16208v1}, 2021,
Section~2.4, Theorem~2.3.

\bibitem{NFMStacksProjective}
The Stacks Project Authors,
\emph{The Stacks Project}, Section~27.21, Projective bundles, Tag~01OA.
\url{https://stacks.math.columbia.edu/tag/01OA}.

\bibitem{NFMStacksAffineDescent}
The Stacks Project Authors,
\emph{The Stacks Project}, Lemma~35.37.1, Tag~0245.
\url{https://stacks.math.columbia.edu/tag/0245}.

\bibitem{NFLZSerre}
X.~Lin and S.~Zhang,
\emph{Serre algebra, matrix factorization and categorical Torelli theorem for hypersurfaces},
Math. Ann. \textbf{391} (2025), 163--177.
\url{https://doi.org/10.1007/s00208-024-02915-8}.
定理番号は \url{https://arxiv.org/abs/2310.09927v2}
の Theorem~3.5, Proposition~3.7, Theorems~4.1 and~5.1 による。

\bibitem{NFBFKKernels}
M.~Ballard, D.~Favero and L.~Katzarkov,
\emph{A category of kernels for equivariant factorizations and its implications for Hodge theory},
Publ. Math. Inst. Hautes \'Etudes Sci. \textbf{120} (2014), 1--111.
\url{https://doi.org/10.1007/s10240-013-0059-9}.
\url{https://numdam.org/item/10.1007/s10240-013-0059-9.pdf},
Theorem~5.39 の一般式とその証明、pp.~82--84。

\bibitem{NFPirozhkov}
D.~Pirozhkov,
\emph{A categorical Torelli theorem for hypersurfaces},
Bull. London Math. Soc. \textbf{56} (2024), no.~10, 3075--3089.
\url{https://doi.org/10.1112/blms.13117}.
本文の Lemmas~5.1 and~5.2、Theorem~5.3 は
\url{https://arxiv.org/abs/2208.13604v1} による。

\bibitem{NFHuybrechtsRennemo}
D.~Huybrechts and J.~V.~Rennemo,
\emph{Hochschild cohomology versus the Jacobian ring and the Torelli theorem for cubic fourfolds},
Algebraic Geometry \textbf{6} (2019), no.~1, 76--99.
\url{https://doi.org/10.14231/AG-2019-005}.
\url{https://content.algebraicgeometry.nl/2019-1/2019-1-005.pdf},
Proposition~3.5 and Remark~3.8.

\bibitem{NFDAGX}
J.~Lurie,
\emph{Derived Algebraic Geometry X: Formal Moduli Problems}, November 5, 2011.
\url{https://www.math.ias.edu/~lurie/papers/DAG-X.pdf},
Section~2, Theorems~5.1.9 and~5.3.16, Proposition~5.3.21.

\bibitem{NFBKP}
A.~Blanc, L.~Katzarkov and P.~Pandit,
\emph{Generators in formal deformations of categories},
Compos. Math. \textbf{154} (2018), 2055--2089.
\url{https://doi.org/10.1112/S0010437X18007303}.
\url{https://arxiv.org/abs/1705.00655v1}, Sections~3.3 and~4.3.

\bibitem{NFKuznetsovFCY}
A.~Kuznetsov,
\emph{Calabi--Yau and fractional Calabi--Yau categories},
J. reine angew. Math. \textbf{753} (2019), 239--267.
\url{https://doi.org/10.1515/crelle-2017-0004}.
\url{https://arxiv.org/abs/1509.07657v2},
Section~3.3, Corollaries~3.18 and~4.1, Lemma~3.19.

\bibitem{NFBFKOrlov}
M.~Ballard, D.~Favero and L.~Katzarkov,
\emph{Orlov spectra: bounds and gaps},
Invent. Math. \textbf{189} (2012), 359--430.
\url{https://doi.org/10.1007/s00222-011-0367-y}.
定理番号は \url{https://arxiv.org/abs/1012.0864v3} の Theorem~5.15 による。

\bibitem{NFKuznetsovCubicFourfolds}
A.~Kuznetsov,
\emph{Derived categories of cubic fourfolds},
in \emph{Cohomological and Geometric Approaches to Rationality Problems},
Progress in Mathematics \textbf{282}, Birkh\"auser, 2010, 219--243.
\url{https://arxiv.org/abs/0808.3351}, Theorem~3.1.
著者稿 \url{https://cims.nyu.edu/~tschinke/books/rationality/kuznetsov.pdf}.

\bibitem{NFMGabrielQuotient}
P.~Gabriel, \emph{Des cat\'egories ab\'eliennes},
Bull. Soc. Math. France \textbf{90} (1962), 323--448,
Chapitre~III, \S1, p.~365.
\url{https://doi.org/10.24033/bsmf.1583}.
\url{https://www.numdam.org/article/BSMF_1962__90__323_0.pdf}.

\bibitem{Fraenkel}
A.~McMillan Fraenkel,
\emph{Extensions of Poisson structures on singular hypersurfaces},
\url{https://arxiv.org/abs/1310.6083v2}, 2013,
Proposition~4.1, Theorem~4.3.

\bibitem{Carles}
R.~Carles,
\emph{D\'eformations dans les sch\'emas d\'efinis par les identit\'es de Jacobi},
Ann. Math. Blaise Pascal \textbf{3} (1996), no.~2, 33--62.
\url{https://doi.org/10.5802/ambp.67}.
\url{https://www.numdam.org/item/AMBP_1996__3_2_33_0.pdf},
Section~8, Lemma~8.35.

\bibitem{CarlesMarquez}
R.~Carles and M.~C.~M\'arquez,
\emph{Different methods for the study of obstructions in the schemes of Jacobi},
Ann. Inst. Fourier \textbf{61} (2011), no.~2, 453--490.
\url{https://doi.org/10.5802/aif.2620}.
\url{https://www.numdam.org/item/AIF_2011__61_2_453_0.pdf},
Section~3, Theorem~3.2, Propositions~3.9 and~3.10.

\end{thebibliography}

\end{document}
